% Similitudes directes du plan — Mathématiques Tle C — Cours signé OnBuch+
% Programme officiel MINESEC. Compiler avec : tectonic M14-similitudes-directes-du-plan.tex
\newcommand{\DOCMATIERE}{Mathématiques}
\newcommand{\DOCNIVEAU}{Tle C}
\newcommand{\DOCMODULE}{Configurations et transformations élémentaires du plan}
\newcommand{\DOCLECON}{14}
\newcommand{\DOCTITRE}{Similitudes directes du plan}
% OnBuch+ — préambule « cours signés » ENRICHI (compilé avec tectonic / XeLaTeX)
% Système visuel « Pop » d'OnBuch+ : fond crème (jamais blanc), bordures encre
% épaisses, ombres « dures » décalées, titres Archivo Black en capitales.
% Version MATHÉMATIQUES : graphes (pgfplots/tikz), boîtes pédagogiques
% (définition, propriété, méthode, attention, le savais-tu, exercice résolu,
% exercice, TP, bilan), page de garde et signature OnBuch+ ancrée en bas.
%
% Macros attendues dans main.tex AVANT \input{preamble} :
%   \DOCMATIERE  \DOCNIVEAU  \DOCMODULE  \DOCLECON (numéro)  \DOCTITRE
\documentclass[11pt]{article}

\usepackage{fontspec}
\usepackage[french]{babel}
\usepackage[a4paper,margin=2cm,top=3.4cm,bottom=2.8cm,headheight=1.4cm,footskip=1.3cm]{geometry}
\usepackage{amsmath,amssymb,amsfonts}
\usepackage{mathtools} % \xmapsto, \coloneqq... souvent utilisés par les modèles
\usepackage{cancel} % simplifications barrées
% \abs{z} (module, valeur absolue) et \norm{u} (norme), utilisés par les modèles
\providecommand{\abs}{}\let\abs\relax\DeclarePairedDelimiter{\abs}{\lvert}{\rvert}
\providecommand{\norm}{}\let\norm\relax\DeclarePairedDelimiter{\norm}{\lVert}{\rVert}
\usepackage[table]{xcolor} % [table] : fournit aussi \rowcolors (alternance de lignes)
\usepackage{pagecolor}
\usepackage{tikz}
\usetikzlibrary{arrows.meta,positioning,calc,shapes.geometric,shapes.misc,shapes.arrows,decorations.pathmorphing,decorations.markings,patterns,shadows,mindmap,trees,fit,backgrounds,matrix,intersections}
\usepackage{pgfplots}
\usepackage{pgf-pie} % diagrammes circulaires \pie (statistiques)
\pgfplotsset{compat=1.18}
\usepgfplotslibrary{fillbetween,groupplots} % aires entre courbes (calcul intégral)
\usepackage{enumitem}
\usepackage{fancyhdr}
% [most] charge toutes les bibliothèques tcolorbox (listings, minted, raster,
% documentation...) et gonfle inutilement la mémoire TeX sur un document long
% (~300 boîtes) : on ne charge que ce qui est réellement utilisé.
\usepackage[skins,breakable]{tcolorbox}
\usepackage{graphicx}
\usepackage{colortbl}
\usepackage{booktabs}
\usepackage{tabularx}
\usepackage{diagbox} % case d'en-tête diagonale des tableaux à double entrée
% Colonnes extensibles centrée / gauche / droite (C, L, R), souvent utilisées par les modèles
\newcolumntype{C}{>{\centering\arraybackslash}X}
\newcolumntype{L}{>{\raggedright\arraybackslash}X}
\newcolumntype{R}{>{\raggedleft\arraybackslash}X}
\usepackage{array}
\usepackage{multicol}
\usepackage{siunitx}
% parse-numbers=false : accepte aussi \SI{2,5 \times 10^{-3}}{...}
\sisetup{locale=FR,output-decimal-marker={,},group-minimum-digits=4,parse-numbers=false}
\DeclareSIUnit\ohm{\ensuremath{\symup{\Omega}}} % U+2126 absent de la police
\usepackage{qrcode}
% \degree : commande fréquemment utilisée par les modèles pour un angle ou une
% température en dehors de \SI{}{} (non fournie par les paquets chargés).
\providecommand{\degree}{^{\circ}}
% \qed (fin de démonstration), utilisé par les modèles sans amsthm
\providecommand{\qed}{\hfill$\square$}
% \barre : double barre des valeurs interdites dans un tableau de variations (tkz-tab)
\providecommand{\barre}{\ensuremath{\|}}
% environnement proof (amsthm non chargé) utilisé par les modèles
\ifdefined\proof\else\newenvironment{proof}[1][Démonstration]{\par\noindent\textit{#1.}\ }{\qed\par}\fi
\usepackage{lastpage}
\usepackage{hyperref}
\hypersetup{hidelinks}

% ============================================================
% Palette « Pop » OnBuch+ — mêmes hex que la classe P (plus_ui.dart)
% ============================================================
\definecolor{poporange}{HTML}{F59321}
\definecolor{popdark}{HTML}{E07A0C}
\definecolor{popcream}{HTML}{FFF6E8}
\definecolor{popink}{HTML}{1C1714}
\definecolor{poppurple}{HTML}{7A5AE0}
\definecolor{popgreen}{HTML}{2E9E6A}
\definecolor{poppink}{HTML}{E0457B}
\definecolor{popblue}{HTML}{2F7DE1}
\definecolor{popgold}{HTML}{C98A12}
\definecolor{popmuted}{HTML}{8A7F76}
\definecolor{popline}{HTML}{EADFD2}
\definecolor{popgray}{HTML}{8A7F76}
\colorlet{popgrey}{popgray}
% Alias tolérants (le modèle nomme parfois les couleurs différemment)
\colorlet{popwhite}{white}
\colorlet{popblack}{popink}
\colorlet{popred}{poppink}
\colorlet{popyellow}{popgold}
% Teintes claires pour fonds de tableaux / zones
\colorlet{popblueL}{popblue!10!white}
\colorlet{poporangeL}{poporange!14!white}
\colorlet{popgreenL}{popgreen!12!white}
\colorlet{poppurpleL}{poppurple!12!white}
\colorlet{poppinkL}{poppink!12!white}
\colorlet{popgoldL}{popgold!14!white}

\pagecolor{popcream}

% ============================================================
% Typographie
% ============================================================
\defaultfontfeatures{Ligatures=TeX}
\setmainfont{PlusJakartaSans-Medium.ttf}[Path=../../../fonts/,BoldFont=PlusJakartaSans-Bold.ttf]
% Maths en sans-serif (Fira Math) avec chiffres et lettres droites de Plus
% Jakarta, pour que les formules se fondent dans le texte.
\usepackage[math-style=ISO,bold-style=ISO]{unicode-math}
\setmathfont{FiraMath-Regular.otf}[Path=../../../fonts/]
\setmathfont{PlusJakartaSans-Medium.ttf}[Path=../../../fonts/,range={up/num,up/{Latin,latin},\mathup}]
\usepackage{icomma} % virgule décimale sans espace en mode math
% Caractères mathématiques Unicode absents des polices de texte (Plus Jakarta,
% Archivo Black) : rendus par leur équivalent mathématique, en texte comme en
% formule — sinon ils s'affichent en carré vide (ex. « Arithmétique dans ℤ »).
\usepackage{newunicodechar}
\newunicodechar{ℝ}{\ensuremath{\mathbb{R}}}
\newunicodechar{ℤ}{\ensuremath{\mathbb{Z}}}
\newunicodechar{ℂ}{\ensuremath{\mathbb{C}}}
\newunicodechar{ℕ}{\ensuremath{\mathbb{N}}}
\newunicodechar{ℚ}{\ensuremath{\mathbb{Q}}}
\newunicodechar{𝔻}{\ensuremath{\mathbb{D}}}
\newunicodechar{α}{\ensuremath{\alpha}}
\newunicodechar{β}{\ensuremath{\beta}}
\newunicodechar{γ}{\ensuremath{\gamma}}
\newunicodechar{δ}{\ensuremath{\delta}}
\newunicodechar{ε}{\ensuremath{\varepsilon}}
\newunicodechar{θ}{\ensuremath{\theta}}
\newunicodechar{λ}{\ensuremath{\lambda}}
\newunicodechar{μ}{\ensuremath{\mu}}
\newunicodechar{σ}{\ensuremath{\sigma}}
\newunicodechar{τ}{\ensuremath{\tau}}
\newunicodechar{φ}{\ensuremath{\varphi}}
\newunicodechar{ψ}{\ensuremath{\psi}}
\newunicodechar{ω}{\ensuremath{\omega}}
\newunicodechar{Σ}{\ensuremath{\Sigma}}
% majuscules produites par \MakeUppercase dans les titres (α-aminés -> Α)
\newunicodechar{Α}{\ensuremath{\alpha}}
\newunicodechar{Β}{\ensuremath{\beta}}
\newunicodechar{Γ}{\ensuremath{\Gamma}}
\newunicodechar{Δ}{\ensuremath{\Delta}}
\newunicodechar{Ε}{\ensuremath{\varepsilon}}
\newunicodechar{Θ}{\ensuremath{\Theta}}
\newunicodechar{Λ}{\ensuremath{\Lambda}}
\newunicodechar{Μ}{\ensuremath{\mu}}
\newunicodechar{Τ}{\ensuremath{\tau}}
\newunicodechar{Φ}{\ensuremath{\Phi}}
\newunicodechar{Ψ}{\ensuremath{\Psi}}
\newunicodechar{Ω}{\ensuremath{\Omega}}
\newunicodechar{↦}{\ensuremath{\mapsto}}
\newunicodechar{∈}{\ensuremath{\in}}
\newunicodechar{∉}{\ensuremath{\notin}}
\newunicodechar{∧}{\ensuremath{\wedge}}
\newunicodechar{⇔}{\ensuremath{\Leftrightarrow}}
\newunicodechar{⟺}{\ensuremath{\Longleftrightarrow}}
\newunicodechar{⇒}{\ensuremath{\Rightarrow}}
\newunicodechar{⟹}{\ensuremath{\Longrightarrow}}
\newunicodechar{∘}{\ensuremath{\circ}}
\newunicodechar{∪}{\ensuremath{\cup}}
\newunicodechar{∩}{\ensuremath{\cap}}
\newunicodechar{≡}{\ensuremath{\equiv}}
\newunicodechar{⊂}{\ensuremath{\subset}}
\newunicodechar{⊥}{\ensuremath{\perp}}
\newunicodechar{∃}{\ensuremath{\exists}}
\newunicodechar{∀}{\ensuremath{\forall}}
\newunicodechar{∛}{\ensuremath{\sqrt[3]{\,}}}
\newunicodechar{∜}{\ensuremath{\sqrt[4]{\,}}}
\newunicodechar{⃗}{\ensuremath{{}^{\to}}}
\newunicodechar{ⁿ}{\ifmmode^{n}\else\textsuperscript{\ensuremath{n}}\fi}
\newunicodechar{ˣ}{\ifmmode^{x}\else\textsuperscript{\ensuremath{x}}\fi}
\newunicodechar{ᵇ}{\ifmmode^{b}\else\textsuperscript{\ensuremath{b}}\fi}
\newunicodechar{ʳ}{\ifmmode^{r}\else\textsuperscript{\ensuremath{r}}\fi}
\newunicodechar{ᵘ}{\ifmmode^{u}\else\textsuperscript{\ensuremath{u}}\fi}
\newunicodechar{ʸ}{\ifmmode^{y}\else\textsuperscript{\ensuremath{y}}\fi}
\newunicodechar{ᵃ}{\ifmmode^{a}\else\textsuperscript{\ensuremath{a}}\fi}
\newunicodechar{ᵉ}{\ifmmode^{e}\else\textsuperscript{\ensuremath{e}}\fi}
\newunicodechar{ᵏ}{\ifmmode^{k}\else\textsuperscript{\ensuremath{k}}\fi}
\newunicodechar{ᵖ}{\ifmmode^{p}\else\textsuperscript{\ensuremath{p}}\fi}
\newunicodechar{ᴺ}{\ifmmode^{N}\else\textsuperscript{\ensuremath{N}}\fi}
\newunicodechar{ᶜ}{\ifmmode^{c}\else\textsuperscript{\ensuremath{c}}\fi}
\newunicodechar{ʰ}{\ifmmode^{h}\else\textsuperscript{\ensuremath{h}}\fi}
\newunicodechar{ᵠ}{\ifmmode^{q}\else\textsuperscript{\ensuremath{q}}\fi}
\newunicodechar{ₙ}{\ifmmode_{n}\else\textsubscript{\ensuremath{n}}\fi}
\newunicodechar{ₐ}{\ifmmode_{a}\else\textsubscript{\ensuremath{a}}\fi}
\newunicodechar{ᵢ}{\ifmmode_{i}\else\textsubscript{\ensuremath{i}}\fi}
\newunicodechar{ᵣ}{\ifmmode_{r}\else\textsubscript{\ensuremath{r}}\fi}
\newunicodechar{ₖ}{\ifmmode_{k}\else\textsubscript{\ensuremath{k}}\fi}
\newunicodechar{ᵦ}{\ifmmode_{\beta}\else\textsubscript{\ensuremath{\beta}}\fi}
\newunicodechar{ₓ}{\ifmmode_{x}\else\textsubscript{\ensuremath{x}}\fi}
\newfontfamily\popdisplay{ArchivoBlack-Regular.ttf}[Path=../../../fonts/]
\newfontfamily\poplogo{SpaceGrotesk-Bold.ttf}[Path=../../../fonts/]
\newfontfamily\popsemi{PlusJakartaSans-SemiBold.ttf}[Path=../../../fonts/]
\newfontfamily\popxbold{PlusJakartaSans-ExtraBold.ttf}[Path=../../../fonts/]
\newfontfamily\popxboldls{PlusJakartaSans-ExtraBold.ttf}[Path=../../../fonts/,LetterSpace=3.0]

\setlist[enumerate]{leftmargin=1.7em,itemsep=5pt,topsep=4pt}
\setlist[itemize]{leftmargin=1.7em,itemsep=4pt,topsep=4pt,label={\color{poporange}\textbullet}}
\setlength{\parindent}{0pt}
\setlength{\parskip}{5pt}
\renewcommand{\arraystretch}{1.25}

% pgfplots : style Pop par défaut
\pgfplotsset{
  every axis/.append style={axis line style={popink,line width=1pt},tick style={popink},
    grid style={popline,line width=0.5pt},label style={font=\small},tick label style={font=\footnotesize}},
  popcurve/.style={poporange,line width=1.8pt,no markers,smooth},
}
% Même style disponible dans un \draw TikZ simple (hors pgfplots)
\tikzset{popcurve/.style={poporange,line width=1.8pt}}

% ============================================================
% Pill / squiggle
% ============================================================
\newcommand{\poppill}[3]{%
  \tikz[baseline=-0.55ex]\node[draw=popink,line width=1.2pt,rounded corners=2.6pt,%
    fill=#2,inner xsep=6.5pt,inner ysep=3pt]{%
    \popxboldls\fontsize{7}{7}\selectfont\color{#3}\MakeUppercase{#1}};%
}
\newcommand{\popsquiggle}[1]{%
  \tikz[baseline=-0.4ex]{
    \draw[#1,line width=2.6pt,line cap=round]
      plot [smooth] coordinates {(0,0.09) (0.34,0.01) (0.68,0.21) (1.02,0.01) (1.36,0.10)};
  }%
}
% Mot-clé surligné (marqueur orange sous le texte)
\newcommand{\cle}[1]{{\popxbold\color{popdark}#1}}

% ============================================================
% En-tête / pied
% ============================================================
\pagestyle{fancy}
\fancyhf{}
\renewcommand{\headrulewidth}{0pt}
\renewcommand{\footrulewidth}{0pt}
\lhead{\raisebox{-0.15ex}{\poplogo\fontsize{13}{13}\selectfont\color{popink}OnBuch\color{poporange}+}}
\rhead{\poppill{\DOCMATIERE\ \textbullet\ \DOCNIVEAU\ \textbullet\ Leçon \DOCLECON}{white}{popink}}
\lfoot{\popxbold\fontsize{7.6}{7.6}\selectfont\color{popmuted}\MakeUppercase{Cours signé OnBuch+}\ \textbullet\ {\popsemi\DOCTITRE}}
\rfoot{\tikz[baseline=-0.6ex]\node[rounded corners=8.5pt,fill=popink,minimum height=17pt,minimum width=17pt,inner xsep=5pt,inner sep=0]{\popxbold\fontsize{8}{8}\selectfont\color{popcream}\thepage\,/\,\pageref*{LastPage}};}

% ============================================================
% Titres
% ============================================================
\newcommand{\chaptitre}[2]{%
  \begin{tcolorbox}[enhanced,colback=poporange,colframe=popink,boxrule=2.6pt,arc=11pt,
      drop shadow=popink,left=15pt,right=15pt,top=12pt,bottom=11pt,before skip=3pt,after skip=18pt]
    \poppill{#2}{popink}{popcream}\par\vspace{9pt}
    {\raggedright\hyphenpenalty=10000\exhyphenpenalty=10000\popdisplay\fontsize{22}{24}\selectfont\color{popcream}\MakeUppercase{#1}\par}
  \end{tcolorbox}
}
% Section numérotée : pastille orange avec le numéro + titre Archivo
\newcounter{popsec}
\newcommand{\coursec}[1]{%
  \stepcounter{popsec}\setcounter{popsub}{0}%
  \par\vspace{16pt}\noindent
  \begin{minipage}{\linewidth}
  \tikz[baseline=-0.7ex]\node[draw=popink,line width=1.6pt,fill=poporange,rounded corners=4pt,minimum size=22pt,inner sep=0]{\popdisplay\fontsize{12}{12}\selectfont\color{popcream}\thepopsec};%
  \hspace{8pt}{\popdisplay\fontsize{15}{17}\selectfont\color{popink}\MakeUppercase{#1}}\par
  \vspace{4pt}\noindent{\color{poporange}\rule{36pt}{3.6pt}}
  \end{minipage}\par\nopagebreak\vspace{8pt}
}
\newcounter{popsub}
\newcommand{\courssub}[1]{%
  \stepcounter{popsub}%
  \par\vspace{9pt}\noindent{\popxbold\fontsize{11.5}{13}\selectfont\color{popink}{\color{poporange}\thepopsec.\thepopsub}\hspace{6pt}#1}\par\nopagebreak\vspace{4pt}
}

% ============================================================
% Boîtes pédagogiques — même recette : fond blanc, bordure épaisse colorée,
% ombre dure encre, badge plein. Titre optionnel [#1].
% ============================================================
\tcbset{popbase/.style={breakable,enhanced,colback=white,boxrule=2.3pt,arc=9pt,
  shadow={3pt}{-3pt}{0pt}{popink},left=14pt,right=14pt,top=11pt,bottom=11pt,before skip=11pt,after skip=15pt,
  fontupper=\popsemi\fontsize{10.6}{14.5}\selectfont\color{popink},
  fontlower=\popsemi\fontsize{10.6}{14.5}\selectfont\color{popink}}}
% Coche dessinée (le glyphe ✓ n'existe pas dans les polices du document)
\newcommand{\popcheck}[1]{\tikz[baseline=-0.5ex]\draw[#1,line width=1.6pt,line cap=round,line join=round](-0.13,0.0)--(-0.04,-0.09)--(0.13,0.1);}
\newcommand{\popboxhead}[3]{\poppill{#1}{#2}{white}\ifx\relax#3\relax\else\hspace{6pt}{\popxbold\fontsize{10.6}{12}\selectfont\color{popink}#3}\fi\par\vspace{7pt}}

\newtcolorbox{aretenir}[1][]{popbase,colframe=popgreen,before upper={\popboxhead{À retenir}{popgreen}{#1}}}
\newtcolorbox{exemplebox}[1][]{popbase,colframe=poporange,before upper={\popboxhead{Exemple}{poporange}{#1}}}
\newtcolorbox{definition}[1][]{popbase,colframe=popblue,before upper={\popboxhead{Définition}{popblue}{#1}}}
\newtcolorbox{propriete}[1][]{popbase,colframe=poppurple,before upper={\popboxhead{Propriété}{poppurple}{#1}}}
\newtcolorbox{methode}[1][]{popbase,colframe=popgold,before upper={\popboxhead{Méthode}{popgold}{#1}}}
\newtcolorbox{attention}[1][]{popbase,colframe=poppink,before upper={\popboxhead{Attention}{poppink}{#1}}}
\newtcolorbox{experience}[1][]{popbase,colframe=popblue,colback=popblueL,before upper={\popboxhead{Expérience}{popblue}{#1}}}
% « Le savais-tu ? » : carte violette pleine (contexte, culture, Cameroun)
\newtcolorbox{savaistu}[1][]{popbase,colback=poppurple,colframe=popink,
  fontupper=\popsemi\fontsize{10.4}{14.2}\selectfont\color{white},
  before upper={\poppill{Le savais-tu\,?}{white}{poppurple}\ifx\relax#1\relax\else\hspace{6pt}{\popxbold\color{white}#1}\fi\par\vspace{7pt}}}
% Exercice résolu : énoncé en haut, \tcblower puis la solution sur fond crème
\newcounter{popexo}\newcounter{popexr}
\newtcolorbox{exoresolu}[1][]{popbase,colframe=poporange,colbacklower=popcream,
  segmentation style={popink,line width=1.2pt,dashed},
  before upper={\stepcounter{popexr}\popboxhead{Exercice résolu \thepopexr}{poporange}{#1}},
  before lower={\poppill{Solution}{popink}{popcream}\par\vspace{6pt}}}
% Exercice (non résolu en place) — la correction est dans la partie Corrigés
\newtcolorbox{exercice}[1][]{popbase,colframe=popink,
  before upper={\stepcounter{popexo}\popboxhead{Exercice \thepopexo}{popink}{#1}}}
% Corrigé
\newtcolorbox{corrige}[1][]{popbase,colframe=popgreen,colback=popgreenL,
  before upper={\popboxhead{Corrigé}{popgreen}{#1}}}
% Objectifs / prérequis / bilan
\newtcolorbox{objectifs}{popbase,colframe=popink,colback=poporangeL,
  before upper={\popboxhead{Objectifs de la leçon}{popink}{}}}
\newtcolorbox{prerequis}{popbase,colframe=popink,colback=popblueL,
  before upper={\popboxhead{Prérequis}{popblue}{}}}
\newtcolorbox{bilan}{popbase,colframe=popink,colback=white,boxrule=2.8pt,
  before upper={\popboxhead{Fiche bilan}{poporange}{L'essentiel à connaître pour l'examen}}}
% Figure encadrée avec légende
\newtcolorbox{popfigure}[1][]{enhanced,breakable=false,colback=white,colframe=popline,boxrule=1.4pt,arc=8pt,
  left=10pt,right=10pt,top=10pt,bottom=8pt,before skip=10pt,after skip=12pt,halign=center,
  lower separated=false,halign lower=center,
  fontlower=\popsemi\fontsize{9}{11.5}\selectfont\color{popmuted},
  #1}
\newcommand{\legende}[1]{\par\vspace{6pt}{\popsemi\fontsize{9}{11.5}\selectfont\color{popmuted}#1}}

% ============================================================
% Signature OnBuch+
% ============================================================
\newcommand{\poplogoinline}[1]{{\poplogo\fontsize{#1}{#1}\selectfont\color{popink}OnBuch\color{poporange}+}}
\newcommand{\signatureonbuch}{%
  \begin{tcolorbox}[enhanced,colback=popink,colframe=popink,boxrule=2.2pt,arc=10pt,
      drop shadow=poporange,left=14pt,right=14pt,top=10pt,bottom=10pt,before skip=0pt,after skip=0pt]
    \tikz[baseline=-0.7ex]\node[circle,fill=popgreen,draw=popcream,line width=1.3pt,minimum size=22pt,inner sep=0]{\popcheck{white}};%
    \hspace{9pt}\begin{minipage}[c]{0.62\linewidth}
      {\popxbold\fontsize{12}{13}\selectfont\color{popcream}Signé\ \textbullet\ {\poplogo OnBuch\color{poporange}+}}\par\vspace{2pt}
      {\raggedright\popsemi\fontsize{8}{10.5}\selectfont\color{popline}\MakeUppercase{Équipe pédagogique OnBuch+}\\\MakeUppercase{Conforme au programme officiel MINESEC}\par}
    \end{minipage}\hfill
    {\poplogo\fontsize{22}{22}\selectfont\color{popcream}OnBuch\color{poporange}+}
  \end{tcolorbox}%
}

% ============================================================
% Page de garde
% #1 titre · #2 matière · #3 niveau · #4 module · #5 n° leçon · #6 durée ·
% #7 description / objectifs (texte ou itemize)
% ============================================================
\newcommand{\pagegarde}[7]{%
  \thispagestyle{empty}
  \newgeometry{a4paper,margin=1.7cm,top=1.6cm,bottom=1.4cm}%
  \begin{tikzpicture}[remember picture,overlay]
    \fill[poporange!18] ([xshift=-3.2cm,yshift=-3.6cm]current page.north east) circle (4.6cm);
    \fill[poppurple!14] ([xshift=2.4cm,yshift=7.6cm]current page.south west) circle (3.8cm);
  \end{tikzpicture}%
  {\poplogo\fontsize{30}{30}\selectfont\color{popink}OnBuch\color{poporange}+}\hfill
  \raisebox{0.6ex}{\poppill{Cours signé \textbullet\ Premium}{popink}{popcream}}\par
  \vspace{1pt}\popsquiggle{poppurple}\par
  \vspace{8pt}
  \begin{tcolorbox}[enhanced,colback=poporange,colframe=popink,boxrule=2.8pt,arc=13pt,
      drop shadow=popink,left=18pt,right=18pt,top=12pt,bottom=13pt,after skip=10pt]
    \poppill{#2 \textbullet\ #3}{popink}{popcream}\hspace{5pt}\poppill{#4}{white}{popink}\par\vspace{8pt}
    {\popdisplay\fontsize{14}{14}\selectfont\color{popink}LEÇON #5}\par\vspace{4pt}
    {\raggedright\hyphenpenalty=10000\exhyphenpenalty=10000\popdisplay\fontsize{25}{27}\selectfont\color{popcream}\MakeUppercase{#1}\par}
    \vspace{8pt}
    {\popxbold\fontsize{9.5}{9.5}\selectfont\color{popink}\MakeUppercase{Durée conseillée : #6}}
  \end{tcolorbox}
  \begin{tcolorbox}[enhanced,colback=white,colframe=popink,boxrule=2.2pt,arc=10pt,
      drop shadow=popink,left=16pt,right=16pt,top=10pt,bottom=10pt,after skip=8pt]
    \poppill{Dans cette leçon}{poppurple}{white}\par\vspace{5pt}
    \popsemi\fontsize{9.3}{12.6}\selectfont\color{popink}#7
  \end{tcolorbox}
  \noindent\begin{minipage}[t]{0.487\linewidth}
    \begin{tcolorbox}[enhanced,colback=popgreen,colframe=popink,boxrule=2.2pt,arc=10pt,
        drop shadow=popink,left=11pt,right=11pt,top=10pt,bottom=10pt,height=3.1cm,valign=center]
      \tcbox[colback=white,colframe=popink,boxrule=1.3pt,arc=5pt,boxsep=0pt,nobeforeafter,
        left=3pt,right=3pt,top=3pt,bottom=3pt,valign=center]{\qrcode[height=1.9cm]{https://play.google.com/store/apps/details?id=com.luvvix.onbuch&hl=fr}}%
      \hspace{8pt}\begin{minipage}[c]{0.5\linewidth}
        {\popdisplay\fontsize{11}{12.5}\selectfont\color{white}\MakeUppercase{Télécharge OnBuch}}\par\vspace{4pt}
        {\raggedright\popsemi\fontsize{8.4}{11}\selectfont\color{white}Gratuit \textbullet\ Google Play\\Tuteur IA, annales, concours, résultats\par}
      \end{minipage}
    \end{tcolorbox}
  \end{minipage}\hfill
  \begin{minipage}[t]{0.487\linewidth}
    \begin{tcolorbox}[enhanced,colback=poppurple,colframe=popink,boxrule=2.2pt,arc=10pt,
        drop shadow=popink,left=11pt,right=11pt,top=10pt,bottom=10pt,height=3.1cm,valign=center]
      \tikz[baseline=-0.7ex]\node[circle,fill=white,draw=popink,line width=1.3pt,minimum size=1.6cm,inner sep=0]{\popdisplay\fontsize{24}{24}\selectfont\color{poppurple}+};%
      \hspace{8pt}\begin{minipage}[c]{0.6\linewidth}
        {\popdisplay\fontsize{11}{12.5}\selectfont\color{white}\MakeUppercase{Passe à OnBuch+}}\par\vspace{4pt}
        {\raggedright\popsemi\fontsize{8.4}{11}\selectfont\color{white}Tous les cours signés, Tuteur IA illimité, fiches PDF — onglet OnBuch+ de l'app\par}
      \end{minipage}
    \end{tcolorbox}
  \end{minipage}
  \vspace{8pt}
  \popcontenu
  \vfill
  \signatureonbuch
  \restoregeometry
  \newpage
}

% Rangée « Ce cours contient » (page de garde)
\newcommand{\poptile}[3]{% #1 couleur #2 gros texte #3 légende
  \begin{tcolorbox}[enhanced,colback=#1,colframe=popink,boxrule=2pt,arc=9pt,shadow={2.5pt}{-2.5pt}{0pt}{popink},
      left=6pt,right=6pt,top=9pt,bottom=9pt,halign=center,valign=center,height=1.75cm,width=\linewidth,nobeforeafter]
    {\popdisplay\fontsize{12.5}{13}\selectfont\color{white}\MakeUppercase{#2}}\par\vspace{3pt}
    {\centering\popsemi\fontsize{7.8}{9.5}\selectfont\color{white}#3\par}
  \end{tcolorbox}}
\newcommand{\popcontenu}{%
  \par\noindent{\popxboldls\fontsize{7.5}{7.5}\selectfont\color{popmuted}CE COURS SIGNÉ CONTIENT}\par\vspace{7pt}
  \noindent\begin{minipage}[t]{0.235\linewidth}\poptile{popblue}{Cours}{complet, illustré, pas à pas}\end{minipage}\hfill
  \begin{minipage}[t]{0.235\linewidth}\poptile{poporange}{Méthodes}{et exercices résolus}\end{minipage}\hfill
  \begin{minipage}[t]{0.235\linewidth}\poptile{poppink}{Exercices}{type examen + corrigés}\end{minipage}\hfill
  \begin{minipage}[t]{0.235\linewidth}\poptile{popgreen}{Fiche bilan}{l'essentiel à retenir}\end{minipage}\par}

% Sommaire
\newtcolorbox{sommaire}{enhanced,colback=white,colframe=popink,boxrule=2.2pt,arc=10pt,
  drop shadow=popink,left=18pt,right=18pt,top=13pt,bottom=13pt,before skip=6pt,after skip=20pt,
  before upper={\poppill{Sommaire}{poppurple}{white}\par\vspace{9pt}\popsemi\fontsize{10.6}{14.5}\selectfont\color{popink}}}

% Signature de fin de cours — poussée tout en bas de la dernière page
\newcommand{\findecours}{%
  \par\vfill\nopagebreak
  \begin{center}\popsquiggle{poporange}\end{center}\vspace{4pt}
  \signatureonbuch
}
\AtBeginDocument{\def\llbracket{\mathopen{[\mkern-3.5mu[}}\def\rrbracket{\mathclose{]\mkern-3.5mu]}}} % intervalles d'entiers (glyphe ⟦ absent de la police)
% Glyphes absents de Fira Math : redéfinis à partir de glyphes disponibles
\AtBeginDocument{%
  \def\setminus{\mathbin{\backslash}}\let\smallsetminus\setminus
  \def\vdots{\mathord{\vbox{\baselineskip3.2pt\lineskiplimit0pt\kern4pt\hbox{.}\hbox{.}\hbox{.}}}}%
  \def\ddots{\mathinner{\mkern1mu\raise6.5pt\hbox{.}\mkern2mu\raise3.6pt\hbox{.}\mkern2mu\raise0.7pt\hbox{.}\mkern1mu}}%
  \def\triangle{\mathord{\scalebox{0.9}{$\symup{\Delta}$}}}%
  \def\nabla{\mathord{\raisebox{\depth}{\scalebox{1}[-1]{$\symup{\Delta}$}}}}%
  \def\checkmark{\popcheck{popdark}}%
  \def\star{\mathbin{\ast}}\def\bigstar{\mathord{\ast}}%
  \def\diamond{\mathbin{\scalebox{0.6}{$\Diamond$}}}%
}

%% FIGFIX-BEGIN v5 — correctifs globaux des figures (docs/onbuch-generation/tools/figfix)
\usepackage{adjustbox}\usepackage{etoolbox}\tcbuselibrary{raster}
% toute figure TikZ/pgfplots plus large que la ligne est réduite à la largeur disponible
\BeforeBeginEnvironment{tikzpicture}{\begin{adjustbox}{max width=\linewidth}}
\AfterEndEnvironment{tikzpicture}{\end{adjustbox}}
% légendes pgfplots lisibles par-dessus les courbes
\pgfplotsset{every axis/.append style={legend style={fill=white,fill opacity=0.92,text opacity=1,draw=black!15,inner sep=3pt}}}
% étiquettes d'arêtes (syntaxe edge["..."]) avec fond blanc
\tikzset{every edge quotes/.append style={fill=white,inner sep=1.2pt}}
%% FIGFIX-END
\begin{document}

\pagegarde{Similitudes directes du plan}{Mathématiques}{Tle C}{Configurations et transformations élémentaires du plan}{14}{10 h}{Cette leçon étudie les similitudes directes du plan, transformations qui conservent les angles orientés et multiplient les distances par un rapport constant. Grâce à l'écriture complexe z' = az + b, on apprend à déterminer leurs éléments caractéristiques (centre, rapport, angle) et à transformer les configurations usuelles du plan.
\vspace{4pt}
\begin{multicols}{2}\small
\begin{itemize}
\item À la fin de cette leçon, je sais donner l'écriture complexe et l'expression analytique d'une translation, d'une homothétie et d'une rotation
\item À la fin de cette leçon, je sais définir une similitude directe et reconnaître son écriture complexe z' = az + b avec a ≠ 0
\item À la fin de cette leçon, je sais déterminer le centre, le rapport et l'angle d'une similitude directe donnée par son écriture complexe
\item À la fin de cette leçon, je sais déterminer l'écriture complexe d'une similitude directe à partir de ses éléments caractéristiques
\item À la fin de cette leçon, je sais passer de l'écriture analytique à l'écriture complexe d'une similitude et réciproquement
\item À la fin de cette leçon, je sais déterminer l'image d'un point, d'une droite et d'un cercle par une similitude directe
\end{itemize}\end{multicols}}

\begin{sommaire}
\begin{multicols}{2}
\begin{enumerate}
\item Rappels : écritures complexes des transformations usuelles
\item Définition et écriture complexe d'une similitude directe
\item Éléments caractéristiques : centre, rapport, angle
\item Propriétés de conservation et images des configurations
\item Détermination d'une similitude et applications aux configurations
\item Similitudes : lieux géométriques, constructions et démonstrations
\item Activité d'intégration
\item Méthodes et exercices résolus
\item Exercices
\item Corrigés des exercices
\item Fiche bilan
\end{enumerate}
\end{multicols}
\end{sommaire}

\begin{objectifs}
\begin{itemize}
\item À la fin de cette leçon, je sais donner l'écriture complexe et l'expression analytique d'une translation, d'une homothétie et d'une rotation
\item À la fin de cette leçon, je sais définir une similitude directe et reconnaître son écriture complexe z' = az + b avec a ≠ 0
\item À la fin de cette leçon, je sais déterminer le centre, le rapport et l'angle d'une similitude directe donnée par son écriture complexe
\item À la fin de cette leçon, je sais déterminer l'écriture complexe d'une similitude directe à partir de ses éléments caractéristiques
\item À la fin de cette leçon, je sais passer de l'écriture analytique à l'écriture complexe d'une similitude et réciproquement
\item À la fin de cette leçon, je sais déterminer l'image d'un point, d'une droite et d'un cercle par une similitude directe
\item À la fin de cette leçon, je sais déterminer la nature et les éléments caractéristiques d'une transformation donnée par son écriture complexe
\item À la fin de cette leçon, je sais utiliser une similitude directe pour résoudre un problème de configuration du plan
\end{itemize}
\end{objectifs}

\begin{prerequis}
\begin{itemize}
\item Nombres complexes : forme algébrique, forme trigonométrique, module et argument (Tle, début d'année)
\item Affixe d'un point, affixe d'un vecteur, interprétation géométrique de z' - z
\item Transformations du plan vues en Première : translation, homothétie, symétries
\item Configurations usuelles : triangle, parallélogramme, droite, cercle (Seconde/Première)
\item Trigonométrie : angles orientés et mesures principales (Première)
\end{itemize}
\end{prerequis}

\newpage

\coursec{Rappels : écritures complexes des transformations usuelles}

À Foumban, capitale du royaume Bamoun, un artisan souhaite reproduire un motif traditionnel gravé sur un panneau : il veut l'agrandir de moitié et le tourner de $45^\circ$ pour l'adapter à une façade plus grande. À Douala, un cartographe doit transformer une carte ancienne de la ville pour l'adapter à une nouvelle échelle, tout en conservant scrupuleusement les angles des rues — sinon la carte deviendrait inutilisable. Dans les deux cas, on transforme une figure en la déplaçant, en l'agrandissant, en la tournant, \emph{sans jamais déformer sa forme}.

Comment décrire mathématiquement de telles transformations ? La réponse est d'une élégance remarquable : grâce aux nombres complexes, une seule et même formule, $z' = az + b$, suffit à coder toutes ces transformations, que l'on appellera \cle{similitudes directes} du plan. Avant de découvrir cette formule unificatrice dans la section suivante, nous devons rafraîchir nos souvenirs sur les trois transformations de base — translation, homothétie, rotation — et sur la manière dont les nombres complexes les traduisent. C'est l'objet de cette première section, qui constitue le socle indispensable de toute la leçon.

\courssub{Affixes : le dictionnaire entre géométrie et nombres complexes}

Toute cette leçon repose sur une idée simple mais féconde : le plan est muni d'un repère orthonormé direct $(O\,;\vec{u},\vec{v})$, et à chaque point $M(x\,;y)$ on associe le nombre complexe $z = x + \mathrm{i}y$, appelé \cle{affixe} de $M$. On note $M(z)$. Réciproquement, tout nombre complexe désigne un unique point du plan. La géométrie devient alors du calcul : au lieu de raisonner sur des figures, on calcule sur des nombres.

\begin{definition}[Affixe d'un point, affixe d'un vecteur]
Soit le plan muni d'un repère orthonormé direct $(O\,;\vec{u},\vec{v})$.
\begin{itemize}
\item L'\cle{affixe} du point $M(x\,;y)$ est le complexe $z_M = x + \mathrm{i}y$.
\item L'\cle{affixe} du vecteur $\overrightarrow{AB}$ est $z_{\overrightarrow{AB}} = z_B - z_A$.
\end{itemize}
\end{definition}

Cette dernière formule est la clé de voûte de tout ce qui suit. Elle traduit la relation de Chasles : $\overrightarrow{AB} = \overrightarrow{OB} - \overrightarrow{OA}$, et elle permet de transformer les égalités vectorielles en égalités entre nombres complexes. Retenons aussi deux conséquences immédiates :

\begin{itemize}
\item La distance $AB$ est le module : $AB = |z_B - z_A|$.
\item L'angle $(\vec{u}\,;\overrightarrow{AB})$ est un argument : $(\vec{u}\,;\overrightarrow{AB}) \equiv \arg(z_B - z_A) \ [2\pi]$.
\end{itemize}

Ainsi, \textbf{les distances se lisent sur les modules, les angles sur les arguments}. C'est précisément cette double lecture qui rendra les complexes si efficaces pour étudier les similitudes, transformations qui agissent à la fois sur les distances et sur les angles.

\begin{exemplebox}[Lecture d'affixes]
Soient $A(2 + \mathrm{i})$ et $B(5 + 3\mathrm{i})$. L'affixe du vecteur $\overrightarrow{AB}$ est
\[z_B - z_A = (5 + 3\mathrm{i}) - (2 + \mathrm{i}) = 3 + 2\mathrm{i}.\]
On en déduit $AB = |3 + 2\mathrm{i}| = \sqrt{9 + 4} = \sqrt{13}$, et le vecteur $\overrightarrow{AB}$ a pour coordonnées $(3\,;2)$.
\end{exemplebox}

\courssub{Translation : l'addition complexe}

\textbf{Intuition.} Une translation, c'est un « glissement » du plan : tous les points se déplacent du même vecteur $\vec{u}$, comme un motif de pagne que l'on reproduit à l'identique un peu plus loin. En termes vectoriels, $M'$ est l'image de $M$ par la translation de vecteur $\vec{u}$ lorsque $\overrightarrow{MM'} = \vec{u}$.

Traduisons cette égalité avec les affixes. Si $\vec{u}$ a pour affixe $b$, alors $\overrightarrow{MM'} = \vec{u}$ s'écrit $z' - z = b$, c'est-à-dire $z' = z + b$. La translation, c'est donc simplement \emph{l'addition d'une constante} !

\begin{propriete}[Écriture complexe de la translation]
La translation de vecteur $\vec{u}$ d'affixe $b$ transforme le point $M(z)$ en le point $M'(z')$ avec :
\[z' = z + b.\]
Si $\vec{u}(a\,;\beta)$, c'est-à-dire $b = a + \mathrm{i}\beta$, et si $M(x\,;y)$ et $M'(x'\,;y')$, l'\cle{expression analytique} de la translation est :
\[\begin{cases} x' = x + a \\ y' = y + \beta \end{cases}\]
\end{propriete}

La démonstration tient en une ligne : $z' - z$ est l'affixe de $\overrightarrow{MM'}$, et dire que $\overrightarrow{MM'} = \vec{u}$ revient exactement à dire que $z' - z = b$. En séparant parties réelle et imaginaire de $z' = z + b$, on obtient immédiatement $x' = x + a$ et $y' = y + \beta$.

\begin{exemplebox}[Image par une translation]
Soit la translation de vecteur $\vec{u}(3\,;-1)$, d'affixe $b = 3 - \mathrm{i}$. L'image de $A(-2 + 4\mathrm{i})$ est :
\[z_{A'} = z_A + b = (-2 + 4\mathrm{i}) + (3 - \mathrm{i}) = 1 + 3\mathrm{i},\]
donc $A'(1\,;3)$. Vérification analytique : $x' = -2 + 3 = 1$ et $y' = 4 + (-1) = 3$. Les deux approches coïncident, heureusement !
\end{exemplebox}

\courssub{Homothétie : la multiplication par un réel}

\textbf{Intuition.} Une homothétie de centre $\Omega$ et de rapport $k$ « dilate » ou « contracte » le plan depuis le point $\Omega$, comme le zoom d'un appareil photo : le point $M$ est envoyé sur la demi-droite $[\Omega M)$ (si $k > 0$) ou sur la demi-droite opposée (si $k < 0$), à une distance multipliée par $|k|$. Le centre $\Omega$, lui, ne bouge pas. Vectoriellement : $\overrightarrow{\Omega M'} = k\,\overrightarrow{\Omega M}$.

Traduisons avec les affixes, en notant $\omega$ l'affixe de $\Omega$ : l'égalité $\overrightarrow{\Omega M'} = k\,\overrightarrow{\Omega M}$ devient $z' - \omega = k(z - \omega)$.

\begin{propriete}[Écriture complexe de l'homothétie]
L'homothétie de centre $\Omega(\omega)$ et de rapport $k \in \mathbb{R}^*$ transforme $M(z)$ en $M'(z')$ avec :
\[z' - \omega = k(z - \omega), \qquad \text{soit} \qquad z' = kz + (1 - k)\omega.\]
\end{propriete}

La forme $z' = kz + (1-k)\omega$ est remarquable : elle a bien l'allure $z' = az + b$, avec $a = k$ \textbf{réel} et $b = (1-k)\omega$. Observons aussi que si $z = \omega$, alors $z' = \omega$ : le centre est bien un point fixe.

\begin{exemplebox}[Image par une homothétie]
Soit $h$ l'homothétie de centre $\Omega(1 + 2\mathrm{i})$ et de rapport $k = -2$. Son écriture complexe est :
\[z' = -2z + (1-(-2))(1 + 2\mathrm{i}) = -2z + 3 + 6\mathrm{i}.\]
L'image de $A(3 - \mathrm{i})$ est $z_{A'} = -2(3 - \mathrm{i}) + 3 + 6\mathrm{i} = -6 + 2\mathrm{i} + 3 + 6\mathrm{i} = -3 + 8\mathrm{i}$, donc $A'(-3\,;8)$. Vérifions la géométrie : $\overrightarrow{\Omega A}$ a pour affixe $(3-\mathrm{i})-(1+2\mathrm{i}) = 2 - 3\mathrm{i}$, et $\overrightarrow{\Omega A'}$ a pour affixe $(-3+8\mathrm{i})-(1+2\mathrm{i}) = -4 + 6\mathrm{i} = -2(2 - 3\mathrm{i})$. On a bien $\overrightarrow{\Omega A'} = -2\,\overrightarrow{\Omega A}$.
\end{exemplebox}

\courssub{Rotation : la multiplication par un complexe de module 1}

\textbf{Intuition.} Une rotation de centre $\Omega$ et d'angle $\theta$ fait « pivoter » le plan autour de $\Omega$, comme l'aiguille d'une montre (en sens inverse, car le sens direct est le sens trigonométrique). Deux choses sont conservées : la distance au centre ($\Omega M' = \Omega M$) et l'angle de pivotement $\left(\overrightarrow{\Omega M}\,;\overrightarrow{\Omega M'}\right) = \theta$.

Or, multiplier un complexe par $\mathrm{e}^{\mathrm{i}\theta}$, c'est exactement conserver son module (car $|\mathrm{e}^{\mathrm{i}\theta}| = 1$) et ajouter $\theta$ à son argument (car $\arg(\mathrm{e}^{\mathrm{i}\theta}) = \theta$). La rotation se traduit donc par une multiplication par $\mathrm{e}^{\mathrm{i}\theta}$ — à condition de mesurer les affixes \emph{depuis le centre} $\Omega$.

\begin{propriete}[Écriture complexe de la rotation]
La rotation de centre $\Omega(\omega)$ et d'angle $\theta$ transforme $M(z)$ en $M'(z')$ avec :
\[z' - \omega = \mathrm{e}^{\mathrm{i}\theta}(z - \omega), \qquad \text{soit} \qquad z' = \mathrm{e}^{\mathrm{i}\theta}\,z + b \quad \text{avec} \quad b = \omega\left(1 - \mathrm{e}^{\mathrm{i}\theta}\right).\]
Son expression analytique, avec $\Omega(x_\Omega\,;y_\Omega)$, est :
\[\begin{cases} x' - x_\Omega = \cos\theta\,(x - x_\Omega) - \sin\theta\,(y - y_\Omega) \\ y' - y_\Omega = \sin\theta\,(x - x_\Omega) + \cos\theta\,(y - y_\Omega) \end{cases}\]
\end{propriete}

\begin{experience}[Démonstration : pourquoi $\mathrm{e}^{\mathrm{i}\theta}$ ?]
Démontrons l'écriture complexe de la rotation — un raisonnement formatrice qui utilise toute la puissance de la forme exponentielle.
\begin{enumerate}
\item Par définition de la rotation, $\Omega M' = \Omega M$, donc $|z' - \omega| = |z - \omega|$. Posons $Z = z - \omega$ et $Z' = z' - \omega$ : on a $|Z'| = |Z|$.
\item Toujours par définition, $\left(\overrightarrow{\Omega M}\,;\overrightarrow{\Omega M'}\right) \equiv \theta \ [2\pi]$, c'est-à-dire $\arg(Z') - \arg(Z) \equiv \theta \ [2\pi]$, autrement dit $\arg\!\left(\dfrac{Z'}{Z}\right) \equiv \theta \ [2\pi]$ (pour $M \neq \Omega$).
\item Le quotient $\dfrac{Z'}{Z}$ a donc module $1$ et argument $\theta$ : c'est exactement $\mathrm{e}^{\mathrm{i}\theta}$. D'où $Z' = \mathrm{e}^{\mathrm{i}\theta} Z$, soit $z' - \omega = \mathrm{e}^{\mathrm{i}\theta}(z - \omega)$.
\item En développant : $z' = \mathrm{e}^{\mathrm{i}\theta} z + \omega - \omega\,\mathrm{e}^{\mathrm{i}\theta} = \mathrm{e}^{\mathrm{i}\theta} z + \omega(1 - \mathrm{e}^{\mathrm{i}\theta})$.
\item Pour l'expression analytique, on écrit $\mathrm{e}^{\mathrm{i}\theta} = \cos\theta + \mathrm{i}\sin\theta$ et on développe :
\[(x' - x_\Omega) + \mathrm{i}(y' - y_\Omega) = (\cos\theta + \mathrm{i}\sin\theta)\big[(x - x_\Omega) + \mathrm{i}(y - y_\Omega)\big].\]
En identifiant parties réelle et imaginaire, on obtient les deux formules annoncées.
\end{enumerate}
\end{experience}

\begin{popfigure}
\begin{tikzpicture}[scale=1.6,>=Stealth]
\coordinate (O) at (0,0);
\coordinate (M) at (2.4,0.6);
\coordinate (Mp) at ({2.4*cos(55)-0.6*sin(55)},{2.4*sin(55)+0.6*cos(55)});
\draw[popline] (-0.6,0) -- (3.2,0);
\draw[popline] (0,-0.6) -- (0,2.8);
\fill[popink] (O) circle (1.5pt) node[below left] {$\Omega(\omega)$};
\fill[popblue] (M) circle (1.5pt) node[right] {$M(z)$};
\fill[poporange] (Mp) circle (1.5pt) node[above right] {$M'(z')$};
\draw[popblue,thick] (O) -- (M) node[midway,below,sloped] {\footnotesize $\Omega M$};
\draw[poporange,thick] (O) -- (Mp) node[midway,above,sloped] {\footnotesize $\Omega M' = \Omega M$};
\draw[popdark,->] (0.9,0.225) arc (14:55:0.925);
\node[popdark] at (1.25,0.62) {$\theta$};
\draw[popgreen,dashed,->] (M) to[bend left=25] (Mp);
\end{tikzpicture}
\legende{Rotation de centre $\Omega$ et d'angle $\theta$ : la distance au centre est conservée et le point pivote de l'angle $\theta$. Complexement : $z' - \omega = \mathrm{e}^{\mathrm{i}\theta}(z - \omega)$.}
\end{popfigure}

\begin{methode}[Calculer l'image d'un point par une transformation]
Pour déterminer l'image de $M(z)$ par une transformation donnée par son écriture complexe :
\begin{enumerate}
\item Identifier l'écriture complexe de la transformation (translation $z' = z + b$, homothétie $z' = kz + (1-k)\omega$, rotation $z' = \mathrm{e}^{\mathrm{i}\theta}z + b$).
\item Remplacer $z$ par l'affixe de $M$ et calculer $z'$ en appliquant les règles de calcul sur les complexes.
\item Identifier les coordonnées de $M'$ : si $z' = x' + \mathrm{i}y'$, alors $M'(x'\,;y')$.
\item \emph{Vérifier} si possible sur un dessin ou par une propriété de conservation (distance au centre, angle).
\end{enumerate}
\end{methode}

\begin{exoresolu}[Rotation de centre $\Omega(1 + \mathrm{i})$ et d'angle $\dfrac{\pi}{3}$]
Soit $r$ la rotation de centre $\Omega(1 + \mathrm{i})$ et d'angle $\theta = \dfrac{\pi}{3}$. Déterminer l'écriture complexe de $r$, puis l'image du point $A(2 + \mathrm{i})$.
\tcblower
\textbf{Étape 1 : écriture complexe.} On a $\mathrm{e}^{\mathrm{i}\pi/3} = \cos\dfrac{\pi}{3} + \mathrm{i}\sin\dfrac{\pi}{3} = \dfrac{1}{2} + \mathrm{i}\dfrac{\sqrt{3}}{2}$. L'écriture est $z' = \mathrm{e}^{\mathrm{i}\pi/3} z + b$ avec :
\begin{align*}
b &= \omega\left(1 - \mathrm{e}^{\mathrm{i}\pi/3}\right) = (1 + \mathrm{i})\left(1 - \frac{1}{2} - \mathrm{i}\frac{\sqrt{3}}{2}\right) = (1 + \mathrm{i})\left(\frac{1}{2} - \mathrm{i}\frac{\sqrt{3}}{2}\right)\\
&= \frac{1}{2} - \mathrm{i}\frac{\sqrt{3}}{2} + \mathrm{i}\frac{1}{2} + \frac{\sqrt{3}}{2} = \frac{1 + \sqrt{3}}{2} + \mathrm{i}\,\frac{1 - \sqrt{3}}{2}.
\end{align*}
\textbf{Étape 2 : image de $A$.} On applique la forme factorisée, plus rapide : $z_{A'} - \omega = \mathrm{e}^{\mathrm{i}\pi/3}(z_A - \omega)$. Or $z_A - \omega = (2 + \mathrm{i}) - (1 + \mathrm{i}) = 1$. Donc :
\[z_{A'} = (1 + \mathrm{i}) + \mathrm{e}^{\mathrm{i}\pi/3} \times 1 = 1 + \mathrm{i} + \frac{1}{2} + \mathrm{i}\frac{\sqrt{3}}{2} = \frac{3}{2} + \mathrm{i}\left(1 + \frac{\sqrt{3}}{2}\right).\]
Ainsi $A'\left(\dfrac{3}{2}\,;\,1 + \dfrac{\sqrt{3}}{2}\right)$, soit numériquement $A'(1{,}5\,;\,1{,}87)$ environ.

\textbf{Vérification :} $\Omega A = |z_A - \omega| = |1| = 1$ et $\Omega A' = |z_{A'} - \omega| = |\mathrm{e}^{\mathrm{i}\pi/3}| = 1$. La distance au centre est bien conservée.
\end{exoresolu}

\begin{attention}[Rapport d'homothétie et facteur de rotation : ne pas confondre !]
Une erreur classique consiste à mélanger les deux multiplications :
\begin{itemize}
\item Dans une \textbf{homothétie}, on multiplie par un \textbf{nombre réel} $k$ : cela change les distances (facteur $|k|$) mais ne fait \emph{pas} tourner (l'argument est conservé si $k > 0$, augmenté de $\pi$ si $k < 0$).
\item Dans une \textbf{rotation}, on multiplie par un \textbf{complexe de module 1}, $\mathrm{e}^{\mathrm{i}\theta}$ : cela fait tourner sans changer aucune distance.
\end{itemize}
Autre piège : dans $z' = \mathrm{e}^{\mathrm{i}\theta}z + b$, le nombre $b$ n'est \textbf{pas} l'affixe du centre ! Le centre $\omega$ se retrouve par la relation $b = \omega(1 - \mathrm{e}^{\mathrm{i}\theta})$, c'est-à-dire $\omega = \dfrac{b}{1 - \mathrm{e}^{\mathrm{i}\theta}}$. Nous y reviendrons en détail dans la section 3.
\end{attention}

\courssub{Points fixes d'une transformation}

\begin{definition}[Point fixe]
Un point $M$ est un \cle{point fixe} (ou point invariant) d'une transformation $f$ lorsque $f(M) = M$, c'est-à-dire lorsque son affixe vérifie l'équation $z' = z$.
\end{definition}

La recherche des points fixes est un réflexe fondamental : elle se ramène toujours à résoudre une \textbf{équation d'inconnue complexe} obtenue en remplaçant $z'$ par $z$ dans l'écriture de la transformation.

\begin{exemplebox}[Points fixes des transformations usuelles]
\begin{itemize}
\item \textbf{Translation} de vecteur non nul : $z = z + b$ donne $b = 0$, impossible. Une translation non triviale n'a \emph{aucun} point fixe (ce qui est intuitif : tout le monde glisse !).
\item \textbf{Homothétie} de rapport $k \neq 1$ : $z = kz + (1-k)\omega$ donne $(1-k)z = (1-k)\omega$, d'où $z = \omega$. Un \emph{unique} point fixe : le centre.
\item \textbf{Rotation} d'angle $\theta \not\equiv 0 \ [2\pi]$ : $z = \mathrm{e}^{\mathrm{i}\theta}z + b$ donne $z(1 - \mathrm{e}^{\mathrm{i}\theta}) = b$, d'où $z = \dfrac{b}{1 - \mathrm{e}^{\mathrm{i}\theta}} = \omega$. Un \emph{unique} point fixe : le centre.
\end{itemize}
\end{exemplebox}

\courssub{Synthèse visuelle : un triangle transformé de trois façons}

La figure suivante résume cette section : un même triangle $ABC$ subit une translation, une homothétie de rapport $2$ et une rotation de $60^\circ$. Observe ce qui change et ce qui est conservé dans chaque cas.

\begin{popfigure}
\begin{tikzpicture}[scale=0.85,>=Stealth]
% Panneau 1 : translation
\begin{scope}
\coordinate (A) at (0,0); \coordinate (B) at (1.6,0); \coordinate (C) at (0.5,1.2);
\coordinate (Ap) at (2.2,0.7); \coordinate (Bp) at (3.8,0.7); \coordinate (Cp) at (2.7,1.9);
\draw[fill=popblueL,draw=popblue,thick] (A)--(B)--(C)--cycle;
\draw[fill=poporangeL,draw=poporange,thick] (Ap)--(Bp)--(Cp)--cycle;
\draw[->,popdark,thick] (0.8,0.5) -- (3.0,1.2) node[midway,above,sloped] {\footnotesize $\vec{u}$};
\node[popblue] at (0,-0.35) {\footnotesize $A$}; \node[popblue] at (1.6,-0.35) {\footnotesize $B$}; \node[popblue] at (0.5,1.45) {\footnotesize $C$};
\node[poporange] at (2.2,0.35) {\footnotesize $A'$}; \node[poporange] at (3.8,0.35) {\footnotesize $B'$}; \node[poporange] at (2.7,2.15) {\footnotesize $C'$};
\node at (2,-1) {\footnotesize \textbf{Translation} : $z' = z + b$};
\end{scope}
% Panneau 2 : homothétie
\begin{scope}[xshift=6.2cm]
\coordinate (Om) at (0,0);
\coordinate (A) at (0.8,0.3); \coordinate (B) at (1.5,0.2); \coordinate (C) at (1.0,0.9);
\coordinate (Ap) at (1.6,0.6); \coordinate (Bp) at (3.0,0.4); \coordinate (Cp) at (2.0,1.8);
\fill[popink] (Om) circle (1.5pt) node[below left] {\footnotesize $\Omega$};
\draw[fill=popblueL,draw=popblue,thick] (A)--(B)--(C)--cycle;
\draw[fill=poporangeL,draw=poporange,thick] (Ap)--(Bp)--(Cp)--cycle;
\draw[popline,dashed] (Om)--(Bp); \draw[popline,dashed] (Om)--(Cp);
\node at (1.8,-0.6) {\footnotesize \textbf{Homothétie} ($k=2$) : $z' = 2z - \omega$};
\end{scope}
% Panneau 3 : rotation
\begin{scope}[xshift=12.4cm]
\coordinate (Om) at (0.9,0.2);
\coordinate (A) at (2.2,0.2); \coordinate (B) at (2.9,0.6); \coordinate (C) at (2.4,1.3);
\coordinate (Ap) at ({0.9+1.3*cos(60)},{0.2+1.3*sin(60)});
\coordinate (Bp) at ({0.9+2*cos(60)+0.4*cos(60+90)},{0.2+2*sin(60)+0.4*sin(60+90)});
\coordinate (Cp) at ({0.9+1.5*cos(60+28)},{0.2+1.5*sin(60+28)});
\fill[popink] (Om) circle (1.5pt) node[below left] {\footnotesize $\Omega$};
\draw[fill=popblueL,draw=popblue,thick] (A)--(B)--(C)--cycle;
\draw[fill=poporangeL,draw=poporange,thick] (Ap)--(Bp)--(Cp)--cycle;
\draw[popdark,->] (1.9,0.2) arc (0:60:1.0);
\node[popdark] at (2.15,0.75) {\footnotesize $60^\circ$};
\node at (1.8,-0.6) {\footnotesize \textbf{Rotation} ($60^\circ$) : $z' = \mathrm{e}^{\mathrm{i}\pi/3}z + b$};
\end{scope}
\end{tikzpicture}
\legende{Le même triangle $ABC$ (bleu) et son image (orange) par les trois transformations usuelles. Dans les trois cas, la forme du triangle est préservée : c'est le point de départ de la notion de similitude.}
\end{popfigure}

\begin{aretenir}[Les trois écritures complexes à connaître par cœur]
\begin{center}
\begin{tabularx}{\linewidth}{|l|X|X|}
\hline
\rowcolor{popblueL} \textbf{Transformation} & \textbf{Écriture complexe} & \textbf{Forme $z' = az + b$} \\
\hline
Translation de vecteur $\vec{u}(b)$ & $z' = z + b$ & $a = 1$ \\
\hline
Homothétie de centre $\Omega(\omega)$, rapport $k \in \mathbb{R}^*$ & $z' - \omega = k(z - \omega)$ & $a = k$ réel, $b = (1-k)\omega$ \\
\hline
Rotation de centre $\Omega(\omega)$, angle $\theta$ & $z' - \omega = \mathrm{e}^{\mathrm{i}\theta}(z - \omega)$ & $a = \mathrm{e}^{\mathrm{i}\theta}$, $|a| = 1$, $b = \omega(1 - \mathrm{e}^{\mathrm{i}\theta})$ \\
\hline
\end{tabularx}
\end{center}
\end{aretenir}

\begin{savaistu}[Une seule formule pour les gouverner toutes]
Regarde attentivement la colonne de droite du tableau ci-dessus : translation, homothétie, rotation s'écrivent \emph{toutes} sous la forme $z' = az + b$ ! Seul le coefficient $a$ change de nature : $a = 1$ pour la translation, $a$ réel pour l'homothétie, $|a| = 1$ pour la rotation. Que se passe-t-il alors si $a$ est un complexe \emph{quelconque}, par exemple $a = 1 + \mathrm{i}$ ? On obtient une transformation qui agrandit \emph{et} tourne à la fois — exactement ce dont l'artisan de Foumban et le cartographe de Douala ont besoin ! C'est la \cle{similitude directe}, que nous définirons rigoureusement dans la section suivante. Cette unification est l'une des plus belles illustrations de la puissance des nombres complexes, introduits au XVI\textsuperscript{e} siècle par les algébristes italiens Cardan et Bombelli, et devenus, trois siècles plus tard avec Gauss et Argand, un langage géométrique à part entière.
\end{savaistu}

\begin{exercice}[Pour t'entraîner]
Le plan est rapporté à un repère orthonormé direct $(O\,;\vec{u},\vec{v})$.
\begin{enumerate}
\item Déterminer l'écriture complexe de la translation qui envoie $A(1 + 2\mathrm{i})$ sur $B(4 - \mathrm{i})$.
\item Soit $h$ l'homothétie de centre $\Omega(2 - \mathrm{i})$ et de rapport $k = 3$. Déterminer l'image de $M(-1 + \mathrm{i})$.
\item Soit $r$ la rotation de centre $O$ et d'angle $\dfrac{\pi}{2}$. Déterminer l'image de $P(3 + 2\mathrm{i})$, puis vérifier que $OP = OP'$ et que $(\overrightarrow{OP}\,;\overrightarrow{OP'}) = \dfrac{\pi}{2}$.
\item Déterminer les points fixes de la transformation d'écriture complexe $z' = \mathrm{i}z + 1 - \mathrm{i}$.
\end{enumerate}
\end{exercice}

\begin{corrige}[Exercice 1]
\begin{enumerate}
\item Le vecteur de translation est $\overrightarrow{AB}$, d'affixe $z_B - z_A = (4 - \mathrm{i}) - (1 + 2\mathrm{i}) = 3 - 3\mathrm{i}$. D'où $z' = z + 3 - 3\mathrm{i}$.
\item $z' = 3z + (1-3)(2 - \mathrm{i}) = 3z - 4 + 2\mathrm{i}$. Pour $M(-1 + \mathrm{i})$ : $z_{M'} = 3(-1 + \mathrm{i}) - 4 + 2\mathrm{i} = -7 + 5\mathrm{i}$, donc $M'(-7\,;5)$.
\item Rotation de centre $O$ : $z' = \mathrm{e}^{\mathrm{i}\pi/2}z = \mathrm{i}z$. Donc $z_{P'} = \mathrm{i}(3 + 2\mathrm{i}) = -2 + 3\mathrm{i}$, soit $P'(-2\,;3)$. Vérification : $OP = |3 + 2\mathrm{i}| = \sqrt{13}$ et $OP' = |-2 + 3\mathrm{i}| = \sqrt{13}$ ; de plus $\arg\!\left(\dfrac{z_{P'}}{z_P}\right) = \arg(\mathrm{i}) = \dfrac{\pi}{2}$.
\item On résout $z = \mathrm{i}z + 1 - \mathrm{i}$, soit $z(1 - \mathrm{i}) = 1 - \mathrm{i}$, d'où $z = 1$. L'unique point fixe est le point d'affixe $1$ (c'est le centre de la rotation d'angle $\frac{\pi}{2}$).
\end{enumerate}
\end{corrige}

Nous voilà armés : affixes, translations, homothéties, rotations et points fixes n'ont plus de secret pour toi. Dans la section suivante, nous assemblerons ces briques pour construire l'objet central de la leçon : la similitude directe, transformation d'écriture complexe $z' = az + b$ avec $a \neq 0$.

\coursec{Définition et écriture complexe d'une similitude directe}

Dans la section précédente, nous avons rappelé les écritures complexes des trois transformations usuelles du plan : la \cle{translation} ($z' = z + b$), l'\cle{homothétie} de centre $\Omega(\omega)$ et de rapport $k$ ($z' - \omega = k(z - \omega)$) et la \cle{rotation} de centre $\Omega(\omega)$ et d'angle $\theta$ ($z' - \omega = e^{i\theta}(z - \omega)$). Observe ces trois écritures attentivement : elles ont toutes la même forme générale $z' = az + b$, avec un coefficient $a$ particulier dans chaque cas ($a = 1$ pour la translation, $a = k$ réel pour l'homothétie, $a = e^{i\theta}$ de module $1$ pour la rotation). Cette observation n'est pas un hasard : elle est le point de départ de toute cette leçon. Nous allons regrouper toutes ces transformations (et d'autres) sous une même bannière : celle des \cle{similitudes directes}.

\courssub{Intuition géométrique : ce qu'est une similitude directe}

\begin{experience}[Découverte : agrandir sans déformer]
Imagine que tu photographies la carte du Cameroun affichée au mur de ta salle de classe, puis que tu affiches cette photo sur l'écran d'un téléphone en l'agrandissant et en la tournant légèrement. La carte agrandie n'est plus à la même échelle : la distance Yaoundé--Douala est deux fois plus grande sur l'écran que sur le mur. Mais la \emph{forme} du pays est restée identique : tous les angles (l'angle entre la frontière ouest et la frontière nord, par exemple) sont conservés, et aucune région n'est « écrasée » ou « étirée » différemment des autres.

Mathématiquement, on vient de décrire une transformation qui :
\begin{itemize}
\item multiplie \textbf{toutes} les distances par un même réel $k > 0$ (ici $k = 2$) ;
\item conserve les angles, y compris leur orientation (le nord reste « à gauche » de l'est, et non à droite : il n'y a pas d'effet miroir).
\end{itemize}
C'est exactement cela, une similitude directe.
\end{experience}

\begin{definition}[Définition géométrique]
Soit $k$ un réel strictement positif. On appelle \cle{similitude directe} du plan de rapport $k$ toute \cle{transformation} $s$ du plan (application bijective du plan dans lui-même) telle que :
\begin{enumerate}
\item pour tous points $M$ et $N$ d'images $M'$ et $N'$ : $M'N' = k \times MN$ (les distances sont multipliées par $k$) ;
\item pour tous points $A$, $B$, $C$ distincts d'images $A'$, $B'$, $C'$ : $\left(\overrightarrow{A'B'}\,;\,\overrightarrow{A'C'}\right) = \left(\overrightarrow{AB}\,;\,\overrightarrow{AC}\right) \ [2\pi]$ (les angles orientés sont conservés).
\end{enumerate}
\end{definition}

\begin{attention}[Directe ou indirecte ?]
La conservation des angles \emph{orientés} est essentielle. Une symétrie axiale (réflexion) conserve les distances et les valeurs absolues des angles, mais elle \textbf{renverse} l'orientation des angles : ce n'est pas une similitude directe, c'est une similitude \emph{indirecte}, qui n'est pas au programme de Terminale. Retiens : « directe » signifie « qui conserve le sens des angles ».
\end{attention}

\courssub{Définition complexe : la forme $z' = az + b$}

L'outil complexe va nous donner une définition opérationnelle remarquablement simple. Le plan est muni d'un repère orthonormé direct $(O\,;\vec{u},\vec{v})$, et chaque point $M$ est repéré par son affixe $z$.

\begin{definition}[Définition complexe d'une similitude directe]
Soient $a \in \mathbb{C}^{*}$ et $b \in \mathbb{C}$. L'application $s$ du plan dans lui-même qui, à tout point $M$ d'affixe $z$, associe le point $M'$ d'affixe $z'$ définie par
\[ z' = az + b \]
est une \cle{similitude directe}. On dit que $z' = az + b$ est l'\cle{écriture complexe} de $s$.
\end{definition}

\begin{attention}[La condition $a \neq 0$ est vitale]
L'erreur la plus fréquente consiste à oublier la condition $a \in \mathbb{C}^{*}$, c'est-à-dire $a \neq 0$. Si $a = 0$, l'écriture devient $z' = b$ : \textbf{tous} les points du plan ont la même image, le point d'affixe $b$. L'application est constante, elle n'est pas bijective, ce n'est donc \textbf{pas une transformation du plan}, et à plus forte raison pas une similitude. Chaque fois que tu écris « soit $s$ la similitude d'écriture complexe $z' = az + b$ », vérifie mentalement que $a \neq 0$.
\end{attention}

\begin{propriete}[Bijectivité et similitude réciproque]
Toute similitude directe $s$ d'écriture complexe $z' = az + b$ ($a \neq 0$) est une bijection du plan. Sa réciproque $s^{-1}$ est aussi une similitude directe, d'écriture complexe
\[ z' = \frac{1}{a}\,z - \frac{b}{a}, \]
de rapport $\dfrac{1}{|a|}$ et d'angle $-\arg a \ [2\pi]$.
\end{propriete}

\begin{experience}[Démonstration guidée]
\textbf{Étape 1 : $s$ est bijective.} Soit $M'(z')$ un point quelconque du plan. Cherchons ses antécédents : $z' = az + b \iff z = \dfrac{z' - b}{a}$, ce qui est possible car $a \neq 0$. Tout point $M'$ possède donc un antécédent, et un seul : $s$ est bijective.

\textbf{Étape 2 : nature de $s^{-1}$.} L'application réciproque associe à $z'$ le nombre $z = \dfrac{1}{a}z' - \dfrac{b}{a}$. En échangeant les rôles des variables, $s^{-1}$ a pour écriture complexe $z' = \dfrac{1}{a}z - \dfrac{b}{a}$. Comme $a \neq 0$, on a $\dfrac{1}{a} \neq 0$ : c'est bien une similitude directe. Son rapport est $\left|\dfrac{1}{a}\right| = \dfrac{1}{|a|}$ et son angle $\arg\left(\dfrac{1}{a}\right) = -\arg a \ [2\pi]$. \hfill$\blacksquare$
\end{experience}

Vérifions maintenant que la définition complexe réalise bien les deux exigences de la définition géométrique.

\begin{propriete}[Effet sur les distances et les angles]
Soit $s$ la similitude directe d'écriture complexe $z' = az + b$ ($a \neq 0$). Alors :
\begin{enumerate}
\item $s$ multiplie les distances par $|a|$ : pour tous points $M$, $N$ d'images $M'$, $N'$, on a $M'N' = |a| \times MN$ ;
\item $s$ conserve les angles orientés : pour tous points $A$, $B$, $C$ distincts d'images $A'$, $B'$, $C'$,
\[ \left(\overrightarrow{A'B'}\,;\,\overrightarrow{A'C'}\right) = \left(\overrightarrow{AB}\,;\,\overrightarrow{AC}\right) \ [2\pi]. \]
\end{enumerate}
\end{propriete}

\begin{experience}[Démonstration guidée]
Soient $M(z_M)$, $N(z_N)$, $A(z_A)$, $B(z_B)$, $C(z_C)$ des points et $M'(z'_M)$, $N'(z'_N)$, \dots leurs images.

\textbf{Étape 1 : les distances.} Exprime les affixes des images : $z'_M = a z_M + b$ et $z'_N = a z_N + b$. Par soustraction, le $b$ disparaît :
\[ z'_N - z'_M = a(z_N - z_M). \]
Passe aux modules, en utilisant $|a \times w| = |a|\,|w|$ :
\[ M'N' = |z'_N - z'_M| = |a|\,|z_N - z_M| = |a| \times MN. \]
Les distances sont bien multipliées par $|a|$, qui est strictement positif puisque $a \neq 0$.

\textbf{Étape 2 : les angles orientés.} De même, $z'_{B} - z'_{A} = a(z_B - z_A)$ et $z'_{C} - z'_{A} = a(z_C - z_A)$. Comme $A$, $B$, $C$ sont distincts et que $s$ est bijective, $A'$, $B'$, $C'$ sont distincts, et ces différences sont non nulles : on peut former le rapport :
\[ \frac{z'_{C} - z'_{A}}{z'_{B} - z'_{A}} = \frac{a(z_C - z_A)}{a(z_B - z_A)} = \frac{z_C - z_A}{z_B - z_A}. \]
Le facteur $a$ se simplifie ! En passant aux arguments :
\[ \left(\overrightarrow{A'B'}\,;\,\overrightarrow{A'C'}\right) = \arg\left(\frac{z'_{C} - z'_{A}}{z'_{B} - z'_{A}}\right) = \arg\left(\frac{z_C - z_A}{z_B - z_A}\right) = \left(\overrightarrow{AB}\,;\,\overrightarrow{AC}\right) \ [2\pi]. \]
Les angles orientés sont conservés. \hfill$\blacksquare$
\end{experience}

\begin{aretenir}[L'essentiel en une phrase]
Une similitude directe, c'est une écriture $z' = az + b$ avec $a \neq 0$ : le coefficient $a$ « pilote » le rapport ($|a|$) et l'angle ($\arg a$), le coefficient $b$ « pilote » le déplacement global de la figure.
\end{aretenir}

\courssub{La réciproque : toute similitude directe s'écrit $z' = az + b$}

Nous venons de voir que toute application $z \mapsto az + b$ (avec $a \neq 0$) est une similitude directe. La réciproque est vraie et fondamentale.

\begin{propriete}[Écriture complexe de toute similitude directe (admise)]
Toute similitude directe du plan admet, dans tout repère orthonormé direct, une écriture complexe de la forme $z' = az + b$, avec $a \in \mathbb{C}^{*}$ et $b \in \mathbb{C}$, et cette écriture est unique.
\end{propriete}

\begin{experience}[Justification guidée (idée de la preuve)]
Soit $s$ une similitude directe de rapport $k$. L'idée est de « ramener » $s$ à une isométrie.

\textbf{Étape 1.} Soit $h$ l'homothétie de centre $O$ et de rapport $\frac{1}{k}$, d'écriture complexe $z' = \frac{1}{k}z$. Considère la composée $f = h \circ s$. Pour tous points $M$, $N$ : la similitude $s$ multiplie les distances par $k$, puis $h$ les multiplie par $\frac{1}{k}$. Donc $f$ conserve les distances : c'est une \cle{isométrie}.

\textbf{Étape 2.} Comme $s$ et $h$ conservent toutes deux les angles orientés, $f$ les conserve aussi : $f$ est un \emph{déplacement} (isométrie directe). Or on sait (cours sur les isométries) que tout déplacement du plan est une translation ou une rotation, et admet donc une écriture complexe $z' = \alpha z + \beta$ avec $|\alpha| = 1$.

\textbf{Étape 3.} De $f = h \circ s$, on tire $s = h^{-1} \circ f$. Or $h^{-1}$ est l'homothétie de rapport $k$, d'écriture $z' = kz$. En composant les écritures :
\[ z \xmapsto{\ f\ } \alpha z + \beta \xmapsto{\ h^{-1}\ } k(\alpha z + \beta) = (k\alpha) z + k\beta. \]
En posant $a = k\alpha$ (non nul, car $|a| = k|\alpha| = k > 0$) et $b = k\beta$, on obtient bien $z' = az + b$. \hfill$\blacksquare$
\end{experience}

\courssub{Les cas particuliers : retrouver nos transformations usuelles}

Toute la richesse de la forme $z' = az + b$ tient dans le coefficient $a$. Selon sa valeur, on retrouve les transformations étudiées dans la section précédente.

\begin{propriete}[Classification selon la valeur de $a$]
Soit $s$ la similitude directe d'écriture complexe $z' = az + b$, $a \in \mathbb{C}^{*}$.
\begin{enumerate}
\item Si $a = 1$ : $z' = z + b$, c'est la \cle{translation} de vecteur $\vec{w}$ d'affixe $b$.
\item Si $a$ est réel et $a \neq 1$ : c'est une \cle{homothétie} de rapport $a$ (et de centre le point fixe de $s$).
\item Si $|a| = 1$ et $a \neq 1$ : en écrivant $a = e^{i\theta}$, c'est une \cle{rotation} d'angle $\theta = \arg a \ [2\pi]$ (et de centre le point fixe de $s$).
\item Si $|a| \neq 1$ et $a$ non réel : c'est une similitude « véritable », composée d'une homothétie de rapport $|a|$ et d'une rotation d'angle $\arg a$, de même centre.
\end{enumerate}
Dans tous les cas, $|a|$ est le rapport de la similitude et, lorsque $a \neq 1$, $\arg a$ est son angle.
\end{propriete}

\begin{popfigure}
\begin{center}
\begin{tikzpicture}[grow=right, level distance=4.2cm,
  edge from parent/.style={draw=popmuted, thick},
  every node/.style={font=\small},
  level 1/.style={sibling distance=2.6cm},
  level 2/.style={sibling distance=1.6cm}]
\node[draw=popdark, fill=popblueL, rounded corners, thick, align=center] {$z' = az + b$\\ $a \in \mathbb{C}^{*},\ b \in \mathbb{C}$}
  child { node[draw=poporange, fill=poporangeL, rounded corners, align=center] {$|a| \neq 1$\\ $a \notin \mathbb{R}$} 
    child { node[draw=popdark, fill=popcream, rounded corners, align=center] {Similitude\\ véritable} } }
  child { node[draw=poppurple, fill=poppurpleL, rounded corners, align=center] {$|a| = 1$\\ $a \neq 1$}
    child { node[draw=popdark, fill=popcream, rounded corners, align=center] {Rotation\\ d'angle $\arg a$} } }
  child { node[draw=popgreen, fill=popgreenL, rounded corners, align=center] {$a \in \mathbb{R}$\\ $a \neq 1$}
    child { node[draw=popdark, fill=popcream, rounded corners, align=center] {Homothétie\\ de rapport $a$} } }
  child { node[draw=popblue, fill=popblueL, rounded corners, align=center] {$a = 1$}
    child { node[draw=popdark, fill=popcream, rounded corners, align=center] {Translation\\ de vecteur $b$} } };
\end{tikzpicture}
\end{center}
\legende{Classification des similitudes directes $z' = az + b$ selon la valeur du coefficient $a$.}
\end{popfigure}

\begin{exemplebox}[Reconnaître une transformation]
L'application $s : M(z) \mapsto M'(z')$ avec $z' = -3z + 2 - i$ est une similitude directe car $a = -3 \neq 0$. Comme $a$ est réel et différent de $1$, $s$ est une homothétie de rapport $-3$. Son centre est l'unique point fixe : $\omega = -3\omega + 2 - i$, soit $4\omega = 2 - i$, donc $\omega = \dfrac{1}{2} - \dfrac{1}{4}i$.
\end{exemplebox}

\courssub{Expression analytique d'une similitude directe}

L'écriture complexe est compacte et élégante, mais dans certains exercices (notamment ceux mêlant géométrie repérée classique), on préfère travailler avec les coordonnées cartésiennes. Le passage de l'une à l'autre est un savoir-faire exigible.

\begin{propriete}[Expression analytique]
Soit $s$ la similitude directe d'écriture complexe $z' = az + b$, avec $a = \alpha + i\beta$ ($\alpha, \beta \in \mathbb{R}$, $(\alpha, \beta) \neq (0,0)$) et $b = b_1 + i b_2$ ($b_1, b_2 \in \mathbb{R}$). Alors le point $M(x\,;y)$ a pour image le point $M'(x'\,;y')$ avec :
\[ \begin{cases} x' = \alpha x - \beta y + b_1 \\ y' = \beta x + \alpha y + b_2 \end{cases} \]
C'est l'\cle{expression analytique} de $s$.
\end{propriete}

\begin{methode}[Passer de l'écriture complexe à l'expression analytique]
\begin{enumerate}
\item Pose $a = \alpha + i\beta$, $b = b_1 + ib_2$, $z = x + iy$ et $z' = x' + iy'$.
\item Développe le produit : $z' = (\alpha + i\beta)(x + iy) + b_1 + ib_2 = \alpha x + i\alpha y + i\beta x + i^2 \beta y + b_1 + ib_2$.
\item Utilise $i^2 = -1$ et regroupe parties réelle et imaginaire : $z' = (\alpha x - \beta y + b_1) + i(\beta x + \alpha y + b_2)$.
\item Identifie : $x'$ est la partie réelle, $y'$ la partie imaginaire.
\end{enumerate}
Pour le sens inverse (analytique $\to$ complexe), on procède par identification : le coefficient de $x$ dans $x'$ et celui de $y$ dans $y'$ donnent $\alpha$ ; le coefficient de $y$ dans $x'$ donne $-\beta$ ; les constantes donnent $b_1$ et $b_2$. On reconstitue alors $a = \alpha + i\beta$ et $b = b_1 + ib_2$, puis on écrit $z' = az + b$.
\end{methode}

\begin{exoresolu}[Expression analytique de $z' = (1 - i)z + 2 + 3i$]
Soit $s$ la transformation du plan d'écriture complexe $z' = (1 - i)z + 2 + 3i$. Déterminer l'expression analytique de $s$.
\tcblower
\textbf{Solution détaillée.} On applique la méthode pas à pas.

\textbf{Étape 1.} On identifie $a = 1 - i$, donc $\alpha = 1$ et $\beta = -1$ ; et $b = 2 + 3i$, donc $b_1 = 2$ et $b_2 = 3$. On pose $z = x + iy$ et $z' = x' + iy'$.

\textbf{Étape 2.} Développons :
\begin{align*}
z' &= (1 - i)(x + iy) + 2 + 3i \\
&= x + iy - ix - i^2 y + 2 + 3i \\
&= x + iy - ix + y + 2 + 3i \quad (\text{car } i^2 = -1).
\end{align*}

\textbf{Étape 3.} Regroupons parties réelle et imaginaire :
\[ z' = (x + y + 2) + i(y - x + 3). \]

\textbf{Étape 4.} Par identification de $z' = x' + iy'$ :
\[ \begin{cases} x' = x + y + 2 \\ y' = -x + y + 3 \end{cases} \]

\textbf{Vérification.} Le point $O(0\,;0)$ a pour image le point d'affixe $b = 2 + 3i$, c'est-à-dire $(2\,;3)$. La formule analytique donne $x' = 0 + 0 + 2 = 2$ et $y' = -0 + 0 + 3 = 3$. Cohérent. \hfill$\blacksquare$
\end{exoresolu}

\begin{attention}[Pièges de l'identification]
\begin{itemize}
\item Attention au signe : dans $x' = \alpha x - \beta y + b_1$, c'est \textbf{moins} $\beta$ devant $y$, car $i^2 = -1$. Beaucoup d'élèves écrivent $x' = \alpha x + \beta y + b_1$ par étourderie.
\item Dans $y'$, l'ordre est inversé : c'est $\beta x + \alpha y$, pas $\alpha x + \beta y$. Une astuce de mémorisation : la matrice $\begin{pmatrix} \alpha & -\beta \\ \beta & \alpha \end{pmatrix}$ est « anti-symétrique » en ses coefficients.
\item Ne confonds pas $b$ (affixe complexe) avec ses coordonnées $b_1$ et $b_2$.
\end{itemize}
\end{attention}

\courssub{Une figure pour visualiser : l'exemple de $z' = (1+i)z$}

Étudions en détail la similitude $s : z' = (1+i)z$ (ici $b = 0$). On a $a = 1 + i$, donc :
\[ |a| = \sqrt{1^2 + 1^2} = \sqrt{2} \qquad \text{et} \qquad \arg(a) = \frac{\pi}{4} \ [2\pi]. \]
La similitude $s$ est donc la composée d'une homothétie de rapport $\sqrt{2}$ et d'une rotation d'angle $\frac{\pi}{4}$, toutes deux de centre $O$ (car $b = 0$, le point $O$ est fixe). Prenons le triangle $ABC$ avec $A(1\,;0)$, $B(2\,;0)$, $C(2\,;1)$, d'affixes $z_A = 1$, $z_B = 2$, $z_C = 2 + i$. Les images sont :
\begin{align*}
z_{A'} &= (1+i)(1) = 1 + i, \\
z_{B'} &= (1+i)(2) = 2 + 2i, \\
z_{C'} &= (1+i)(2+i) = 2 + 2i + i + i^2 = 1 + 3i.
\end{align*}
Donc $A'(1\,;1)$, $B'(2\,;2)$, $C'(1\,;3)$. Vérifions sur une distance : $AB = 1$ et $A'B' = |z_{B'} - z_{A'}| = |1 + i| = \sqrt{2} = \sqrt{2}\,AB$. Parfait.

\begin{popfigure}
\begin{center}
\begin{tikzpicture}[scale=1.15, >=Stealth]
% Grille et axes
\draw[popline, very thin, step=1] (-0.5,-0.5) grid (3.4,3.4);
\draw[->, popdark, thick] (-0.5,0) -- (3.6,0) node[below left] {$x$};
\draw[->, popdark, thick] (0,-0.5) -- (0,3.6) node[below left] {$y$};
\node[below left, popdark] at (0,0) {$O$};
% Triangle ABC
\coordinate (A) at (1,0);
\coordinate (B) at (2,0);
\coordinate (C) at (2,1);
\fill[popblueL, opacity=0.6] (A) -- (B) -- (C) -- cycle;
\draw[popblue, very thick] (A) -- (B) -- (C) -- cycle;
\node[below, popblue] at (A) {$A(1\,;0)$};
\node[below right, popblue] at (B) {$B(2\,;0)$};
\node[right, popblue] at (C) {$C(2\,;1)$};
% Triangle A'B'C'
\coordinate (A2) at (1,1);
\coordinate (B2) at (2,2);
\coordinate (C2) at (1,3);
\fill[poporangeL, opacity=0.6] (A2) -- (B2) -- (C2) -- cycle;
\draw[poporange, very thick] (A2) -- (B2) -- (C2) -- cycle;
\node[left, poporange] at (A2) {$A'(1\,;1)$};
\node[right, poporange] at (B2) {$B'(2\,;2)$};
\node[left, poporange] at (C2) {$C'(1\,;3)$};
% Arcs de rotation depuis O
\draw[->, poppink, thick] (1.7,0) arc (0:45:1.7) node[midway, right, xshift=2pt] {\small $\frac{\pi}{4}$};
\draw[poppink, dashed, thick] (0,0) -- (B);
\draw[poppink, dashed, thick] (0,0) -- (B2);
% Légende rapport
\node[popdark, align=center, fill=popcream, rounded corners, draw=popmuted] at (2.85,0.85) {\small rapport $\sqrt{2}$\\ \small angle $\frac{\pi}{4}$};
\end{tikzpicture}
\end{center}
\legende{Le triangle $ABC$ (bleu) et son image $A'B'C'$ (orange) par la similitude $z' = (1+i)z$ : chaque point tourne de $\frac{\pi}{4}$ autour de $O$ et sa distance à $O$ est multipliée par $\sqrt{2}$.}
\end{popfigure}

Observe la figure : le triangle image a exactement la même \emph{forme} que le triangle initial (les angles sont conservés, en particulier l'angle droit en $B$ se retrouve en $B'$), mais il est agrandi d'un facteur $\sqrt{2}$ et pivoté de $45^{\circ}$. C'est toute la philosophie des similitudes.

\courssub{Composée de deux similitudes directes}

Que se passe-t-il lorsqu'on enchaîne deux similitudes directes ? Par exemple, un agrandissement suivi d'une rotation ? Le calcul complexe répond en une ligne.

\begin{propriete}[Composée de deux similitudes directes]
Soient $s_1$ et $s_2$ deux similitudes directes d'écritures complexes respectives $z' = a_1 z + b_1$ et $z' = a_2 z + b_2$ ($a_1, a_2 \in \mathbb{C}^{*}$). Alors la composée $s_2 \circ s_1$ est une similitude directe d'écriture complexe
\[ z' = (a_2 a_1)\, z + (a_2 b_1 + b_2). \]
En particulier :
\begin{itemize}
\item son \cle{rapport} est le \textbf{produit} des rapports : $k = |a_2 a_1| = |a_2| \times |a_1|$ ;
\item son \cle{angle} est la \textbf{somme} des angles : $\theta = \arg(a_2 a_1) = \arg a_1 + \arg a_2 \ [2\pi]$.
\end{itemize}
L'ensemble des similitudes directes du plan est donc stable par composition.
\end{propriete}

\begin{experience}[Démonstration guidée]
Soit $M(z)$ un point quelconque. Notons $M_1(z_1)$ son image par $s_1$, puis $M'(z')$ l'image de $M_1$ par $s_2$.

\textbf{Étape 1.} Par définition de $s_1$ : $z_1 = a_1 z + b_1$.

\textbf{Étape 2.} Par définition de $s_2$ appliquée à $M_1$ : $z' = a_2 z_1 + b_2$.

\textbf{Étape 3.} Substitue l'expression de $z_1$ et développe :
\begin{align*}
z' &= a_2 (a_1 z + b_1) + b_2 \\
&= a_2 a_1 z + a_2 b_1 + b_2 \\
&= (a_2 a_1)\, z + (a_2 b_1 + b_2).
\end{align*}

\textbf{Étape 4.} Pose $a = a_2 a_1$ et $b = a_2 b_1 + b_2$. Comme $a_1 \neq 0$ et $a_2 \neq 0$, on a $a \neq 0$ (un produit de complexes non nuls est non nul). L'écriture $z' = az + b$ est donc bien celle d'une similitude directe. Les propriétés du module et de l'argument d'un produit donnent immédiatement le rapport produit et l'angle somme. \hfill$\blacksquare$
\end{experience}

\begin{attention}[L'ordre de composition compte]
En général, $s_2 \circ s_1 \neq s_1 \circ s_2$ : la composition des similitudes n'est \textbf{pas commutative}. Le rapport ($k_1 k_2$) et l'angle ($\theta_1 + \theta_2$) sont les mêmes dans les deux ordres, mais le terme $b$ change : pour $s_2 \circ s_1$ c'est $a_2 b_1 + b_2$, pour $s_1 \circ s_2$ c'est $a_1 b_2 + b_1$. Le centre de la similitude composée peut donc différer. Au Baccalauréat, lis toujours soigneusement l'ordre indiqué dans l'énoncé : dans $s_2 \circ s_1$, c'est $s_1$ qui agit \emph{en premier}.
\end{attention}

\begin{exemplebox}[Composer une homothétie et une rotation]
Soient $h$ l'homothétie d'écriture $z' = 2z + i$ et $r$ la rotation d'écriture $z' = iz + 1$. Déterminons la nature de $r \circ h$. On a $a_1 = 2$, $b_1 = i$, $a_2 = i$, $b_2 = 1$, donc :
\[ z' = i(2z + i) + 1 = 2iz + i^2 + 1 = 2iz. \]
La composée $r \circ h$ est la similitude d'écriture $z' = 2iz$ : rapport $|2i| = 2$, angle $\arg(2i) = \frac{\pi}{2}$, et elle fixe $O$ (car $b = 0$). C'est la composée commutative de l'homothétie de centre $O$ et de rapport $2$ et de la rotation de centre $O$ et d'angle $\frac{\pi}{2}$.
\end{exemplebox}

\begin{savaistu}[Pourquoi « similitude » ?]
Le mot vient du latin \emph{similis}, « semblable ». Deux figures sont dites \emph{semblables} lorsque l'une est l'image de l'autre par une similitude : elles ont la même forme, mais pas nécessairement la même taille. C'est exactement la notion que tu as rencontrée au collège avec les \emph{triangles semblables} (triangles ayant les mêmes angles) et les agrandissements/réductions. Les similitudes sont aussi au cœur des \emph{fractales} : le flocon de Koch ou l'ensemble de Mandelbrot sont constitués de motifs qui se répètent à toutes les échelles, chaque copie étant l'image du motif initial par une similitude. Les ingénieurs télécoms utilisent d'ailleurs des antennes fractales, dont la géométrie auto-similaire permet de capter plusieurs fréquences avec un seul dispositif compact !
\end{savaistu}

\begin{exercice}[À toi de jouer]
On considère les transformations $s_1$ et $s_2$ d'écritures complexes respectives $z' = (2 + 2i)z - 1$ et $z' = \dfrac{1}{2}i\, z + 3 - i$.
\begin{enumerate}
\item Donner l'expression analytique de $s_1$.
\item Déterminer l'écriture complexe de $s_2 \circ s_1$, puis préciser son rapport et son angle.
\item La composée $s_2 \circ s_1$ est-elle une transformation usuelle ? Justifier.
\end{enumerate}
\end{exercice}

\begin{corrige}[Exercice : À toi de jouer]
\textbf{1.} Avec $a = 2 + 2i$ ($\alpha = 2$, $\beta = 2$) et $b = -1$ ($b_1 = -1$, $b_2 = 0$) :
\[ z' = (2+2i)(x+iy) - 1 = 2x + 2iy + 2ix - 2y - 1 = (2x - 2y - 1) + i(2x + 2y). \]
Donc $\begin{cases} x' = 2x - 2y - 1 \\ y' = 2x + 2y \end{cases}$.

\textbf{2.} Composée $s_2 \circ s_1$ : $a = \dfrac{i}{2}(2 + 2i) = i + i^2 = -1 + i$ et $b = \dfrac{i}{2}(-1) + 3 - i = 3 - \dfrac{3}{2}i$. D'où $z' = (-1+i)z + 3 - \dfrac{3}{2}i$. Rapport : $|-1+i| = \sqrt{2}$. Angle : $\arg(-1+i) = \dfrac{3\pi}{4} \ [2\pi]$.

\textbf{3.} Le rapport $\sqrt{2} \neq 1$ et $a = -1 + i$ n'est pas réel : ce n'est ni une translation, ni une homothétie, ni une rotation. C'est une similitude directe « véritable », de rapport $\sqrt{2}$ et d'angle $\dfrac{3\pi}{4}$.
\end{corrige}

\begin{aretenir}[Bilan de la section]
\begin{itemize}
\item Une \cle{similitude directe} est une transformation du plan d'écriture complexe $z' = az + b$ avec $a \in \mathbb{C}^{*}$ et $b \in \mathbb{C}$.
\item Elle multiplie les distances par $|a|$ et conserve les angles orientés.
\item Cas particuliers : $a = 1$ (translation), $a$ réel $\neq 1$ (homothétie), $|a| = 1$ et $a \neq 1$ (rotation).
\item Expression analytique : $\begin{cases} x' = \alpha x - \beta y + b_1 \\ y' = \beta x + \alpha y + b_2 \end{cases}$ où $a = \alpha + i\beta$ et $b = b_1 + ib_2$.
\item La composée de deux similitudes directes est une similitude directe de rapport produit et d'angle somme ; la réciproque d'une similitude directe est une similitude directe.
\end{itemize}
Dans la section suivante, nous apprendrons à déterminer les \emph{éléments caractéristiques} (centre, rapport, angle) d'une similitude directe quelconque à partir de son écriture complexe.
\end{aretenir}

\coursec{Éléments caractéristiques : centre, rapport, angle}

Dans la section précédente, nous avons défini une similitude directe comme une transformation du plan dont l'écriture complexe est de la forme $z' = az + b$, avec $a \in \mathbb{C}^*$ et $b \in \mathbb{C}$. Nous avons aussi remarqué que les cas $a = 1$ (translation) et $b = 0$ avec $a$ réel positif (homothétie) ou $|a| = 1$ (rotation) étaient déjà connus. La grande question est maintenant la suivante : \emph{comment lire géométriquement une écriture complexe $z' = az + b$ quelconque ?} Autrement dit, quels nombres, cachés dans $a$ et $b$, décrivent complètement l'action de la similitude sur le plan ?

L'intuition est la suivante. Une translation « glisse » tout le plan : aucun point ne reste en place. Mais dès que la similitude tourne ou dilate, on s'attend à ce qu'il existe un point \emph{pivot}, immobile, autour duquel tout le plan tourne et se dilate à la fois — comme une photo que l'on agrandit en la faisant pivoter autour d'une punaise plantée en son centre. Ce point privilégié, c'est le \cle{centre} de la similitude. Toute cette section consiste à le débusquer, puis à en déduire le \cle{rapport} et l'\cle{angle} de la similitude.

\courssub{Le théorème du point fixe}

\begin{propriete}[Théorème fondamental — existence et unicité du centre]
Soit $s$ une similitude directe d'écriture complexe $z' = az + b$, avec $a \neq 1$. Alors $s$ admet un \textbf{unique point fixe} $\Omega$, c'est-à-dire un unique point tel que $s(\Omega) = \Omega$. Son affixe est
\[ \omega = \frac{b}{1 - a}. \]
Ce point $\Omega$ est appelé le \cle{centre} de la similitude. De plus, l'écriture complexe de $s$ se met sous la \textbf{forme réduite} :
\[ z' - \omega = a\,(z - \omega). \]
\emph{(Démonstration exigible au Baccalauréat.)}
\end{propriete}

\begin{experience}[Démonstration guidée du théorème du point fixe]
Cherchons les points fixes de $s$, c'est-à-dire les points $M$ d'affixe $z$ tels que $s(M) = M$.
\begin{enumerate}
\item \textbf{Traduire la condition.} $M$ est fixe si et seulement si son affixe $z$ vérifie $z = az + b$.
\item \textbf{Résoudre l'équation.} On regroupe les termes en $z$ :
\[ z - az = b \quad \Longleftrightarrow \quad z(1 - a) = b. \]
\item \textbf{Conclure grâce à l'hypothèse $a \neq 1$.} Comme $a \neq 1$, le nombre $1 - a$ est non nul, donc on peut diviser : l'équation admet une \textbf{unique} solution
\[ \omega = \frac{b}{1 - a}. \]
Il existe donc un unique point fixe $\Omega$, d'affixe $\omega$.
\item \textbf{Établir la forme réduite.} Pour tout point $M$ d'affixe $z$ et son image $M'$ d'affixe $z'$, on a le système
\[ \begin{cases} z' = az + b \\ \omega = a\omega + b \end{cases} \]
(la seconde ligne exprime que $\Omega$ est fixe). En soustrayant membre à membre, $b$ disparaît :
\[ z' - \omega = a(z - \omega). \]
\end{enumerate}
La démonstration est terminée. Remarque au passage : si $a = 1$, l'équation $z(1-a) = b$ devient $0 = b$ ; elle n'a aucune solution si $b \neq 0$ (la translation n'a aucun point fixe) et en a une infinité si $b = 0$ (l'identité). C'est pourquoi le cas $a = 1$ est exclu de l'énoncé.
\end{experience}

\begin{aretenir}[Forme réduite — la clé de toute la leçon]
Toute similitude directe qui n'est pas une translation s'écrit, par rapport à son centre $\Omega$ :
\[ z' - \omega = a\,(z - \omega). \]
Géométriquement, cette égalité dit que le vecteur $\overrightarrow{\Omega M'}$ s'obtient à partir du vecteur $\overrightarrow{\Omega M}$ en multipliant son affixe par $a$. Toute la géométrie de la similitude est donc contenue dans le seul nombre $a$ : son module dilate, son argument fait tourner.
\end{aretenir}

\courssub{Rapport et angle d'une similitude directe}

Écrivons le nombre $a$ sous forme trigonométrique : $a = k\,e^{i\theta}$ avec $k = |a| > 0$ et $\theta = \arg(a) \pmod{2\pi}$. La forme réduite devient
\[ z' - \omega = k\,e^{i\theta}\,(z - \omega). \]
Traduisons cette égalité en distances et en angles. En passant aux modules :
\[ \Omega M' = |z' - \omega| = |a|\cdot|z - \omega| = k \cdot \Omega M. \]
En passant aux arguments, pour $M \neq \Omega$ :
\[ \left(\overrightarrow{\Omega M},\, \overrightarrow{\Omega M'}\right) = \arg\!\left(\frac{z' - \omega}{z - \omega}\right) = \arg(a) = \theta \pmod{2\pi}. \]

\begin{definition}[Rapport et angle d'une similitude directe]
Soit $s$ la similitude directe d'écriture $z' = az + b$ avec $a \neq 1$, de centre $\Omega$.
\begin{itemize}
\item Le \cle{rapport} de $s$ est le réel strictement positif $k = |a|$. Pour tout point $M$ d'image $M'$ : $\Omega M' = k \cdot \Omega M$.
\item L'\cle{angle} de $s$ est le réel $\theta = \arg(a) \pmod{2\pi}$. Pour tout point $M$ distinct de $\Omega$, d'image $M'$ : $\left(\overrightarrow{\Omega M}, \overrightarrow{\Omega M'}\right) = \theta \pmod{2\pi}$.
\end{itemize}
Le triplet $(\Omega, k, \theta)$ — centre, rapport, angle — constitue les \cle{éléments caractéristiques} de la similitude. Ils la déterminent entièrement.
\end{definition}

\begin{attention}[Le piège du module et de l'argument de $a$]
Le rapport doit être un réel \textbf{strictement positif}. Si $a$ est un réel négatif, par exemple $a = -2$, il est \textbf{faux} de dire « rapport $-2$, angle $0$ ». Il faut écrire $a$ sous forme trigonométrique :
\[ a = -2 = 2\,e^{i\pi}, \]
donc le rapport est $k = 2$ et l'angle est $\theta = \pi$. De même, si $a = -3i$, alors $a = 3e^{-i\pi/2}$ : rapport $3$, angle $-\dfrac{\pi}{2}$. Règle d'or : \textbf{toujours} mettre $a$ sous la forme $k e^{i\theta}$ avec $k > 0$ avant de lire le rapport et l'angle.
\end{attention}

\begin{exemplebox}[Lecture directe des éléments caractéristiques]
Soit $s$ la transformation d'écriture complexe $z' = i\sqrt{3}\,z + 2 - 2i\sqrt{3}$.
\begin{itemize}
\item Ici $a = i\sqrt{3} \neq 1$ : $s$ n'est pas une translation, c'est une similitude de centre $\Omega$ d'affixe
\[ \omega = \frac{b}{1 - a} = \frac{2 - 2i\sqrt{3}}{1 - i\sqrt{3}} = \frac{(2 - 2i\sqrt{3})(1 + i\sqrt{3})}{1^2 + (\sqrt{3})^2} = \frac{2 + 2i\sqrt{3} - 2i\sqrt{3} + 6}{4} = \frac{8}{4} = 2. \]
\item Rapport : $k = |a| = |i\sqrt{3}| = \sqrt{3}$.
\item Angle : $\theta = \arg(i\sqrt{3}) = \dfrac{\pi}{2} \pmod{2\pi}$.
\end{itemize}
Conclusion : $s$ est la similitude directe de centre $\Omega(2)$, de rapport $\sqrt{3}$ et d'angle $\dfrac{\pi}{2}$.
\end{exemplebox}

\begin{popfigure}
\begin{center}
\begin{tikzpicture}[scale=1.05, >=Stealth]
\coordinate (O) at (0,0);
\coordinate (M) at (3.4,0);
\coordinate (Mp) at (1.47,2.55);
\coordinate (Mpp) at (-1.91,1.91);
\coordinate (Mppp) at (-2.7,-1.1);
\fill[popink] (O) circle (2pt) node[below left] {$\Omega$};
\fill[popblue] (M) circle (2pt) node[right] {$M$};
\fill[poporange] (Mp) circle (2pt) node[above right] {$M'$};
\fill[poppurple] (Mpp) circle (2pt) node[above left] {$M''$};
\fill[popgreen] (Mppp) circle (2pt) node[below left] {$M'''$};
\draw[popline] (O) -- (M);
\draw[popline] (O) -- (Mp);
\draw[popline] (O) -- (Mpp);
\draw[popline] (O) -- (Mppp);
\draw[->, popblue, thick] (M) .. controls (2.9,1.3) and (2.3,2.2) .. (Mp);
\draw[->, poporange, thick] (Mp) .. controls (0.4,2.9) and (-0.9,2.6) .. (Mpp);
\draw[->, poppurple, thick] (Mpp) .. controls (-2.9,1.2) and (-3.1,0.2) .. (Mppp);
\draw[->, popdark] (1.1,0) arc (0:60:1.1) node[midway, right] {$\theta$};
\node[popmuted] at (4.1,2.6) {$\Omega M' = k \cdot \Omega M$};
\end{tikzpicture}
\end{center}
\legende{Similitude de centre $\Omega$ : chaque point tourne de l'angle $\theta$ et s'éloigne (ou se rapproche) du centre selon le rapport $k$. Les images successives $M, M', M'', M'''$ dessinent une spirale.}
\end{popfigure}

\courssub{Décomposition : homothétie et rotation de même centre}

La forme réduite $z' - \omega = k e^{i\theta}(z - \omega)$ peut se lire en deux temps :
\begin{itemize}
\item multiplier par $k$ : c'est l'action de l'homothétie $h(\Omega, k)$ ;
\item multiplier par $e^{i\theta}$ : c'est l'action de la rotation $r(\Omega, \theta)$.
\end{itemize}
Comme la multiplication des nombres complexes est commutative ($k \times e^{i\theta} = e^{i\theta} \times k$), l'ordre des deux opérations n'a aucune importance.

\begin{definition}[Similitude définie par ses éléments caractéristiques]
Soit $\Omega$ un point du plan, $k$ un réel strictement positif et $\theta$ un réel. La \cle{similitude directe} de centre $\Omega$, de rapport $k$ et d'angle $\theta$ est la composée
\[ s = r(\Omega, \theta) \circ h(\Omega, k) = h(\Omega, k) \circ r(\Omega, \theta). \]
L'homothétie et la rotation de même centre \textbf{commutent} : on peut dilater puis tourner, ou tourner puis dilater, le résultat est identique. Son écriture complexe est
\[ z' - \omega = k e^{i\theta}\,(z - \omega). \]
\end{definition}

\begin{popfigure}
\begin{center}
\begin{tikzpicture}[scale=1.0, >=Stealth]
\coordinate (O) at (0,0);
\coordinate (M) at (3.2,0);
\coordinate (H) at (4.6,0);
\coordinate (Mp) at (2.12,2.12);
\coordinate (R) at (1.48,1.48);
\fill[popink] (O) circle (2pt) node[below left] {$\Omega$};
\fill[popblue] (M) circle (2pt) node[below] {$M$};
\fill[poporange] (Mp) circle (2pt) node[above] {$M'$};
\fill[popgreen] (H) circle (2pt) node[below right] {$M_1 = h(M)$};
\fill[poppurple] (R) circle (2pt) node[left] {$M_2 = r(M)$};
\draw[->, popgreen, thick] (M) -- (H) node[midway, below] {$h(\Omega, k)$};
\draw[->, popgreen, thick, dashed] (H) -- (Mp) node[midway, right] {$r(\Omega, \theta)$};
\draw[->, poppurple, thick] (M) -- (R) node[midway, above left] {$r(\Omega, \theta)$};
\draw[->, poppurple, thick, dashed] (R) -- (Mp) node[midway, above left] {$h(\Omega, k)$};
\draw[popline] (O) -- (M);
\draw[popline] (O) -- (Mp);
\draw[->, popdark] (0.9,0) arc (0:45:0.9) node[midway, right] {$\theta$};
\end{tikzpicture}
\end{center}
\legende{Les deux chemins de $M$ vers $M'$ : d'abord l'homothétie puis la rotation (chemin vert), ou d'abord la rotation puis l'homothétie (chemin violet). Le résultat est le même : $s = r \circ h = h \circ r$.}
\end{popfigure}

\begin{savaistu}[Pourquoi « spirale » ?]
En appliquant plusieurs fois une similitude de rapport $k \neq 1$ à un point $M$, on obtient des points $M, M', M'', \dots$ alignés sur une courbe appelée \emph{spirale logarithmique}. Cette spirale, étudiée par Descartes puis par Jacques Bernoulli (qui la surnommait \emph{spira mirabilis}, la spirale merveilleuse), apparaît dans la nature : coquilles de nautile, bras des galaxies, disposition des graines de tournesol. Bernoulli la trouvait si belle qu'il la fit graver sur sa tombe !
\end{savaistu}

\courssub{Des éléments caractéristiques à l'écriture complexe, et réciproquement}

\begin{methode}[Identifier une similitude donnée par son écriture complexe]
Pour déterminer la nature et les éléments caractéristiques de la transformation $z' = az + b$ :
\begin{enumerate}
\item \textbf{Tester $a$.} Si $a = 1$, c'est la translation de vecteur d'affixe $b$ (ou l'identité si $b = 0$). Stop.
\item \textbf{Sinon, calculer le centre.} $\omega = \dfrac{b}{1 - a}$.
\item \textbf{Calculer le rapport.} $k = |a|$.
\item \textbf{Calculer l'angle.} $\theta = \arg(a) \pmod{2\pi}$, en écrivant d'abord $a$ sous forme trigonométrique $k e^{i\theta}$ avec $k > 0$.
\item \textbf{Conclure} en énonçant : « $s$ est la similitude directe de centre $\Omega(\omega)$, de rapport $k$ et d'angle $\theta$ ». Cas particuliers à reconnaître : si $\theta = 0 \pmod{2\pi}$, c'est une homothétie de rapport $k$ ; si $k = 1$, c'est une rotation d'angle $\theta$.
\end{enumerate}
\end{methode}

\begin{methode}[Écrire $z' = az + b$ connaissant le centre, le rapport et l'angle]
Pour déterminer l'écriture complexe de la similitude de centre $\Omega(\omega)$, de rapport $k$ et d'angle $\theta$ :
\begin{enumerate}
\item Poser $a = k\,e^{i\theta}$ (calculer explicitement la valeur algébrique de $a$).
\item Calculer $b = \omega(1 - a)$, qui traduit que $\Omega$ est le point fixe.
\item Conclure : $z' = az + b$. On peut vérifier aussitôt que $\omega = a\omega + b$.
\end{enumerate}
\end{methode}

\begin{exoresolu}[Deux sens du même exercice]
\textbf{Partie 1.} Déterminer la nature et les éléments caractéristiques de la transformation $s$ d'écriture complexe $z' = i\sqrt{3}\,z + 2 - 2i\sqrt{3}$.

\textbf{Partie 2.} Déterminer l'écriture complexe de la similitude $t$ de centre $\Omega$ d'affixe $1 + i$, de rapport $2$ et d'angle $-\dfrac{\pi}{3}$.
\tcblower
\textbf{Partie 1.} On applique la méthode.
\begin{enumerate}
\item $a = i\sqrt{3} \neq 1$ : ce n'est pas une translation.
\item Centre : $\omega = \dfrac{b}{1-a} = \dfrac{2 - 2i\sqrt{3}}{1 - i\sqrt{3}}$. On multiplie numérateur et dénominateur par le conjugué du dénominateur :
\[ \omega = \frac{(2 - 2i\sqrt{3})(1 + i\sqrt{3})}{1 + 3} = \frac{2 + 2i\sqrt{3} - 2i\sqrt{3} - 2i^2 \times 3}{4} = \frac{2 + 6}{4} = 2. \]
\item Rapport : $k = |i\sqrt{3}| = \sqrt{3}$.
\item Angle : $a = i\sqrt{3} = \sqrt{3}\,e^{i\pi/2}$, donc $\theta = \dfrac{\pi}{2} \pmod{2\pi}$.
\end{enumerate}
\textbf{Conclusion :} $s$ est la similitude directe de centre $\Omega$ d'affixe $2$, de rapport $\sqrt{3}$ et d'angle $\dfrac{\pi}{2}$.

\textbf{Partie 2.}
\begin{enumerate}
\item $a = 2e^{-i\pi/3} = 2\left(\dfrac{1}{2} - i\dfrac{\sqrt{3}}{2}\right) = 1 - i\sqrt{3}$.
\item $b = \omega(1 - a) = (1 + i)\bigl(1 - (1 - i\sqrt{3})\bigr) = (1 + i)(i\sqrt{3}) = i\sqrt{3} + i^2\sqrt{3} = -\sqrt{3} + i\sqrt{3}$.
\end{enumerate}
\textbf{Conclusion :} $t$ a pour écriture complexe $z' = (1 - i\sqrt{3})\,z - \sqrt{3} + i\sqrt{3}$. Vérification : $a\omega + b = (1 - i\sqrt{3})(1 + i) - \sqrt{3} + i\sqrt{3} = 1 + i - i\sqrt{3} + \sqrt{3} - \sqrt{3} + i\sqrt{3} = 1 + i = \omega$. Le point $\Omega$ est bien fixe.
\end{exoresolu}

\courssub{Similitude définie par deux points et leurs images}

Une dernière situation, très fréquente au Baccalauréat : on ne donne ni $a$ ni les éléments caractéristiques, mais simplement deux points et leurs images. C'est naturel : une similitude directe est entièrement déterminée par l'image de deux points distincts, car l'écriture $z' = az + b$ contient exactement deux inconnues complexes $a$ et $b$, et chaque condition $s(A) = A'$ fournit une équation.

\begin{propriete}[Similitude définie par deux couples de points]
Soient $A$ et $B$ deux points distincts, $A'$ et $B'$ deux points du plan. Il existe une \textbf{unique} similitude directe $s$ telle que $s(A) = A'$ et $s(B) = B'$. Ses coefficients sont
\[ a = \frac{z_{A'} - z_{B'}}{z_A - z_B} \qquad \text{et} \qquad b = z_{A'} - a\,z_A. \]
\end{propriete}

\begin{experience}[Démonstration guidée]
\begin{enumerate}
\item \textbf{Écrire les conditions.} $s(A) = A'$ et $s(B) = B'$ se traduisent par
\[ \begin{cases} z_{A'} = a\,z_A + b \\ z_{B'} = a\,z_B + b \end{cases} \]
\item \textbf{Éliminer $b$.} En soustrayant membre à membre : $z_{A'} - z_{B'} = a(z_A - z_B)$. Comme $A \neq B$, on a $z_A - z_B \neq 0$, d'où l'unique valeur $a = \dfrac{z_{A'} - z_{B'}}{z_A - z_B}$.
\item \textbf{En déduire $b$.} La première équation donne $b = z_{A'} - a\,z_A$.
\end{enumerate}
L'unicité de $a$ et $b$ garantit l'unicité de la similitude.
\end{experience}

\begin{exoresolu}[Similitude envoyant deux points sur deux points]
Le plan est muni d'un repère orthonormé direct. Soient $A(1)$, $B(2 + i)$, $A'(2i)$ et $B'(-2 + 3i)$. Déterminer l'écriture complexe de la similitude directe $s$ telle que $s(A) = A'$ et $s(B) = B'$, puis ses éléments caractéristiques.
\tcblower
\textbf{Étape 1 : calcul de $a$.}
\[ a = \frac{z_{A'} - z_{B'}}{z_A - z_B} = \frac{2i - (-2 + 3i)}{1 - (2 + i)} = \frac{2 - i}{-1 - i} = \frac{(2 - i)(-1 + i)}{(-1)^2 + 1^2} = \frac{-2 + 2i + i - i^2}{2} = \frac{-1 + 3i}{2}. \]
\textbf{Étape 2 : calcul de $b$.}
\[ b = z_{A'} - a\,z_A = 2i - \frac{-1 + 3i}{2} \times 1 = \frac{4i + 1 - 3i}{2} = \frac{1 + i}{2}. \]
Donc $s$ a pour écriture $z' = \dfrac{-1 + 3i}{2}\,z + \dfrac{1 + i}{2}$.

\textbf{Étape 3 : éléments caractéristiques.} Comme $a \neq 1$ :
\[ \omega = \frac{b}{1 - a} = \frac{\dfrac{1+i}{2}}{1 - \dfrac{-1+3i}{2}} = \frac{\dfrac{1+i}{2}}{\dfrac{3 - 3i}{2}} = \frac{1 + i}{3 - 3i} = \frac{(1+i)(1+i)}{3(1-i)(1+i)} = \frac{2i}{6} = \frac{i}{3}. \]
\[ k = |a| = \frac{|-1 + 3i|}{2} = \frac{\sqrt{1 + 9}}{2} = \frac{\sqrt{10}}{2}, \qquad \theta = \arg(-1 + 3i) = \pi - \arctan(3) \pmod{2\pi}. \]
\textbf{Conclusion :} $s$ est la similitude de centre $\Omega\left(\dfrac{i}{3}\right)$, de rapport $\dfrac{\sqrt{10}}{2}$ et d'angle $\pi - \arctan(3)$.
\end{exoresolu}

\begin{attention}[Erreurs fréquentes à bannir]
\begin{itemize}
\item \textbf{Oublier de tester $a = 1$.} Si $a = 1$, la formule $\omega = \dfrac{b}{1-a}$ fait diviser par zéro : la transformation est une translation, il n'y a pas de centre.
\item \textbf{Confondre $a$ et le rapport.} Le rapport est $k = |a|$, jamais $a$ lui-même. Si $a = -2$, le rapport est $2$ et l'angle est $\pi$ (voir l'attention précédente).
\item \textbf{Se tromper de formule pour $b$.} C'est $b = \omega(1 - a)$ et non $b = \omega(a - 1)$. Pour ne plus hésiter, retenez la phrase : « le centre est fixe », c'est-à-dire $\omega = a\omega + b$, d'où $b = \omega - a\omega = \omega(1 - a)$.
\item \textbf{Oublier le centre dans la forme réduite.} La relation $z' - \omega = a(z - \omega)$ porte sur les affixes \emph{relatives au centre}. Écrire $z' = az$ n'est correct que si le centre est l'origine.
\end{itemize}
\end{attention}

\begin{exercice}[À toi de jouer]
Dans chaque cas, déterminer la nature et les éléments caractéristiques de la transformation d'écriture complexe :
\begin{enumerate}
\item $z' = 3z - 2 + 4i$ ;
\item $z' = \left(\dfrac{1}{2} + i\dfrac{\sqrt{3}}{2}\right)z + 1$ ;
\item $z' = -2iz + 3 - i$.
\end{enumerate}
\end{exercice}

\begin{corrige}[Exercice : identification de similitudes]
\textbf{1.} $a = 3 \neq 1$. Centre : $\omega = \dfrac{-2 + 4i}{1 - 3} = \dfrac{-2 + 4i}{-2} = 1 - 2i$. Rapport $k = 3$, angle $\theta = 0$ : c'est l'\textbf{homothétie} de centre $\Omega(1 - 2i)$ et de rapport $3$.

\textbf{2.} $a = \dfrac{1}{2} + i\dfrac{\sqrt{3}}{2} = e^{i\pi/3}$, donc $|a| = 1$ : c'est une \textbf{rotation}. Centre : $\omega = \dfrac{1}{1 - e^{i\pi/3}} = \dfrac{1}{\dfrac{1}{2} - i\dfrac{\sqrt{3}}{2}} = \dfrac{1}{e^{-i\pi/3}} = e^{i\pi/3} = \dfrac{1}{2} + i\dfrac{\sqrt{3}}{2}$. Rotation de centre $\Omega\left(e^{i\pi/3}\right)$ et d'angle $\dfrac{\pi}{3}$.

\textbf{3.} $a = -2i = 2e^{-i\pi/2}$, donc $k = 2$ et $\theta = -\dfrac{\pi}{2}$. Centre : $\omega = \dfrac{3 - i}{1 + 2i} = \dfrac{(3 - i)(1 - 2i)}{5} = \dfrac{3 - 6i - i + 2i^2}{5} = \dfrac{1 - 7i}{5}$. Similitude de centre $\Omega\left(\dfrac{1 - 7i}{5}\right)$, de rapport $2$ et d'angle $-\dfrac{\pi}{2}$.
\end{corrige}

\begin{aretenir}[L'essentiel de la section]
\begin{itemize}
\item Toute similitude directe $z' = az + b$ avec $a \neq 1$ admet un unique point fixe $\Omega$ d'affixe $\omega = \dfrac{b}{1 - a}$, et s'écrit sous forme réduite $z' - \omega = a(z - \omega)$.
\item Rapport : $k = |a|$ ; angle : $\theta = \arg(a) \pmod{2\pi}$. Toujours écrire $a = k e^{i\theta}$ avec $k > 0$.
\item $s = h(\Omega, k) \circ r(\Omega, \theta) = r(\Omega, \theta) \circ h(\Omega, k)$ : homothétie et rotation de même centre commutent.
\item Connaissant $\Omega(\omega)$, $k$, $\theta$ : $a = k e^{i\theta}$ puis $b = \omega(1 - a)$.
\item Connaissant $s(A) = A'$ et $s(B) = B'$ : $a = \dfrac{z_{A'} - z_{B'}}{z_A - z_B}$ puis $b = z_{A'} - a z_A$.
\end{itemize}
\end{aretenir}

\coursec{Propriétés de conservation et images des configurations}

Dans la section précédente, tu as appris à reconnaître une similitude directe $s$ d'écriture complexe $z' = az + b$ ($a \neq 0$) et à en extraire ses \cle{éléments caractéristiques} : le \cle{rapport} $k = |a|$, l'\cle{angle} $\theta = \operatorname{arg}(a)$ et le \cle{centre} $\Omega$ d'affixe $\omega = \dfrac{b}{1-a}$ (lorsque $a \neq 1$). Tu sais donc maintenant \emph{qui} est $s$. La question naturelle est : \emph{que fait-elle aux figures} ? Quand on applique une similitude à un cercle, à un triangle, à un parallélogramme, que devient la figure ?

L'intuition est simple : une similitude directe, c'est un « agrandissement-rotation ». Pense à la photocopieuse du cybercafé de quartier qui agrandit un document à $150\,\%$ : les formes sont préservées, seules les tailles changent. Cette section transforme cette intuition en théorèmes précis, tous conséquences d'un seul fait fondamental : la \cle{conservation du barycentre}.

\courssub{Le fait fondamental : conservation du barycentre}

\begin{experience}[Découverte guidée : pourquoi le barycentre est-il conservé ?]
Soit $s$ la similitude directe d'écriture complexe $z' = az + b$. Soient $A$ et $B$ deux points d'affixes $z_A$ et $z_B$, d'images $A'$ et $B'$ d'affixes $z_{A'} = a z_A + b$ et $z_{B'} = a z_B + b$.
\begin{enumerate}
\item Soit $G$ le milieu de $[AB]$. Quelle est l'affixe de $G$ ?
\item Calcule l'affixe de $G' = s(G)$ en fonction de $z_{A'}$ et $z_{B'}$.
\item Que peux-tu en conclure sur $G'$ ?
\end{enumerate}
\textbf{Réponses.} (1) $z_G = \dfrac{z_A + z_B}{2}$. (2) On calcule :
\[ z_{G'} = a z_G + b = a\,\frac{z_A + z_B}{2} + b = \frac{a z_A + a z_B + 2b}{2} = \frac{(a z_A + b) + (a z_B + b)}{2} = \frac{z_{A'} + z_{B'}}{2}. \]
(3) $G'$ est le milieu de $[A'B']$ : \textbf{l'image du milieu est le milieu des images}. Le calcul n'a utilisé que la distributivité de la multiplication par $a$ et l'addition de $b$ : rien d'autre !
\end{experience}

Ce petit calcul se généralise à n'importe quel barycentre, et c'est la clé de toute la section.

\begin{propriete}[Conservation du barycentre]
Soit $s$ une similitude directe du plan. Pour tous points $A_1, A_2, \ldots, A_n$ et tous réels $\alpha_1, \alpha_2, \ldots, \alpha_n$ de somme non nulle, si $G$ est le barycentre des points pondérés $(A_i, \alpha_i)$, alors $s(G)$ est le barycentre des points pondérés $(s(A_i), \alpha_i)$, avec les \emph{mêmes coefficients}.
\end{propriete}

\begin{experience}[Démonstration guidée (exigible)]
Notons $z_i$ l'affixe de $A_i$, $z_i' = a z_i + b$ l'affixe de $A_i' = s(A_i)$, et $S = \alpha_1 + \alpha_2 + \cdots + \alpha_n \neq 0$.
\begin{enumerate}
\item Écris l'affixe $g$ de $G$ : $g = \dfrac{1}{S}\displaystyle\sum_{i=1}^{n} \alpha_i z_i$.
\item Calcule l'affixe $g'$ de $G' = s(G)$ :
\begin{align*}
g' &= a g + b = \frac{a}{S}\sum_{i=1}^{n} \alpha_i z_i + b = \frac{1}{S}\left(\sum_{i=1}^{n} \alpha_i\, a z_i + S\, b\right) \\
&= \frac{1}{S}\left(\sum_{i=1}^{n} \alpha_i\, a z_i + \sum_{i=1}^{n} \alpha_i b\right) = \frac{1}{S}\sum_{i=1}^{n} \alpha_i (a z_i + b) = \frac{1}{S}\sum_{i=1}^{n} \alpha_i z_i'.
\end{align*}
\item Conclus : $g'$ est l'affixe du barycentre des $(A_i', \alpha_i)$. \hfill $\blacksquare$
\end{enumerate}
L'astuce de l'étape 2 consiste à écrire $b = \dfrac{S\, b}{S} = \dfrac{1}{S}\sum \alpha_i b$ pour faire « rentrer » le $b$ dans la somme.
\end{experience}

\begin{propriete}[Conséquences immédiates]
Toute similitude directe :
\begin{itemize}
\item \cle{conserve les milieux} : l'image du milieu de $[AB]$ est le milieu de $[A'B']$ (barycentre de coefficients $(1,1)$) ;
\item \cle{conserve l'alignement} : si $A$, $B$, $C$ sont alignés, alors $A'$, $B'$, $C'$ sont alignés (un point d'une droite $(AB)$ est un barycentre de $A$ et $B$) ;
\item \cle{conserve le parallélisme} : les images de deux droites parallèles sont deux droites parallèles ;
\item \cle{conserve les angles orientés} : pour tous points $A$, $B$, $C$ avec $A \neq B$ et $A \neq C$,
\[ \left(\overrightarrow{A'B'}\,;\,\overrightarrow{A'C'}\right) = \left(\overrightarrow{AB}\,;\,\overrightarrow{AC}\right) \quad [2\pi]. \]
\end{itemize}
\end{propriete}

\textbf{Justification de la conservation des angles.} Avec $z' = az + b$, on a $z_{C'} - z_{A'} = a(z_C - z_A)$ et $z_{B'} - z_{A'} = a(z_B - z_A)$, donc
\[ \frac{z_{C'} - z_{A'}}{z_{B'} - z_{A'}} = \frac{z_C - z_A}{z_B - z_A}, \]
et en prenant les arguments : $(\overrightarrow{A'B'};\overrightarrow{A'C'}) = (\overrightarrow{AB};\overrightarrow{AC})$ $[2\pi]$. On a au passage l'égalité des modules : $A'C' = |a|\cdot AC = k \cdot AC$ : \textbf{toutes les distances sont multipliées par le rapport $k$}.

\begin{attention}[Directe ou indirecte : le mot qui change tout]
Seules les similitudes \emph{directes} conservent les angles \textbf{orientés}. Une symétrie axiale, par exemple, conserve les angles géométriques mais \emph{renverse} les angles orientés. Au Bac, quand tu invoques la conservation des angles orientés, vérifie que la transformation est bien de la forme $z' = az + b$ (et non $z' = a\bar{z} + b$).
\end{attention}

\courssub{Image d'une droite}

\begin{propriete}[Image d'une droite]
L'image d'une droite $\mathcal{D}$ par une similitude directe $s$ est une droite $\mathcal{D}'$. De plus, si $\vec{u}$ dirige $\mathcal{D}$, alors $\vec{u}'$, image de $\vec{u}$ par la partie linéaire de $s$ (rotation d'angle $\theta$ composée avec l'homothétie de rapport $k$), dirige $\mathcal{D}'$.
\end{propriete}

\textbf{Démonstration.} Soient $A$ et $B$ deux points distincts de $\mathcal{D}$. Tout point $M$ de $\mathcal{D}$ est un barycentre de $A$ et $B$ (par exemple $M = \text{bar}\{(A, 1-t), (B, t)\}$ pour un réel $t$). Par conservation du barycentre, $M' = s(M)$ est un barycentre de $A' = s(A)$ et $B' = s(B)$, donc $M'$ appartient à la droite $(A'B')$. Réciproquement, en appliquant le même raisonnement à la similitude réciproque $s^{-1}$ (qui existe car $a \neq 0$), tout point de $(A'B')$ est l'image d'un point de $\mathcal{D}$. Donc $s(\mathcal{D}) = (A'B')$. \hfill $\blacksquare$

\begin{methode}[Déterminer l'image d'une droite]
Pour déterminer l'équation de la droite $\mathcal{D}'$, image d'une droite $\mathcal{D}$ par une similitude $s$ :
\begin{enumerate}
\item choisir \textbf{deux points} $A$ et $B$ de $\mathcal{D}$ (prends les points aux coordonnées les plus simples, par exemple les intersections avec les axes) ;
\item calculer leurs images $A' = s(A)$ et $B' = s(B)$ grâce à l'écriture complexe ou analytique ;
\item écrire une équation de la droite $(A'B')$ (vecteur directeur $\overrightarrow{A'B'}$, puis équation cartésienne).
\end{enumerate}
Inutile de transformer l'équation de $\mathcal{D}$ point par point : deux images suffisent !
\end{methode}

\begin{exemplebox}[Image d'une droite]
Soit $s$ la similitude d'écriture complexe $z' = (1+i)z - 1$ et $\mathcal{D}$ l'axe des réels (équation $y = 0$).
\begin{enumerate}
\item Prenons $O(0)$ et $I(1)$ sur $\mathcal{D}$. Leurs images : $z_{O'} = -1$ et $z_{I'} = (1+i) - 1 = i$.
\item $\mathcal{D}' = (O'I')$ passe par $O'(-1\,;\,0)$ et $I'(0\,;\,1)$. Un vecteur directeur est $\overrightarrow{O'I'}\begin{pmatrix} 1 \\ 1 \end{pmatrix}$.
\item Équation : $M(x\,;\,y) \in \mathcal{D}' \iff \det(\overrightarrow{O'M}, \overrightarrow{O'I'}) = 0 \iff (x+1) - y = 0$, soit $\mathcal{D}' : y = x + 1$.
\end{enumerate}
Remarque : la direction de $\mathcal{D}$ (angle $0$) est tournée de $\frac{\pi}{4}$ (angle de $s$), ce qui donne une droite de direction $\frac{\pi}{4}$ (vecteur $(1,1)$), cohérent avec le résultat.
\end{exemplebox}

\courssub{Image d'un cercle}

\begin{propriete}[Image d'un cercle]
Soit $s$ une similitude directe de rapport $k$. L'image du cercle $\mathcal{C}$ de centre $\Omega$ et de rayon $r$ est le cercle $\mathcal{C}'$ de centre $\Omega' = s(\Omega)$ et de rayon $k\,r$.
\end{propriete}

\textbf{Démonstration (exigible).} Notons $\omega$ l'affixe de $\Omega$ et $\omega' = a\omega + b$ celle de $\Omega'$. Pour tout point $M$ d'affixe $z$, d'image $M'$ d'affixe $z' = az + b$ :
\[ \Omega'M' = |z' - \omega'| = |az + b - (a\omega + b)| = |a(z - \omega)| = |a|\cdot|z - \omega| = k \cdot \Omega M. \]
Ainsi : $M \in \mathcal{C} \iff \Omega M = r \iff \Omega'M' = kr \iff M' \in \mathcal{C}(\Omega', kr)$. \hfill $\blacksquare$

\begin{attention}[Le centre de l'image est l'image du centre]
Le centre de $\mathcal{C}'$ est $s(\Omega)$, et \emph{non pas} le centre $\Omega_s$ de la similitude ! Ce sont deux notions différentes qui portent malheureusement le même nom. Le centre de la similitude est le point fixe ; le centre du cercle image est l'image du centre du cercle de départ.
\end{attention}

\begin{exoresolu}[Image d'un cercle]
\textbf{Énoncé.} On considère la similitude directe $s$ d'écriture complexe $z' = (1+i)z - 1$ et le cercle $\mathcal{C}$ de centre $A$ d'affixe $1 - i$ et de rayon $2$. Déterminer l'image de $\mathcal{C}$ par $s$.
\tcblower
\textbf{Solution.} \textbf{Étape 1 : rapport de la similitude.} $a = 1+i$, donc $k = |a| = \sqrt{1^2 + 1^2} = \sqrt{2}$.

\textbf{Étape 2 : image du centre.} $z_{A'} = (1+i)(1-i) - 1 = (1 - i^2) - 1 = 2 - 1 = 1$. Donc $A'(1\,;\,0)$.

\textbf{Étape 3 : rayon image.} $r' = k \cdot r = 2\sqrt{2}$.

\textbf{Conclusion.} $s(\mathcal{C})$ est le cercle de centre $A'(1\,;\,0)$ et de rayon $2\sqrt{2}$. Son équation cartésienne est $(x-1)^2 + y^2 = 8$.
\end{exoresolu}

\begin{popfigure}
\begin{center}
\begin{tikzpicture}[scale=0.85]
\% Cercle initial
\coordinate (Om) at (0,0);
\coordinate (Omp) at (5.2,1.2);
\draw[popline, dashed] (Om) -- (Omp);
\draw[popblue, thick] (Om) circle (1.4);
\draw[poporange, thick] (Omp) circle (2.4);
\% Rayons
\draw[popblue, thick, -{Stealth}] (Om) -- (60:1.4) node[midway, above left, popblue] {$r$};
\draw[poporange, thick, -{Stealth}] (Omp) -- ++(60:2.4) node[midway, above left, poporange] {$kr$};
\% Points
\fill[popblue] (Om) circle (1.8pt) node[below left] {$\Omega$};
\fill[poporange] (Omp) circle (1.8pt) node[below right] {$\Omega' = s(\Omega)$};
\fill[popblue] (60:1.4) circle (1.5pt) node[right, popblue] {$M$};
\fill[poporange] ($(Omp)+(60:2.4)$) circle (1.5pt) node[right, poporange] {$M' = s(M)$};
\% Flèche similitude
\draw[popdark, -{Stealth}, thick] (2.1,2.6) to[bend left=25] node[above] {$s$ : rapport $k$} (4.6,3.4);
\node[popblue] at (-1.2,-1.3) {$\mathcal{C}(\Omega, r)$};
\node[poporange] at (7.6,-0.6) {$\mathcal{C}'(\Omega', kr)$};
\end{tikzpicture}
\end{center}
\legende{L'image d'un cercle par une similitude directe de rapport $k$ est un cercle : le centre est l'image du centre, le rayon est multiplié par $k$.}
\end{popfigure}

\courssub{Image d'un triangle : triangles directement semblables}

\begin{propriete}[Image d'un triangle]
Soit $s$ une similitude directe de rapport $k$, et $ABC$ un triangle. Son image $A'B'C'$ est un triangle tel que :
\begin{itemize}
\item les distances sont multipliées par $k$ : $A'B' = k\,AB$, $B'C' = k\,BC$, $A'C' = k\,AC$ ;
\item les angles orientés sont conservés : $(\overrightarrow{A'B'};\overrightarrow{A'C'}) = (\overrightarrow{AB};\overrightarrow{AC})$ $[2\pi]$, et de même pour les autres sommets.
\end{itemize}
On dit que $A'B'C'$ est \cle{directement semblable} à $ABC$ : mêmes angles, mêmes orientations, côtés proportionnels dans le rapport $k$.
\end{propriete}

\textbf{Démonstration.} La proportionnalité des côtés vient de $\Omega'M' = k\,\Omega M$ appliquée à chaque paire de sommets ; la conservation des angles orientés a été démontrée plus haut via les rapports d'affixes. \hfill $\blacksquare$

\begin{popfigure}
\begin{center}
\begin{tikzpicture}[scale=0.9]
\% Triangle initial
\coordinate (A) at (0,0);
\coordinate (B) at (2.6,0);
\coordinate (C) at (0.9,1.9);
\draw[popblue, thick] (A) -- (B) -- (C) -- cycle;
\fill[popblue] (A) circle (1.6pt) node[below left] {$A$};
\fill[popblue] (B) circle (1.6pt) node[below right] {$B$};
\fill[popblue] (C) circle (1.6pt) node[above] {$C$};
\% Angle en A
\draw[popblue] (0.55,0) arc (0:64:0.55) node[midway, right, popblue] {$\alpha$};
\% Triangle image (homothétie-rotation)
\coordinate (Ap) at (5.6,0.6);
\coordinate (Bp) at ($(Ap) + (25:3.4)$);
\coordinate (Cp) at ($(Ap) + (89:2.5)$);
\draw[poporange, thick] (Ap) -- (Bp) -- (Cp) -- cycle;
\fill[poporange] (Ap) circle (1.6pt) node[below left] {$A'$};
\fill[poporange] (Bp) circle (1.6pt) node[right] {$B'$};
\fill[poporange] (Cp) circle (1.6pt) node[above] {$C'$};
\% Angle en A'
\draw[poporange] ($(Ap)+(25:0.55)$) arc (25:89:0.55) node[midway, right, poporange] {$\alpha$};
\% Flèche
\draw[popdark, -{Stealth}, thick] (2.9,2.3) to[bend left=20] node[above] {$s$} (5.2,2.6);
\node[popmuted] at (4.3,-0.5) {$A'B' = k\,AB$, angles conservés};
\end{tikzpicture}
\end{center}
\legende{L'image d'un triangle par une similitude directe est un triangle directement semblable : mêmes angles (ici $\widehat{A} = \widehat{A'} = \alpha$), côtés multipliés par $k$, même sens de parcours.}
\end{popfigure}

\begin{aretenir}[Conséquences pour les triangles particuliers]
Puisque les angles et les rapports de longueurs sont conservés, une similitude directe transforme :
\begin{itemize}
\item un triangle \textbf{rectangle} en un triangle rectangle (l'angle droit est conservé) ;
\item un triangle \textbf{isocèle} en un triangle isocèle (l'égalité de deux longueurs est conservée) ;
\item un triangle \textbf{équilatéral} en un triangle équilatéral ;
\item un triangle rectangle isocèle en un triangle rectangle isocèle.
\end{itemize}
\end{aretenir}

\courssub{Image d'un parallélogramme}

\begin{propriete}[Image d'un parallélogramme]
L'image d'un parallélogramme $ABCD$ par une similitude directe est un parallélogramme $A'B'C'D'$. En particulier, l'image d'un carré est un carré, et l'image d'un rectangle est un rectangle.
\end{propriete}

\textbf{Démonstration.} $ABCD$ est un parallélogramme si et seulement si ses diagonales $[AC]$ et $[BD]$ ont le même milieu $I$. Or $s$ conserve les milieux : $I' = s(I)$ est le milieu commun de $[A'C']$ et $[B'D']$. Un quadrilatère dont les diagonales se coupent en leur milieu est un parallélogramme, donc $A'B'C'D'$ est un parallélogramme. Pour le carré : on ajoute la conservation des angles droits (rectangle) et l'égalité des côtés (losange), donc l'image d'un carré est un carré. \hfill $\blacksquare$

\begin{popfigure}
\begin{center}
\begin{tikzpicture}[scale=0.8]
\% Parallélogramme initial
\coordinate (A) at (0,0);
\coordinate (B) at (2.8,0);
\coordinate (D) at (1,1.7);
\coordinate (C) at ($(B)+(D)-(A)$);
\draw[popblue, thick] (A) -- (B) -- (C) -- (D) -- cycle;
\draw[popblue, dashed] (A) -- (C);
\draw[popblue, dashed] (B) -- (D);
\coordinate (I) at ($(A)!0.5!(C)$);
\fill[popblue] (A) circle (1.5pt) node[below left] {$A$};
\fill[popblue] (B) circle (1.5pt) node[below right] {$B$};
\fill[popblue] (C) circle (1.5pt) node[above right] {$C$};
\fill[popblue] (D) circle (1.5pt) node[above left] {$D$};
\fill[popblue] (I) circle (1.8pt) node[below, popblue] {$I$};
\% Image
\coordinate (Ap) at (6.2,0.4);
\coordinate (Bp) at ($(Ap)+(15:3.6)$);
\coordinate (Dp) at ($(Ap)+(75:2.2)$);
\coordinate (Cp) at ($(Bp)+(Dp)-(Ap)$);
\draw[poporange, thick] (Ap) -- (Bp) -- (Cp) -- (Dp) -- cycle;
\draw[poporange, dashed] (Ap) -- (Cp);
\draw[poporange, dashed] (Bp) -- (Dp);
\coordinate (Ip) at ($(Ap)!0.5!(Cp)$);
\fill[poporange] (Ap) circle (1.5pt) node[below left] {$A'$};
\fill[poporange] (Bp) circle (1.5pt) node[below right] {$B'$};
\fill[poporange] (Cp) circle (1.5pt) node[above right] {$C'$};
\fill[poporange] (Dp) circle (1.5pt) node[above left] {$D'$};
\fill[poporange] (Ip) circle (1.8pt) node[right, poporange] {$I' = s(I)$};
\draw[popdark, -{Stealth}, thick] (3.6,2.4) to[bend left=20] node[above] {$s$} (5.9,2.7);
\end{tikzpicture}
\end{center}
\legende{Le milieu $I$ des diagonales est envoyé sur le milieu $I'$ des diagonales images : l'image d'un parallélogramme est un parallélogramme.}
\end{popfigure}

\begin{exoresolu}[Nature d'un quadrilatère image]
\textbf{Énoncé.} Soit $s$ la similitude d'écriture complexe $z' = 2i\,z + 1$. On donne $A(0)$, $B(1)$, $C(1+i)$, $D(i)$. Montrer que $ABCD$ est un carré, puis déterminer la nature de $A'B'C'D'$.
\tcblower
\textbf{Solution.} $AB = |1| = 1$, $AD = |i| = 1$, $\overrightarrow{AB} \cdot \overrightarrow{AD} = 0$ (car $(\overrightarrow{AB};\overrightarrow{AD}) = \operatorname{arg}\!\left(\dfrac{i}{1}\right) = \dfrac{\pi}{2}$), et $\overrightarrow{AB} + \overrightarrow{AD} = 1 + i = \overrightarrow{AC}$ : $ABCD$ est un carré.

Ses images : $z_{A'} = 1$, $z_{B'} = 2i + 1$, $z_{C'} = 2i(1+i) + 1 = -1 + 2i$, $z_{D'} = 2i \cdot i + 1 = -1$.

$s$ est une similitude directe ($a = 2i \neq 0$) de rapport $k = |2i| = 2$ et d'angle $\dfrac{\pi}{2}$. Par conservation des angles et proportionnalité des longueurs, $A'B'C'D'$ est un \textbf{carré} de côté $k \cdot AB = 2$. Vérification : $A'B' = |2i| = 2$ et $A'D' = |-2| = 2$, avec $(\overrightarrow{A'B'};\overrightarrow{A'D'}) = \operatorname{arg}\!\left(\dfrac{-2}{2i}\right) = \operatorname{arg}(i) = \dfrac{\pi}{2}$. \hfill $\blacksquare$
\end{exoresolu}

\begin{attention}[Pièges fréquents au Bac]
\begin{itemize}
\item Dire « l'image d'un cercle est un cercle de même centre » : \textbf{faux} en général ! Le centre bouge, sauf si $\Omega$ est le point fixe de $s$.
\item Oublier de multiplier le rayon par $k$ : le rayon de l'image est $kr$, pas $r$.
\item Affirmer qu'une similitude conserve les distances : seules les \emph{isométries} ($k = 1$) les conservent ; une similitude conserve les \textbf{rapports} de distances.
\item Conclure « parallélogramme » sans justifier : la démonstration passe par la conservation du milieu des diagonales, cite-la.
\end{itemize}
\end{attention}

\begin{savaistu}[La conservation des angles : l'arme fatale du Bac]
Au Baccalauréat, la conservation des angles orientés permet de démontrer qu'un triangle est rectangle ou isocèle \emph{sans aucun calcul de distance}. Stratégie type : on te donne une similitude $s$ envoyant $A$ sur $A'$ et $B$ sur $B'$, et on te demande la nature du triangle $A'B'C'$. Plutôt que de calculer trois longueurs, montre que $(\overrightarrow{A'B'};\overrightarrow{A'C'}) = (\overrightarrow{AB};\overrightarrow{AC})$ et lis l'angle sur la figure initiale ! Cette idée — remplacer le calcul par la structure — est au cœur des mathématiques modernes : Felix Klein l'a théorisée dès 1872 dans son \textbf{programme d'Erlangen}, qui définit chaque géométrie par les propriétés que ses transformations conservent. Les similitudes, elles, conservent les \textbf{formes} : c'est pourquoi une carte routière du Cameroun, à n'importe quelle échelle, respecte les angles des routes et la forme du triangle Yaoundé--Douala--Bafoussam.
\end{savaistu}

\begin{exercice}[Pour t'entraîner]
Soit $s$ la similitude d'écriture complexe $z' = (1 - i\sqrt{3})\,z + 2i$.
\begin{enumerate}
\item Détermine le rapport et l'angle de $s$.
\item Détermine l'image du cercle de centre $\Omega(1 + i)$ et de rayon $3$.
\item $ABC$ est un triangle équilatéral direct. Quelle est la nature de son image $A'B'C'$ ? Justifie.
\end{enumerate}
\end{exercice}

\begin{corrige}[Exercice 1]
\textbf{1.} $a = 1 - i\sqrt{3}$, donc $k = |a| = \sqrt{1 + 3} = 2$ et $\theta = \operatorname{arg}(a) = -\dfrac{\pi}{3}$ $[2\pi]$.

\textbf{2.} Image du centre : $z_{\Omega'} = (1 - i\sqrt{3})(1+i) + 2i = 1 + i - i\sqrt{3} + \sqrt{3} + 2i = (1+\sqrt{3}) + i(3 - \sqrt{3})$. Le rayon image est $kr = 2 \times 3 = 6$. L'image est le cercle de centre $\Omega'\big(1+\sqrt{3}\,;\,3-\sqrt{3}\big)$ et de rayon $6$.

\textbf{3.} $s$ conserve les angles orientés et multiplie les distances par $k = 2$. Les trois côtés de $ABC$ étant égaux, ceux de $A'B'C'$ le sont aussi ($A'B' = B'C' = A'C' = 2\,AB$), et le sens direct est conservé : $A'B'C'$ est un \textbf{triangle équilatéral direct}.
\end{corrige}

\coursec{Détermination d'une similitude et applications aux configurations}

Dans la section précédente, nous avons vu que les similitudes directes conservent les angles orientés, les milieux, l'alignement, l'orthogonalité, et transforment droites en droites et cercles en cercles. Autrement dit : une similitude directe \emph{préserve la forme} des figures, elle n'en change que la taille, la position et l'orientation. C'est précisément cette fidélité qui en fait un outil redoutable pour résoudre des problèmes de configuration. Mais avant de pouvoir l'utiliser, il faut savoir la \cle{fabriquer} : c'est l'objet de cette dernière section. Nous allons apprendre à déterminer une similitude à partir de données variées (ses éléments caractéristiques, ou les images de deux points), à reconnaître sa nature à partir de son écriture complexe, puis à l'utiliser comme instrument de démonstration géométrique.

\courssub{Problème type 1 : similitude donnée par son centre, son rapport et son angle}

Commençons par le problème le plus direct. On te donne le centre $\Omega$ d'affixe $\omega$, le rapport $k > 0$ et l'angle $\theta$ d'une similitude directe $s$, et on te demande son écriture complexe $z' = az + b$.

L'intuition est simple : dire que $M'$ est l'image de $M$ par la similitude de centre $\Omega$, de rapport $k$ et d'angle $\theta$, c'est dire que le vecteur $\overrightarrow{\Omega M'}$ s'obtient à partir de $\overrightarrow{\Omega M}$ en le tournant de $\theta$ puis en le dilatant de $k$. En langage complexe, tourner de $\theta$ c'est multiplier par $e^{i\theta}$, dilater de $k$ c'est multiplier par $k$. D'où :

\begin{propriete}[Écriture complexe à partir des éléments caractéristiques]
La similitude directe de centre $\Omega(\omega)$, de rapport $k > 0$ et d'angle $\theta$ a pour écriture complexe :
\[ z' - \omega = k\,e^{i\theta}(z - \omega), \]
c'est-à-dire $z' = az + b$ avec $a = k e^{i\theta}$ et $b = \omega - a\omega = \omega(1 - a)$.
\end{propriete}

\begin{methode}[Écrire $z' = az + b$ à partir de $(\Omega, k, \theta)$]
\begin{enumerate}
\item Calculer $a = k e^{i\theta}$ (mettre $e^{i\theta}$ sous forme algébrique si nécessaire).
\item Calculer $b = \omega(1 - a)$.
\item Écrire $z' = az + b$, puis \textbf{vérifier} que le point fixe est bien $\Omega$ : remplacer $z$ par $\omega$ et contrôler que $z' = \omega$.
\end{enumerate}
\end{methode}

\begin{exemplebox}[Similitude de centre $\Omega(1 + i)$, rapport $2$, angle $\dfrac{\pi}{3}$]
On a $a = 2e^{i\pi/3} = 2\left(\dfrac{1}{2} + i\dfrac{\sqrt{3}}{2}\right) = 1 + i\sqrt{3}$.

Puis $b = (1+i)(1 - a) = (1+i)(-i\sqrt{3}) = -i\sqrt{3} - i^2\sqrt{3} = \sqrt{3} - i\sqrt{3}$.

L'écriture complexe est donc :
\[ z' = (1 + i\sqrt{3})\,z + \sqrt{3} - i\sqrt{3}. \]
Vérification : pour $z = 1+i$, on obtient $(1+i\sqrt{3})(1+i) + \sqrt{3} - i\sqrt{3} = 1 + i + i\sqrt{3} - \sqrt{3} + \sqrt{3} - i\sqrt{3} = 1 + i$. Le point $\Omega$ est bien fixe.
\end{exemplebox}

\courssub{Problème type 2 : similitude déterminée par les images de deux points}

Voici maintenant le problème le plus fréquent au Baccalauréat. On connaît deux points distincts $A$ et $B$ et leurs images $A'$ et $B'$, et on cherche la similitude directe $s$ telle que $s(A) = A'$ et $s(B) = B'$. Existe-t-elle toujours ? Est-elle unique ?

\begin{propriete}[Détermination par deux points et leurs images — admis]
Étant donnés quatre points $A$, $B$, $A'$, $B'$ avec $A \neq B$, il existe une \textbf{unique} similitude directe $s$ telle que $s(A) = A'$ et $s(B) = B'$. Autrement dit, une similitude directe est entièrement déterminée par les images de deux points distincts.
\end{propriete}

Ce résultat est admis ; il est la conséquence directe de l'unicité du couple $(a, b)$ solution du système ci-dessous, système qui possède toujours une solution unique dès que $z_A \neq z_B$.

\begin{methode}[Déterminer la similitude telle que $s(A) = A'$ et $s(B) = B'$]
\begin{enumerate}
\item Écrire le système : $\begin{cases} z_{A'} = a\,z_A + b \\ z_{B'} = a\,z_B + b \end{cases}$
\item \textbf{Soustraire membre à membre} pour éliminer $b$ : $z_{A'} - z_{B'} = a(z_A - z_B)$, d'où
\[ a = \frac{z_{A'} - z_{B'}}{z_A - z_B}. \]
\item Reporter $a$ dans l'une des équations pour obtenir $b = z_{A'} - a\,z_A$.
\item \textbf{Vérifier} en recalculant $a z_B + b$ : on doit retrouver $z_{B'}$.
\item Si on demande les éléments caractéristiques : rapport $k = |a|$, angle $\theta = \arg(a)$, centre $\Omega$ d'affixe $\omega = \dfrac{b}{1 - a}$ (si $a \neq 1$).
\end{enumerate}
\end{methode}

\begin{exoresolu}[Détermination complète d'une similitude]
Le plan est muni d'un repère orthonormé direct. On considère les points $A$, $B$, $A'$, $B'$ d'affixes respectives $z_A = 1$, $z_B = i$, $z_{A'} = 2 + i$, $z_{B'} = 3i$. Déterminer l'écriture complexe de la similitude directe $s$ telle que $s(A) = A'$ et $s(B) = B'$, puis ses éléments caractéristiques.
\tcblower
\textbf{Étape 1 : le système.} On cherche $a$ et $b$ tels que
\[ \begin{cases} 2 + i = a \times 1 + b \\ 3i = a \times i + b \end{cases} \]
\textbf{Étape 2 : soustraction.} En retranchant la deuxième équation de la première :
\[ (2 + i) - 3i = a(1 - i) \iff 2 - 2i = a(1 - i) \iff a = \frac{2 - 2i}{1 - i} = \frac{2(1-i)}{1-i} = 2. \]
\textbf{Étape 3 : calcul de $b$.} De la première équation : $b = 2 + i - a = 2 + i - 2 = i$.

\textbf{Étape 4 : vérification.} $a z_B + b = 2i + i = 3i = z_{B'}$. Correct.

L'écriture complexe est donc $\boxed{z' = 2z + i}$.

\textbf{Étape 5 : éléments caractéristiques.} $a = 2$ est un réel strictement positif : c'est une \textbf{homothétie} de rapport $k = |a| = 2$ (angle $\theta = 0$). Le centre est le point fixe :
\[ \omega = \frac{b}{1 - a} = \frac{i}{1 - 2} = -i. \]
Conclusion : $s$ est l'homothétie de centre $\Omega(-i)$ et de rapport $2$.
\end{exoresolu}

\begin{attention}[Le piège de la soustraction]
L'erreur la plus fréquente est de soustraire les équations dans des ordres incohérents : par exemple calculer $z_{A'} - z_{B'}$ au numérateur mais $z_B - z_A$ au dénominateur, ce qui donne $-a$ au lieu de $a$. Retiens la règle : \textbf{on soustrait toujours dans le même ordre en haut et en bas} :
\[ a = \frac{z_{A'} - z_{B'}}{z_A - z_B} \quad \text{(les primes en haut, dans le même ordre que les non-primes en bas).} \]
Deuxième réflexe de survie : après avoir trouvé $a$ et $b$, \textbf{recalcule toujours} $a z_B + b$ et compare à $z_{B'}$. Cette vérification prend dix secondes et détecte toutes les erreurs de signe.
\end{attention}

\courssub{Problème type 3 : reconnaître la nature d'une transformation donnée par $z' = az + b$}

On te donne maintenant l'écriture complexe, et on te demande : « quelle est la nature et quels sont les éléments caractéristiques de cette transformation ? » Tout se lit sur le nombre $a$.

\begin{propriete}[Classification selon la valeur de $a$]
Soit $f$ la transformation du plan d'écriture complexe $z' = az + b$ avec $a \neq 0$.
\begin{itemize}
\item Si $a = 1$ : $f$ est la \cle{translation} de vecteur $\vec{u}$ d'affixe $b$.
\item Si $a \neq 1$ : $f$ est une \cle{similitude directe} de rapport $k = |a|$, d'angle $\theta = \arg(a)$ et de centre $\Omega$ d'affixe $\omega = \dfrac{b}{1-a}$. En particulier :
\begin{itemize}
\item si $a$ est réel positif ($a > 0$, $a \neq 1$) : \cle{homothétie} de rapport $a$ ;
\item si $a$ est réel négatif : homothétie de rapport $-a$ composée... plus simplement, similitude de rapport $|a|$ et d'angle $\pi$ ;
\item si $|a| = 1$ (et $a \neq 1$) : \cle{rotation} d'angle $\arg(a)$.
\end{itemize}
\end{itemize}
\end{propriete}

\begin{methode}[Réflexe de reconnaissance en trois questions]
Face à $z' = az + b$, pose-toi dans l'ordre :
\begin{enumerate}
\item $a = 1$ ? Si oui : translation de vecteur d'affixe $b$. Fin.
\item Sinon, calcule $|a|$. Si $|a| = 1$ : rotation d'angle $\arg(a)$, de centre $\omega = \dfrac{b}{1-a}$.
\item Sinon : similitude de rapport $|a|$, d'angle $\arg(a)$, de centre $\omega = \dfrac{b}{1-a}$ (et si $a \in \mathbb{R}_+^*$, c'est une homothétie).
\end{enumerate}
\end{methode}

\begin{exemplebox}[Trois reconnaissances éclair]
\begin{itemize}
\item $z' = z + 2 - 3i$ : $a = 1$, c'est la translation de vecteur d'affixe $2 - 3i$.
\item $z' = iz + 1 - i$ : $a = i$, $|a| = 1$, $\arg(a) = \dfrac{\pi}{2}$ : rotation d'angle $\dfrac{\pi}{2}$, de centre $\omega = \dfrac{1-i}{1-i} = 1$, donc $\Omega(1)$.
\item $z' = \left(\dfrac{1}{2} + i\dfrac{\sqrt{3}}{2}\right) z + 2$ : $|a| = 1$ et $\arg(a) = \dfrac{\pi}{3}$ : rotation d'angle $\dfrac{\pi}{3}$, de centre $\omega = \dfrac{2}{1 - e^{i\pi/3}}$.
\end{itemize}
\end{exemplebox}

\courssub{Lecture géométrique d'une écriture complexe}

Voici l'outil le plus puissant de cette leçon, celui qui transforme les calculs en géométrie. Toute l'idée tient en une phrase : une égalité entre différences d'affixes \emph{se lit} comme une similitude.

\begin{propriete}[Lecture géométrique fondamentale]
Soient $A$, $B$, $C$ trois points d'affixes $z_A$, $z_B$, $z_C$. Si
\[ z_C - z_A = k\,e^{i\theta}\,(z_B - z_A) \quad (k > 0), \]
alors $C$ est l'image de $B$ par la similitude directe de centre $A$, de rapport $k$ et d'angle $\theta$. En particulier :
\begin{itemize}
\item $AC = k \times AB$ (en prenant les modules) ;
\item $\left(\overrightarrow{AB}, \overrightarrow{AC}\right) = \theta \ [2\pi]$ (en prenant les arguments).
\end{itemize}
\end{propriete}

\begin{experience}[Découverte sur un cas particulier : le triangle équilatéral]
Plaçons $A$ d'affixe $0$ et $B$ d'affixe $2$, et définissons $C$ par $z_C - z_A = e^{i\pi/3}(z_B - z_A)$, c'est-à-dire $z_C = 2e^{i\pi/3} = 1 + i\sqrt{3}$.
\begin{enumerate}
\item Calcule $AC = |z_C - z_A|$ et compare à $AB$. Que constates-tu ?
\item Calcule une mesure de l'angle $\left(\overrightarrow{AB}, \overrightarrow{AC}\right)$.
\item Calcule $BC = |z_C - z_B|$. Quelle est la nature du triangle $ABC$ ?
\end{enumerate}
\emph{Réponses :} $AC = |2e^{i\pi/3}| = 2 = AB$ ; l'angle vaut $\dfrac{\pi}{3}$ ; $BC = |1 + i\sqrt{3} - 2| = |-1 + i\sqrt{3}| = 2$. Le triangle $ABC$ est \textbf{équilatéral direct}. Retiens : $z_C - z_A = e^{i\pi/3}(z_B - z_A)$ est la traduction complexe de « $ABC$ équilatéral direct ».
\end{experience}

\begin{popfigure}
\begin{tikzpicture}[scale=1.6,>=Stealth]
\coordinate (A) at (0,0);
\coordinate (B) at (2,0);
\coordinate (C) at (1,1.732);
\draw[popline] (-0.4,0) -- (2.6,0);
\draw[popline] (0,-0.4) -- (0,2.3);
\draw[thick,popblue] (A) -- (B) -- (C) -- cycle;
\fill[popblue] (A) circle (1.5pt) node[below left] {$A\ (0)$};
\fill[popblue] (B) circle (1.5pt) node[below right] {$B\ (2)$};
\fill[poporange] (C) circle (1.5pt) node[above] {$C\ (1+i\sqrt{3})$};
\draw[->,poporange,thick] (0.7,0) arc (0:60:0.7) node[midway,right] {$\frac{\pi}{3}$};
\draw[popgreen,thick] (A) -- node[midway,left,sloped] {} (C);
\node[popgreen] at (0.28,1.05) {$AC = AB$};
\end{tikzpicture}
\legende{Lecture géométrique : $z_C - z_A = e^{i\pi/3}(z_B - z_A)$ signifie que $C$ est l'image de $B$ par la rotation de centre $A$ et d'angle $\frac{\pi}{3}$ : le triangle $ABC$ est équilatéral direct.}
\end{popfigure}

\begin{methode}[Prouver qu'un triangle est rectangle isocèle]
Pour prouver qu'un triangle $ABC$ est \textbf{rectangle isocèle en $A$} (avec $\left(\overrightarrow{AB}, \overrightarrow{AC}\right) = +\dfrac{\pi}{2}$), il suffit de montrer que
\[ z_C - z_A = i\,(z_B - z_A), \]
c'est-à-dire que $C$ est l'image de $B$ par la rotation de centre $A$ et d'angle $\dfrac{\pi}{2}$. En effet, multiplier par $i$ conserve le module (donc $AC = AB$) et ajoute $\dfrac{\pi}{2}$ à l'argument (donc angle droit en $A$).

Variantes utiles : $z_C - z_A = -i(z_B - z_A)$ pour un angle de $-\dfrac{\pi}{2}$ ; $z_C - z_A = 2i(z_B - z_A)$ pour un triangle rectangle en $A$ avec $AC = 2AB$.
\end{methode}

\courssub{Stratégie : utiliser une similitude pour résoudre un problème de configuration}

Face à un problème de géométrie (alignement, perpendicularité, nature d'un triangle), la stratégie gagnante consiste à \emph{choisir} une similitude adaptée — le plus souvent de centre un sommet remarquable de la figure — et à exploiter ses propriétés de conservation.

\begin{methode}[Méthode générale de résolution]
\begin{enumerate}
\item \textbf{Identifier} dans la figure une similitude naturelle : souvent celle de centre un sommet commun à deux segments, qui envoie un point sur un autre (par exemple $s : B \mapsto C$ de centre $A$).
\item \textbf{Déterminer} son écriture complexe ou ses éléments caractéristiques.
\item \textbf{Transporter} la propriété à démontrer : une similitude conserve l'alignement, l'orthogonalité, les angles, les rapports de longueurs. Si l'image d'une configuration est plus simple, démontre la propriété sur l'image et conclus.
\item \textbf{Rédiger} la conclusion en revenant à la figure initiale.
\end{enumerate}
\end{methode}

\begin{exoresolu}[Montrer qu'un triangle est rectangle isocèle]
Le plan complexe est muni d'un repère orthonormé direct. On considère les points $A$, $B$, $C$ d'affixes $z_A = 1 + i$, $z_B = 3 + i$, $z_C = 1 + 3i$. Démontrer que le triangle $ABC$ est rectangle isocèle en $A$.
\tcblower
Calculons les différences d'affixes :
\[ z_B - z_A = (3 + i) - (1 + i) = 2, \qquad z_C - z_A = (1 + 3i) - (1 + i) = 2i. \]
Or $2i = i \times 2$, donc $z_C - z_A = i(z_B - z_A)$.

Par lecture géométrique, $C$ est l'image de $B$ par la rotation de centre $A$ et d'angle $\dfrac{\pi}{2}$. Par conséquent $AC = AB$ (la rotation conserve les distances) et $\left(\overrightarrow{AB}, \overrightarrow{AC}\right) = \dfrac{\pi}{2}$ : le triangle $ABC$ est rectangle isocèle en $A$.
\end{exoresolu}

\begin{exoresolu}[Alignement et angle par une similitude]
Soit $s$ la similitude directe d'écriture complexe $z' = (1 + i)z - 2i$. On donne $D(2)$, $E(1 + i)$ et $F(2i)$.
\begin{enumerate}
\item Déterminer les images $D'$, $E'$, $F'$ de ces points par $s$.
\item Comparer les directions des droites $(DE)$ et $(D'E')$.
\end{enumerate}
\tcblower
\textbf{1.} Calculs directs :
\begin{align*}
z_{D'} &= (1+i)(2) - 2i = 2 + 2i - 2i = 2,\\
z_{E'} &= (1+i)(1+i) - 2i = 2i - 2i = 0,\\
z_{F'} &= (1+i)(2i) - 2i = -2 + 2i - 2i = -2.
\end{align*}
Donc $D' = D(2)$, $E' = O(0)$ et $F'(-2)$.

\textbf{2.} Remarquons d'abord que $a = 1+i$ a pour module $|a| = \sqrt{2}$ et pour argument $\arg(a) = \dfrac{\pi}{4}$ : $s$ est la similitude de rapport $\sqrt{2}$ et d'angle $\dfrac{\pi}{4}$. Son centre est le point fixe $\omega = \dfrac{b}{1-a} = \dfrac{-2i}{1-(1+i)} = \dfrac{-2i}{-i} = 2$, c'est-à-dire le point $D$ (ce que l'on voit aussi car $D'=D$).

La direction de $(D'E')$ est l'image de celle de $(DE)$ par la rotation d'angle $\dfrac{\pi}{4}$ (centrée en $D$). L'angle orienté des droites $(DE)$ et $(D'E')$ vaut donc $\dfrac{\pi}{4}$.

Vérification par les affixes :
\begin{itemize}
\item Vecteur $\overrightarrow{DE}$ : $z_E - z_D = (1+i) - 2 = -1 + i$, son argument est $\dfrac{3\pi}{4}$.
\item Vecteur $\overrightarrow{D'E'}$ : $z_{E'} - z_{D'} = 0 - 2 = -2$, son argument est $\pi$.
\item Différence : $\pi - \dfrac{3\pi}{4} = \dfrac{\pi}{4}$. Cohérent.
\end{itemize}
On note au passage que $E'$, $D'$, $F'$ sont alignés sur l'axe réel (affixes $0, 2, -2$), ce qui confirme que $D, E, F$ étaient alignés (car $s$ conserve l'alignement).
\end{exoresolu}

\begin{attention}[Rédaction : ne pas confondre les rôles des points]
Quand tu écris $z_C - z_A = i(z_B - z_A)$, le \textbf{centre} est le point dont l'affixe est soustraite des deux côtés (ici $A$), le point \textbf{image} est celui dont l'affixe est isolée à gauche (ici $C$), et le point \textbf{antécédent} est dans la parenthèse (ici $B$). Inverser $B$ et $C$ change l'angle en son opposé : au Bac, cela peut faire basculer toute la conclusion (triangle rectangle isocèle « direct » ou « indirect »). Prends l'habitude de lire l'égalité à voix haute : « $C$ est l'image de $B$ par la rotation de centre $A$... ».
\end{attention}

\courssub{Complément Bac : similitudes et suites de points (hors programme strict, mais utile)}

\begin{savaistu}[Des similitudes aux fractales, et aux cases de l'Ouest camerounais]
Les similitudes directes sont le moteur mathématique des \emph{fractales} : une figure autosimilaire (flocon de von Koch, fougère de Barnsley) est obtenue en répétant indéfiniment des copies réduites et tournées d'elle-même — c'est-à-dire des images par des similitudes. On retrouve la même idée dans les motifs décoratifs africains : les gravures en spirale des cases traditionnelles de l'Ouest camerounais, les tissus Ndop des Bamilékés, ou certains motifs de vannerie, s'obtiennent en répétant un motif de base par rotation et réduction — autrement dit, par itération d'une similitude directe. Les mathématiciens parlent de \emph{spirales logarithmiques}, que l'on observe aussi dans les coquillages et les galaxies.
\end{savaistu}

Ce complément, bien que hors du programme strict, tombe régulièrement dans les sujets de Bac sous forme de question ouverte. On considère une suite de points $(M_n)$ définie par $M_0$ donné et la relation de récurrence
\[ z_{n+1} = a z_n + b, \]
où $z' = az + b$ est l'écriture d'une similitude directe $s$ de centre $\Omega$. Alors chaque point est l'image du précédent : $M_{n+1} = s(M_n)$. En retranchant l'égalité de point fixe $\omega = a\omega + b$, on obtient
\[ z_{n+1} - \omega = a(z_n - \omega), \]
donc la suite $(z_n - \omega)$ est \textbf{géométrique de raison $a$} : $z_n - \omega = a^n (z_0 - \omega)$. Géométriquement, chaque point s'obtient à partir du précédent en tournant de $\theta = \arg(a)$ autour de $\Omega$ et en multipliant la distance à $\Omega$ par $k = |a|$. Si $k < 1$, les points $M_n$ \emph{spiralent} vers $\Omega$ ; si $k > 1$, ils s'en éloignent en spirale ; si $k = 1$, ils tournent sur un cercle de centre $\Omega$.

\begin{popfigure}
\begin{tikzpicture}[scale=2.6,>=Stealth]
\coordinate (O) at (0,0);
\coordinate (M0) at (1.6,0);
\coordinate (M1) at (0.980,0.823);
\coordinate (M2) at (0.297,0.957);
\coordinate (M3) at (-0.252,0.736);
\coordinate (M4) at (-0.573,0.376);
\coordinate (M5) at (-0.635,0.028);
\draw[popline,->] (-1,0) -- (2,0);
\draw[popline,->] (0,-0.6) -- (0,1.4);
\draw[thick,popblue] (M0) -- (M1) -- (M2) -- (M3) -- (M4) -- (M5);
\foreach \p/\n/\pos in {M0/{M_0}/right, M1/{M_1}/above right, M2/{M_2}/above, M3/{M_3}/above left, M4/{M_4}/left, M5/{M_5}/below left}{
  \fill[poporange] (\p) circle (1.3pt) node[\pos,popdark] {$\n$};
}
\fill[popgreen] (O) circle (1.6pt) node[below left] {$\Omega$};
\draw[->,popgreen] (0.5,0) arc (0:40:0.5) node[midway,right] {$40^\circ$};
\end{tikzpicture}
\legende{Suite de points $M_{n+1} = s(M_n)$ où $s$ est la similitude de centre $\Omega$, de rapport $0{,}8$ et d'angle $40^\circ$ : les points spiralent vers le centre $\Omega$.}
\end{popfigure}

\begin{exercice}[Pour s'entraîner]
On considère la suite $(M_n)$ de points d'affixes définies par $z_0 = 8$ et $z_{n+1} = \dfrac{1+i}{2}\,z_n$ pour tout $n \in \mathbb{N}$.
\begin{enumerate}
\item Déterminer la nature et les éléments caractéristiques de la similitude $s$ telle que $M_{n+1} = s(M_n)$.
\item Montrer que $z_n = 8\left(\dfrac{\sqrt{2}}{2}\right)^n e^{in\pi/4}$.
\item Déterminer la limite de $|z_n|$ et interpréter géométriquement.
\end{enumerate}
\end{exercice}

\begin{corrige}[Exercice 1]
\textbf{1.} L'écriture est de la forme $z' = az$ avec $a = \dfrac{1+i}{2}$, donc $b = 0$ et le centre est $\Omega = O$ (origine). On a $|a| = \dfrac{\sqrt{2}}{2}$ et $\arg(a) = \dfrac{\pi}{4}$ : $s$ est la similitude de centre $O$, de rapport $\dfrac{\sqrt{2}}{2}$ et d'angle $\dfrac{\pi}{4}$.

\textbf{2.} La suite $(z_n)$ est géométrique de raison $a$, donc $z_n = 8a^n$. Or $a = \dfrac{\sqrt{2}}{2}e^{i\pi/4}$, donc $a^n = \left(\dfrac{\sqrt{2}}{2}\right)^n e^{in\pi/4}$, d'où le résultat.

\textbf{3.} $|z_n| = 8\left(\dfrac{\sqrt{2}}{2}\right)^n$. Comme $0 < \dfrac{\sqrt{2}}{2} < 1$, on a $\lim |z_n| = 0$ : les points $M_n$ spiralent vers l'origine $O$, en tournant d'un angle de $\dfrac{\pi}{4}$ à chaque étape.
\end{corrige}

\begin{aretenir}[L'essentiel de la section]
\begin{itemize}
\item Similitude de centre $\Omega(\omega)$, rapport $k$, angle $\theta$ : $z' - \omega = ke^{i\theta}(z - \omega)$.
\item Similitude envoyant $A$ sur $A'$ et $B$ sur $B'$ : résoudre le système, avec $a = \dfrac{z_{A'} - z_{B'}}{z_A - z_B}$, puis vérifier.
\item Reconnaissance de $z' = az + b$ : $a = 1$ translation ; $|a| = 1$ rotation ; $a \in \mathbb{R}_+^*$ homothétie ; sinon similitude de rapport $|a|$, angle $\arg(a)$, centre $\dfrac{b}{1-a}$.
\item Lecture géométrique : $z_C - z_A = ke^{i\theta}(z_B - z_A) \iff C = s(B)$ avec $s$ de centre $A$, rapport $k$, angle $\theta$. Cas clé : $z_C - z_A = i(z_B - z_A) \iff ABC$ rectangle isocèle en $A$.
\item Suite $z_{n+1} = az_n + b$ : les points itérés spiralent autour du point fixe $\Omega$.
\end{itemize}
\end{aretenir}

\newpage

\coursec{Similitudes : lieux géométriques, constructions et démonstrations}

Dans la leçon commune, tu as appris à reconnaître une similitude directe donnée par $z' = az + b$, à déterminer son centre, son rapport et son angle, et à trouver l'image d'une droite ou d'un cercle. Nous allons maintenant franchir un pas : utiliser les similitudes comme un \emph{outil de résolution}, pour transformer des lieux géométriques, construire des figures et démontrer des propriétés. C'est exactement l'esprit des problèmes de géométrie du Baccalauréat.

\courssub{Transformation d'un lieu géométrique}

\begin{propriete}[Image d'un lieu par une similitude directe]
Soit $s$ une similitude directe de rapport $k > 0$ et d'angle $\theta$.
\begin{itemize}
\item L'image d'une droite $(d)$ est une droite $(d')$ ; de plus, pour tout vecteur directeur $\vec{u}$ de $(d)$, un vecteur directeur $\vec{u'}$ de $(d')$ vérifie $\left(\vec{u}, \vec{u'}\right) = \theta \ [2\pi]$.
\item L'image d'un cercle $\mathcal{C}$ de centre $\Omega$ et de rayon $R$ est le cercle $\mathcal{C}'$ de centre $\Omega' = s(\Omega)$ et de rayon $kR$.
\item L'image d'un lieu de points $\mathcal{L}$ est le lieu $\mathcal{L}' = \{s(M) \mid M \in \mathcal{L}\}$ : on transforme l'\emph{ensemble}, pas seulement quelques points.
\end{itemize}
\end{propriete}

\begin{methode}[Déterminer l'image d'un lieu géométrique]
\begin{enumerate}
\item Identifier la nature du lieu de départ (droite, cercle, segment, arc de cercle...).
\item Déterminer les éléments caractéristiques de la similitude : rapport $k = |a|$, angle $\theta = \arg(a)$, centre $\omega = \frac{b}{1-a}$ (si $a \neq 1$).
\item Conclure sur la nature du lieu image, puis préciser ses éléments : direction de la droite image (tournée de $\theta$), centre et rayon du cercle image.
\item Pour une équation cartésienne ou complexe du lieu image, utiliser l'écriture réciproque $z = \dfrac{z' - b}{a}$ et substituer dans la condition définissant le lieu.
\end{enumerate}
\end{methode}

\begin{exoresolu}[Image d'un cercle par une similitude]
Le plan est rapporté à un repère orthonormé direct $(O\,;\vec{u},\vec{v})$. Soit $s$ la transformation d'écriture complexe $z' = (1 + \mathrm{i})z + 2$. Déterminer l'image du cercle $\mathcal{C}$ de centre $A(1\,;\,0)$ et de rayon $2$.
\tcblower
\textbf{Solution.} On a $a = 1 + \mathrm{i}$, donc $|a| = \sqrt{2} \neq 1$ : $s$ est une similitude directe de rapport $k = \sqrt{2}$ et d'angle $\theta = \arg(1+\mathrm{i}) = \dfrac{\pi}{4}$.

L'image d'un cercle est un cercle. L'affixe de $A$ est $z_A = 1$, donc :
\[ z_{A'} = (1+\mathrm{i})\times 1 + 2 = 3 + \mathrm{i}. \]
Le rayon image est $R' = kR = 2\sqrt{2}$.

\textbf{Conclusion :} l'image de $\mathcal{C}$ est le cercle de centre $A'(3\,;\,1)$ et de rayon $2\sqrt{2}$.
\end{exoresolu}

\begin{attention}[Piège : confondre point image et lieu image]
Une erreur classique consiste à calculer l'image d'\emph{un seul} point du lieu et à conclure trop vite. Pour l'image d'un cercle, il faut transformer le \cle{centre} \emph{et} multiplier le \cle{rayon} par $k$. Autre piège : le centre de la similitude n'est \textbf{pas} en général le centre du cercle image ; c'est le centre du cercle de départ qu'il faut transformer.
\end{attention}

\courssub{Constructions à l'aide d'une similitude}

L'idée directrice est simple : pour construire un point $M$ soumis à des conditions, on introduit une similitude $s$ qui envoie une figure connue sur la figure cherchée. Les conditions imposées à $M$ se traduisent alors par : $M$ appartient à un lieu $\mathcal{L}$ \emph{et} à l'image réciproque $s^{-1}(\mathcal{L}')$ d'un autre lieu. Le point cherché est à l'intersection.

\begin{methode}[Résoudre un problème de construction]
\begin{enumerate}
\item Traduire les données pour identifier la similitude $s$ (souvent définie par deux couples de points $s(A) = A'$, $s(B) = B'$).
\item Exprimer chaque condition sur le point cherché $M$ comme une appartenance à un lieu (droite, cercle).
\item Transformer l'un des lieux par $s$ (ou $s^{-1}$) pour ramener les conditions à des lieux constructibles à la règle et au compas (droites, cercles).
\item Construire l'intersection, puis vérifier que la solution convient (il peut y avoir zéro, une ou plusieurs solutions).
\end{enumerate}
\end{methode}

\begin{exemplebox}[Construire un carré inscrit dans un triangle]
Soit un triangle $ABC$. On veut construire un carré $MNPQ$ tel que $M, N \in [AB]$, $P \in [BC]$ et $Q \in [AC]$.
\emph{Astuce :} on construit d'abord un carré $mnpq$ « flottant » avec $m, n \in (AB)$ et $q \in (AC)$, puis l'homothétie $h$ de centre $A$ qui envoie $p$ sur le point $P'$ où $(Ap)$ coupe $(BC)$ envoie le petit carré sur le carré solution. Une homothétie étant une similitude directe, l'image d'un carré est un carré : la construction est justifiée rigoureusement.
\end{exemplebox}

\begin{popfigure}
\begin{center}
\begin{tikzpicture}[scale=1.1]
% Triangle ABC
\coordinate (A) at (0,0);
\coordinate (B) at (5,0);
\coordinate (C) at (2,3);
\draw[thick,popdark] (A)--(B)--(C)--cycle;

% Small square mnpq : m on AB, q on AC
% Choose m at 20% of AB
\coordinate (m) at ($(A)!0.2!(B)$);
% q = intersection of perpendicular to AB at m and line AC
\path[name path=perp_m] (m) -- ++(0,3);
\path[name path=AC] (A) -- (C);
\path[name intersections={of=perp_m and AC}] (intersection-1) coordinate (q);

% Complete square m-n-p-q (direct sense)
% n = rotation of q around m by -90 deg
\coordinate (n) at ($(m)!1!-90:(q)$);
% p = rotation of m around q by +90 deg (or translation)
\coordinate (p) at ($(q)!1!-90:(m)$);

\draw[popblue,thick] (m)--(n)--(p)--(q)--cycle;

% Homothety center A sending p to P on BC
\path[name path=Ap] (A) -- (p);
\path[name path=BC] (B) -- (C);
\path[name intersections={of=Ap and BC}] (intersection-1) coordinate (P);

% Q = intersection of horizontal through P (parallel to pq) and AC
% Since pq is horizontal (AB horizontal), PQ is horizontal.
\path[name path=horiz_P] (P) -- ++(-5,0);
\path[name intersections={of=horiz_P and AC}] (intersection-1) coordinate (Q);

% M, N projections of Q, P on AB
\coordinate (M) at ($(A)!(Q)!(B)$);
\coordinate (N) at ($(A)!(P)!(B)$);

\draw[poporange,very thick] (M)--(N)--(P)--(Q)--cycle;
\draw[dashed,popmuted] (A)--(P);
\draw[dashed,popmuted] (A)--(Q);

% Labels
\fill[popink] (A) circle (1.2pt) node[below left]{$A$};
\fill[popink] (B) circle (1.2pt) node[below right]{$B$};
\fill[popink] (C) circle (1.2pt) node[above]{$C$};
\fill[poporange] (M) circle (1.2pt) node[below]{$M$};
\fill[poporange] (N) circle (1.2pt) node[below]{$N$};
\fill[poporange] (P) circle (1.2pt) node[right]{$P$};
\fill[poporange] (Q) circle (1.2pt) node[left]{$Q$};
\fill[popblue] (m) circle (1pt) node[below,scale=0.7]{$m$};
\fill[popblue] (n) circle (1pt) node[below,scale=0.7]{$n$};
\fill[popblue] (p) circle (1pt) node[right,scale=0.7]{$p$};
\fill[popblue] (q) circle (1pt) node[left,scale=0.7]{$q$};
\end{tikzpicture}
\end{center}
\legende{Le petit carré bleu est envoyé par l'homothétie de centre $A$ sur le carré orange solution.}
\end{popfigure}

\courssub{Démontrer avec une similitude}

Les similitudes directes conservent les angles orientés et multiplient les distances par $k$. Elles sont donc parfaites pour démontrer des propriétés d'alignement, de perpendicularité, de cocyclicité ou d'égalité d'angles.

\begin{methode}[Rédiger une démonstration par similitude]
\begin{enumerate}
\item Deviner ou exhiber la similitude $s$ qui envoie une configuration sur une autre (souvent définie par $s(A)=A'$ et $s(B)=B'$).
\item Justifier son existence et son unicité : deux couples de points distincts définissent une unique similitude directe.
\item Utiliser les conservations : $s$ conserve les angles orientés, donc transforme droites perpendiculaires en droites perpendiculaires, et préserve l'alignement.
\item Conclure sur la propriété demandée en revenant à la configuration initiale.
\end{enumerate}
\end{methode}

\begin{exoresolu}[Perpendicularité démontrée par une rotation]
$ABCD$ est un carré de centre $O$, de sens direct. Soit $r$ la rotation de centre $O$ et d'angle $\dfrac{\pi}{2}$. Montrer que les diagonales $(AC)$ et $(BD)$ sont perpendiculaires.
\tcblower
\textbf{Solution.} Une rotation est une similitude directe de rapport $1$. Comme $ABCD$ est un carré de sens direct, $r(A) = B$ et $r(C) = D$ (le centre $O$ est fixe et $OA = OB$, $(\overrightarrow{OA},\overrightarrow{OB}) = \dfrac{\pi}{2}$).

L'image de la droite $(AC)$ par $r$ est donc la droite $(BD)$. Or une rotation d'angle $\dfrac{\pi}{2}$ envoie toute droite sur une droite qui lui est perpendiculaire (les vecteurs directeurs tournent de $\dfrac{\pi}{2}$). Donc $(AC) \perp (BD)$. \hfill$\blacksquare$
\end{exoresolu}

\begin{attention}[Rédaction : l'existence avant l'usage]
Dans une démonstration, on ne peut pas utiliser une similitude avant d'avoir justifié qu'elle existe : précise toujours ses éléments caractéristiques ou les deux couples de points qui la définissent. Et rappelle-toi : une similitude directe conserve les angles \emph{orientés} ; une similitude indirecte (hors programme ici) les changerait en leur opposé.
\end{attention}

\begin{savaistu}[Des pavages aux cartes numériques]
Les logiciels de cartographie utilisés pour les plans de Douala ou de Yaoundé reposent sur des similitudes : zoomer sur une carte, c'est appliquer une homothétie ; orienter la carte vers le nord, c'est composer avec une rotation. Les distances sont multipliées par un rapport constant, les angles (carrefours, rues perpendiculaires) sont préservés : la carte reste fidèle à la réalité.
\end{savaistu}

\begin{exercice}[Lieu transformé]
Soit $s$ la similitude directe d'écriture complexe $z' = 2\mathrm{i}\,z + 1 - \mathrm{i}$. On note $\mathcal{D}$ la droite d'équation $y = x - 1$. Déterminer la nature et une équation de l'image de $\mathcal{D}$ par $s$.
\end{exercice}

\begin{corrige}[Exercice 1]
On a $a = 2\mathrm{i}$, donc $k = 2$ et $\theta = \dfrac{\pi}{2}$ : l'image de $\mathcal{D}$ est une droite $\mathcal{D}'$ perpendiculaire à $\mathcal{D}$. De $z' = 2\mathrm{i}z + 1 - \mathrm{i}$, on tire $z = \dfrac{z' - 1 + \mathrm{i}}{2\mathrm{i}}$. En écrivant $z = x + \mathrm{i}y$ et $z' = x' + \mathrm{i}y'$, on obtient $x = \dfrac{y' + 1}{2}$ et $y = \dfrac{1 - x'}{2}$. La condition $y = x - 1$ devient $\dfrac{1 - x'}{2} = \dfrac{y' + 1}{2} - 1$, soit $1 - x' = y' - 1$, c'est-à-dire $y' = -x' + 2$. L'image est la droite d'équation $y = -x + 2$, bien perpendiculaire à $\mathcal{D}$ (coefficients directeurs $1$ et $-1$).
\end{corrige}

\newpage

\coursec{Activité d'intégration}

\coursec{Activité d'intégration : Le motif du pagne de Foumban}

\courssub{Objectifs de l'activité}

À l'issue de cette activité, tu seras capable de :
\begin{itemize}
\item déterminer la \cle{nature} et les \cle{éléments caractéristiques} (centre, rapport, angle) d'une similitude directe donnée par son écriture complexe ;
\item calculer les \cle{affixes des images} de points par une similitude directe et construire l'image d'une configuration usuelle (triangle équilatéral) ;
\item déterminer l'\cle{image d'un cercle} par une similitude directe ;
\item \cle{composer} une similitude avec elle-même (itération) et interpréter géométriquement le résultat (spirale) ;
\item utiliser les propriétés métriques d'une similitude pour résoudre un problème concret de seuil (inéquation exponentielle).
\end{itemize}

\courssub{Situation}

À Foumban, capitale du royaume Bamoun, les artisans tisserands sont réputés pour leurs pagnes décorés de motifs géométriques répétés en spirale. Mama Aïcha, artisane de la cité des arts, veut créer un nouveau pagne dont le motif de base est un triangle équilatéral $ABC$ brodé au fil doré. Pour obtenir l'effet « tourbillon » caractéristique, chaque motif suivant est l'image du précédent par une même transformation : une réduction accompagnée d'une rotation.

Pour préparer son travail, elle modélise le problème sur une feuille quadrillée munie d'un repère orthonormé direct $(O\,;\vec{u},\vec{v})$ d'unité \SI{1}{cm}. Le motif de base est le triangle $ABC$ dont les sommets ont pour affixes :
\[ z_A = 0, \qquad z_B = 2, \qquad z_C = 1 + i\sqrt{3}. \]
La transformation utilisée pour passer d'un motif au suivant est la similitude directe $s$ d'écriture complexe :
\[ z' = \frac{1}{2}(1+i)\,z + 1. \]
Mama Aïcha veut savoir où placer exactement les motifs successifs, où sera brodé le petit cercle décoratif inscrit dans chaque triangle, et surtout \textbf{combien de motifs} elle peut faire tenir sur la bande de tissu avant que le motif ne devienne trop petit pour être brodé (elle ne sait plus broder un motif dont le rayon est inférieur à \SI{0,1}{cm}).

\courssub{Consignes}

\begin{enumerate}
\item \textbf{Partie A -- Étude de la transformation} : Vérifier que le triangle $ABC$ est équilatéral. Donner la nature de $s$, son rapport $k$, son angle $\theta$ et l'affixe $\omega$ de son centre $\Omega$.
\item \textbf{Partie B -- Construction du deuxième motif} : Calculer les affixes $z_{A'}$, $z_{B'}$, $z_{C'}$ des images de $A$, $B$, $C$ par $s$. Justifier que $A'B'C'$ est équilatéral et donner la longueur de son côté. Sur papier millimétré, placer $\Omega$, $A$, $B$, $C$, $A'$, $B'$, $C'$ et vérifier graphiquement que $(\overrightarrow{\Omega A},\overrightarrow{\Omega A'}) = \frac{\pi}{4}$ et $\Omega A' = \frac{\sqrt{2}}{2}\Omega A$.
\item \textbf{Partie C -- Le cercle décoratif} : Déterminer l'affixe du centre $G$ du cercle inscrit dans $ABC$ (centre de gravité) et son rayon $r$. En déduire le centre et le rayon du cercle image par $s$.
\item \textbf{Partie D -- Itération et spirale} : Calculer les affixes des sommets du troisième motif $A''B''C''$. Exprimer la distance $\Omega A^{(n)}$ en fonction de $n$ et expliquer l'enroulement en spirale. Déterminer le plus petit entier $n$ tel que le rayon du $n$-ième motif (défini comme $\max(\Omega A^{(n-1)},\Omega B^{(n-1)},\Omega C^{(n-1)})$) soit strictement inférieur à \SI{0,1}{cm}. Conclure pour Mama Aïcha.
\end{enumerate}

\courssub{Ressources à mobiliser}

\begin{propriete}[Écriture complexe d'une similitude directe]
Toute similitude directe $s$ du plan admet une écriture complexe de la forme $z' = a z + b$ avec $a \in \mathbb{C}^*$, $b \in \mathbb{C}$.
\begin{itemize}
\item Son \cle{rapport} est $k = |a|$.
\item Son \cle{angle} est $\theta = \arg(a) \pmod{2\pi}$.
\item Si $a \neq 1$, $s$ admet un unique point fixe, son \cle{centre} $\Omega$ d'affixe $\omega = \dfrac{b}{1-a}$. La transformation est la composée de l'homothétie $h(\Omega, k)$ et de la rotation $r(\Omega, \theta)$ (dans n'importe quel ordre).
\item Si $a = 1$, $s$ est la translation de vecteur d'affixe $b$.
\end{itemize}
\end{propriete}

\begin{propriete}[Propriétés métriques et conservation]
Une similitude directe de rapport $k$ multiplie toutes les distances par $k$ et \cle{conserve les angles orientés} (et les rapports de longueurs).
\begin{itemize}
\item L'image d'un segment est un segment de longueur multipliée par $k$.
\item L'image d'un triangle est un triangle semblable (mêmes angles, côtés proportionnels au rapport $k$). Un triangle équilatéral reste équilatéral.
\item L'image d'un cercle $\mathcal{C}(G, r)$ est le cercle $\mathcal{C}(s(G), k\,r)$.
\item L'image d'une droite est une droite ; l'image d'une demi-droite est une demi-droite.
\end{itemize}
\end{propriete}

\begin{aretenir}[Itération d'une similitude directe]
Soit $s$ une similitude directe de centre $\Omega$, rapport $k$ et angle $\theta$ ($k>0$, $k\neq 1$). Pour tout entier $n \ge 1$, la composée $s^n = \underbrace{s \circ s \circ \dots \circ s}_{n \text{ fois}}$ est la similitude directe de centre $\Omega$, de rapport $k^n$ et d'angle $n\theta \pmod{2\pi}$.
Si $k < 1$, les itérées $s^n(M)$ convergent vers $\Omega$ en décrivant une \cle{spirale logarithmique}.
\end{aretenir}

\courssub{Démarche et corrigé commenté}

\begin{exoresolu}[Le motif du pagne de Foumban : corrigé détaillé]
\textbf{Rappel des questions.}
\begin{enumerate}
\item Vérifier que $ABC$ est équilatéral. Nature, rapport, angle, centre de $s$.
\item Affixes de $A', B', C'$ ; nature et côté de $A'B'C'$ ; construction graphique et vérifications.
\item Centre $G$ et rayon $r$ du cercle inscrit dans $ABC$ ; image par $s$.
\item Affixes de $A'', B'', C''$ ; expression de $\Omega A^{(n)}$ ; explication spirale ; seuil $0,1$ cm.
\end{enumerate}
\tcblower
\textbf{Partie A -- Étude de la transformation}

\textbf{1. Triangle $ABC$ équilatéral.}
On calcule les longueurs des côtés à l'aide du module :
\begin{align*}
AB &= |z_B - z_A| = |2 - 0| = 2,\\
AC &= |z_C - z_A| = |1 + i\sqrt{3}| = \sqrt{1^2 + (\sqrt{3})^2} = \sqrt{4} = 2,\\
BC &= |z_C - z_B| = |1 + i\sqrt{3} - 2| = |-1 + i\sqrt{3}| = \sqrt{(-1)^2 + (\sqrt{3})^2} = 2.
\end{align*}
Les trois côtés mesurent $2$ cm, donc $ABC$ est un triangle équilatéral.

\textbf{2. Nature et éléments caractéristiques de $s$.}
L'écriture complexe est $z' = a z + b$ avec $a = \dfrac{1+i}{2}$ et $b = 1$.
Comme $a \neq 1$, $s$ est une \cle{similitude directe} (ni translation, ni homothétie pure, ni rotation pure).
\begin{itemize}
\item \cle{Rapport} : $k = |a| = \left|\dfrac{1+i}{2}\right| = \dfrac{|1+i|}{2} = \dfrac{\sqrt{2}}{2}$.
\item \cle{Angle} : $\theta = \arg(a) = \arg(1+i) = \dfrac{\pi}{4} \pmod{2\pi}$ (car $a$ est un réel positif $\frac{1}{2}$ fois $1+i$).
\end{itemize}

\textbf{3. Centre $\Omega$ de la similitude.}
Le centre est l'unique point fixe : $\omega = a\omega + b \iff \omega(1-a) = b \iff \omega = \dfrac{b}{1-a}$.
\[
\omega = \frac{1}{1 - \frac{1+i}{2}} = \frac{1}{\frac{1-i}{2}} = \frac{2}{1-i} = \frac{2(1+i)}{(1-i)(1+i)} = \frac{2(1+i)}{2} = 1+i.
\]
Le centre $\Omega$ a pour affixe $\omega = 1+i$, soit les coordonnées $(1\,;\,1)$ dans le repère.

\bigskip
\textbf{Partie B -- Construction du deuxième motif}

\textbf{4. Affixes des images $A', B', C'$.}
On applique $z' = \frac{1+i}{2}z + 1$ :
\begin{align*}
z_{A'} &= \frac{1+i}{2}\times 0 + 1 = 1,\\[4pt]
z_{B'} &= \frac{1+i}{2}\times 2 + 1 = (1+i) + 1 = 2 + i,\\[4pt]
z_{C'} &= \frac{1+i}{2}(1 + i\sqrt{3}) + 1 
= \frac{1}{2}\bigl[(1-\sqrt{3}) + i(1+\sqrt{3})\bigr] + 1 \\
&= \frac{3-\sqrt{3}}{2} + i\,\frac{1+\sqrt{3}}{2}.
\end{align*}

\textbf{5. Nature et dimensions de $A'B'C'$.}
Puisque $s$ est une similitude directe de rapport $k = \frac{\sqrt{2}}{2}$, elle conserve les angles et multiplie les longueurs par $k$.
$ABC$ est équilatéral $\implies$ $A'B'C'$ est équilatéral.
Longueur du côté : $A'B' = k \times AB = \frac{\sqrt{2}}{2} \times 2 = \sqrt{2}$ cm.

\textbf{6. Construction graphique et vérifications.}

\begin{popfigure}
\begin{tikzpicture}[scale=1.2, >=Stealth]
% Axes
\draw[popmuted, thin] (-0.5,-0.5) grid (3,2.5);
\draw[->, popdark] (-0.5,0) -- (3.2,0) node[right] {$x$};
\draw[->, popdark] (0,-0.5) -- (0,2.7) node[above] {$y$};
\foreach \x in {0,1,2,3} \draw (\x,0.05) -- (\x,-0.05) node[below, font=\tiny] {\x};
\foreach \y in {0,1,2} \draw (0.05,\y) -- (-0.05,\y) node[left, font=\tiny] {\y};

% Points originaux
\coordinate (O) at (0,0);
\coordinate (Om) at (1,1);
\coordinate (A) at (0,0);
\coordinate (B) at (2,0);
\coordinate (C) at (1,{sqrt(3)});
\coordinate (Ap) at (1,0);
\coordinate (Bp) at (2,1);
\coordinate (Cp) at ({(3-sqrt(3))/2},{(1+sqrt(3))/2});

% Triangle ABC
\draw[popblue, thick] (A) -- (B) -- (C) -- cycle;
% Triangle A'B'C'
\draw[poporange, thick, dashed] (Ap) -- (Bp) -- (Cp) -- cycle;

% Centre Omega
\fill[poppurple] (Om) circle (2pt) node[above left] {$\Omega(1,1)$};

% Points labels
\fill[popblue] (A) circle (1.5pt) node[below left] {$A(0,0)$};
\fill[popblue] (B) circle (1.5pt) node[below right] {$B(2,0)$};
\fill[popblue] (C) circle (1.5pt) node[above] {$C(1,\sqrt{3})$};

\fill[poporange] (Ap) circle (1.5pt) node[below] {$A'(1,0)$};
\fill[poporange] (Bp) circle (1.5pt) node[right] {$B'(2,1)$};
\fill[poporange] (Cp) circle (1.5pt) node[above left] {$C'$};

% Vecteurs Omega A et Omega A'
\draw[popgreen, ->, thick] (Om) -- (A) node[midway, left] {$\overrightarrow{\Omega A}$};
\draw[popgreen, ->, thick] (Om) -- (Ap) node[midway, right] {$\overrightarrow{\Omega A'}$};

% Angle arc
\draw[popgreen, thick] (Om) ++(225:0.5) arc (225:270:0.5) node[midway, below left, font=\small] {$\frac{\pi}{4}$};

% Distances
\draw[popmuted, densely dotted] (Om) -- (A);
\draw[popmuted, densely dotted] (Om) -- (Ap);
\end{tikzpicture}
\legende{Construction des deux premiers motifs : $ABC$ (bleu) et $A'B'C'$ (orange pointillé). Le centre $\Omega$ est en violet. L'angle $(\overrightarrow{\Omega A},\overrightarrow{\Omega A'}) = \pi/4$ et le rapport $\Omega A'/\Omega A = \sqrt{2}/2$ sont visibles.}
\end{popfigure}

\textit{Vérifications numériques :}
\begin{itemize}
\item $\Omega A = |z_A - \omega| = |0 - (1+i)| = \sqrt{2}$.
\item $\Omega A' = |z_{A'} - \omega| = |1 - (1+i)| = |-i| = 1$.
\item Rapport : $\dfrac{\Omega A'}{\Omega A} = \dfrac{1}{\sqrt{2}} = \dfrac{\sqrt{2}}{2} = k$. \checkmark
\item Angle : $\arg\left(\dfrac{z_{A'}-\omega}{z_A-\omega}\right) = \arg\left(\dfrac{-i}{-1-i}\right) = \arg\left(\dfrac{i}{1+i}\right) = \arg\left(\dfrac{i(1-i)}{2}\right) = \arg\left(\dfrac{1+i}{2}\right) = \dfrac{\pi}{4} = \theta$. \checkmark
\end{itemize}

\bigskip
\textbf{Partie C -- Le cercle décoratif}

\textbf{7. Cercle inscrit dans $ABC$.}
Dans un triangle équilatéral, le centre du cercle inscrit (incentre) coïncide avec le centre de gravité (isobarycentre des sommets).
\[
z_G = \frac{z_A + z_B + z_C}{3} = \frac{0 + 2 + (1+i\sqrt{3})}{3} = \frac{3 + i\sqrt{3}}{3} = 1 + i\frac{\sqrt{3}}{3}.
\]
Le rayon $r$ du cercle inscrit dans un triangle équilatéral de côté $a=2$ est :
\[
r = \frac{a\sqrt{3}}{6} = \frac{2\sqrt{3}}{6} = \frac{\sqrt{3}}{3} \text{ cm} \approx 0,577 \text{ cm}.
\]

\textbf{8. Image du cercle par $s$.}
L'image est un cercle de centre $G' = s(G)$ et de rayon $r' = k\,r$.
\begin{align*}
z_{G'} &= \frac{1+i}{2}\left(1 + i\frac{\sqrt{3}}{3}\right) + 1 \\
&= \frac{1}{2}\left(1 + i\frac{\sqrt{3}}{3} + i - \frac{\sqrt{3}}{3}\right) + 1 \\
&= \frac{1}{2}\left(1 - \frac{\sqrt{3}}{3}\right) + 1 + i\,\frac{1}{2}\left(1 + \frac{\sqrt{3}}{3}\right) \\
&= \frac{3 - \sqrt{3}}{6} + 1 + i\,\frac{3 + \sqrt{3}}{6} = \frac{9 - \sqrt{3}}{6} + i\,\frac{3 + \sqrt{3}}{6}.
\end{align*}
\[
r' = k\,r = \frac{\sqrt{2}}{2} \times \frac{\sqrt{3}}{3} = \frac{\sqrt{6}}{6} \approx 0,408 \text{ cm}.
\]

\bigskip
\textbf{Partie D -- Itération et spirale}

\textbf{9. Troisième motif $A''B''C'' = s(A'B'C')$.}
On applique $s$ aux affixes de $A', B', C'$ :
\begin{align*}
z_{A''} &= \frac{1+i}{2}(1) + 1 = \frac{3+i}{2},\\[4pt]
z_{B''} &= \frac{1+i}{2}(2+i) + 1 = \frac{1}{2}(2+i+2i-1) + 1 = \frac{1+3i}{2} + 1 = \frac{3+3i}{2},\\[4pt]
z_{C''} &= \frac{1+i}{2}\,z_{C'} + 1 
= \frac{1+i}{2}\left(\frac{3-\sqrt{3}}{2} + i\frac{1+\sqrt{3}}{2}\right) + 1 \\
&= \frac{1}{4}\Bigl[(1+i)(3-\sqrt{3}) + i(1+i)(1+\sqrt{3})\Bigr] + 1 \\
&= \frac{1}{4}\Bigl[ (3-\sqrt{3}) + i(3-\sqrt{3}) + i(1+\sqrt{3}) - (1+\sqrt{3}) \Bigr] + 1 \\
&= \frac{1}{4}\Bigl[ 2 - 2\sqrt{3} + 4i \Bigr] + 1 = \frac{1-\sqrt{3}}{2} + i + 1 = \frac{3-\sqrt{3}}{2} + i.
\end{align*}

\textbf{10. Distance $\Omega A^{(n)}$ et spirale.}
Notons $A^{(n)}$ l'image de $A$ par $s^n$ ($A^{(0)}=A$, $A^{(1)}=A'$, $A^{(2)}=A''$, etc.).
Puisque $s$ est une similitude de centre $\Omega$ et de rapport $k$, pour tout point $M$, $\Omega s(M) = k \cdot \Omega M$.
Par récurrence immédiate : $\Omega A^{(n)} = k^n \cdot \Omega A = \left(\frac{\sqrt{2}}{2}\right)^n \sqrt{2}$.
De plus, l'angle $(\overrightarrow{\Omega A^{(n-1)}}, \overrightarrow{\Omega A^{(n)}}) = \theta = \frac{\pi}{4}$.
Comme $k = \frac{\sqrt{2}}{2} < 1$, les distances à $\Omega$ diminuent géométriquement vers $0$, tandis que les points tournent d'un angle constant $\frac{\pi}{4}$ à chaque étape. Les sommets successifs décrivent donc une \cle{spirale logarithmique} (ou spirale de Bernoulli) convergente vers $\Omega$.

\textbf{11. Seuil de broderie (question de synthèse).}
Le « rayon » d'un motif est défini comme la plus grande distance de $\Omega$ à ses sommets.
Pour le motif de base ($n=1$, triangle $ABC$), les distances sont :
$\Omega A = \sqrt{2}$, $\Omega B = |2-(1+i)| = |1-i| = \sqrt{2}$, $\Omega C = |1+i\sqrt{3}-1-i| = \sqrt{3}-1 \approx 0,73$.
Le maximum est $R_1 = \sqrt{2}$ cm (atteint pour $A$ et $B$).

Le $n$-ième motif est l'image du motif de base par $s^{n-1}$. Son rayon est donc :
\[
R_n = k^{n-1} R_1 = \left(\frac{\sqrt{2}}{2}\right)^{n-1} \sqrt{2}.
\]
On cherche le plus petit entier $n$ tel que $R_n < 0,1$.
\begin{align*}
\sqrt{2}\left(\frac{\sqrt{2}}{2}\right)^{n-1} &< 0,1 \\
\left(\frac{\sqrt{2}}{2}\right)^{n-1} &< \frac{0,1}{\sqrt{2}} = \frac{\sqrt{2}}{20} \approx 0,0707.
\end{align*}
On applique le logarithme népérien (fonction strictement croissante) ; comme $\ln(\sqrt{2}/2) = -\frac{1}{2}\ln 2 < 0$, le sens de l'inégalité s'inverse :
\[
(n-1)\ln\left(\frac{\sqrt{2}}{2}\right) < \ln\left(\frac{\sqrt{2}}{20}\right)
\iff n-1 > \frac{\ln(\sqrt{2}/20)}{\ln(\sqrt{2}/2)}.
\]
Calcul numérique :
$\ln(\sqrt{2}/2) \approx -0,3466$, $\ln(\sqrt{2}/20) \approx -2,649$.
$n-1 > \dfrac{-2,649}{-0,3466} \approx 7,64 \implies n \ge 9$.

\textit{Vérification exacte sans calculatrice :}
$R_n = \sqrt{2} \cdot 2^{-(n-1)/2} = 2^{1/2 - (n-1)/2} = 2^{(2-n)/2} = 2^{1 - n/2}$.
$R_8 = 2^{1-4} = 2^{-3} = \frac{1}{8} = 0,125 > 0,1$.
$R_9 = 2^{1-4,5} = 2^{-3,5} = \frac{1}{2^{3,5}} = \frac{1}{8\sqrt{2}} \approx 0,0884 < 0,1$.

\medskip
\noindent\textbf{Conclusion :} Mama Aïcha pourra broder \textbf{8 motifs} complets (du 1er au 8ème) sur sa bande de tissu. Le 9ème motif aurait un rayon inférieur à \SI{0,1}{cm} et serait trop fin pour être brodé.
\end{exoresolu}

\courssub{Bilan de l'activité}

Cette activité a permis de mobiliser l'ensemble des savoir-faire de la leçon « Similitudes directes du plan » dans un contexte authentique et culturellement ancré :
\begin{itemize}
\item La lecture d'une écriture complexe $z' = az + b$ donne immédiatement la \cle{nature} (similitude directe car $a\notin\mathbb{R}_+^*$ et $a\neq 1$), le \cle{rapport} $k=|a|$, l'\cle{angle} $\theta=\arg(a)$, et le \cle{centre} $\omega = b/(1-a)$.
\item Les images de configurations usuelles se déduisent directement des propriétés de conservation : un triangle équilatéral reste équilatéral (côtés multipliés par $k$), un cercle reste un cercle (centre image, rayon multiplié par $k$).
\item L'\cle{itération} d'une similitude de rapport $k<1$ produit une spirale convergente vers le centre ; les problèmes de « seuil » (taille minimale) se ramènent à une inéquation exponentielle $k^{n-1}R_1 < \varepsilon$, résolue rigoureusement grâce au logarithme.
\end{itemize}
Tu as ainsi découvert comment les mathématiques de Terminale permettent de modéliser — et de prévoir avec précision — les créations artisanales du patrimoine camerounais.

\newpage

\coursec{Méthodes et exercices résolus}

\coursec{Similitudes directes du plan}

\courssub{Méthode 1 : Déterminer l'écriture complexe d'une similitude directe à partir de ses éléments caractéristiques}

\begin{methode}[De « centre, rapport, angle » vers $z' = az + b$]
Une similitude directe $s$ est donnée par son \cle{centre} $\Omega$ d'affixe $\omega$, son \cle{rapport} $k > 0$ et son \cle{angle} $\theta$.
\begin{enumerate}
\item Écrire la relation vectorielle complexe caractéristique :
\[ z' - \omega = k e^{i\theta}(z - \omega). \]
\item Identifier $a = k e^{i\theta}$ : on a donc $|a| = k$ et $\arg(a) \equiv \theta \ [2\pi]$.
\item Développer pour obtenir la forme $z' = az + b$ avec $b = \omega - a\omega = \omega(1 - a)$.
\item Vérifier : $b = \omega(1-a)$ doit être cohérent avec le fait que $\Omega$ est le point fixe ($\omega = a\omega + b$).
\end{enumerate}
\textbf{Cas particuliers à reconnaître :} $a = 1$ : translation ; $a \in \mathbb{R}_+^* \setminus \{1\}$ : homothétie ; $|a| = 1, a \neq 1$ : rotation.
\end{methode}

\begin{exoresolu}[Écriture complexe à partir des éléments caractéristiques]
Le plan est muni d'un repère orthonormé direct $(O\,;\vec{u},\vec{v})$. Soit $s$ la similitude directe de centre $\Omega$ d'affixe $\omega = 1 + 2i$, de rapport $k = 2$ et d'angle $\theta = \dfrac{\pi}{3}$.
\begin{enumerate}
\item Déterminer l'écriture complexe de $s$.
\item Déterminer l'affixe de l'image $A'$ du point $A$ d'affixe $z_A = -1 + i$.
\end{enumerate}
\tcblower
\textbf{1. Écriture complexe de $s$.}

On applique la relation caractéristique d'une similitude directe de centre $\Omega(\omega)$, de rapport $k$ et d'angle $\theta$ :
\[ z' - \omega = k e^{i\theta}(z - \omega). \]
Calculons le coefficient $a = k e^{i\theta}$ :
\[ a = 2 e^{i\pi/3} = 2\left(\cos\frac{\pi}{3} + i\sin\frac{\pi}{3}\right) = 2\left(\frac{1}{2} + i\frac{\sqrt{3}}{2}\right) = 1 + i\sqrt{3}. \]
On a donc :
\[ z' - (1 + 2i) = (1 + i\sqrt{3})\bigl(z - (1+2i)\bigr). \]
Développons pour obtenir la forme $z' = az + b$ :
\begin{align*}
z' &= (1 + i\sqrt{3})z - (1+i\sqrt{3})(1+2i) + (1+2i) \\
&= (1 + i\sqrt{3})z + (1+2i)\bigl(1 - (1+i\sqrt{3})\bigr) \\
&= (1 + i\sqrt{3})z + (1+2i)(-i\sqrt{3}).
\end{align*}
Or $(1+2i)(-i\sqrt{3}) = -i\sqrt{3} - 2i^2\sqrt{3} = 2\sqrt{3} - i\sqrt{3}$. D'où :
\[ \boxed{z' = (1 + i\sqrt{3})\,z + 2\sqrt{3} - i\sqrt{3}}. \]
\textbf{Vérification du point fixe :} $a\omega + b = (1+i\sqrt{3})(1+2i) + 2\sqrt{3} - i\sqrt{3} = 1 + 2i + i\sqrt{3} - 2\sqrt{3} + 2\sqrt{3} - i\sqrt{3} = 1 + 2i = \omega$. $\Omega$ est bien invariant.

\textbf{2. Image de $A(-1+i)$.}
\begin{align*}
z_{A'} &= (1 + i\sqrt{3})(-1 + i) + 2\sqrt{3} - i\sqrt{3} \\
&= -1 + i - i\sqrt{3} + i^2\sqrt{3} + 2\sqrt{3} - i\sqrt{3} \\
&= -1 - \sqrt{3} + 2\sqrt{3} + i\bigl(1 - \sqrt{3} - \sqrt{3}\bigr) \\
&= (\sqrt{3} - 1) + i(1 - 2\sqrt{3}).
\end{align*}
\textbf{Conclusion :} l'écriture complexe de $s$ est $z' = (1+i\sqrt{3})z + 2\sqrt{3} - i\sqrt{3}$, et l'image de $A$ a pour affixe $z_{A'} = (\sqrt{3}-1) + i(1 - 2\sqrt{3})$.
\end{exoresolu}

\courssub{Méthode 2 : Déterminer la nature et les éléments caractéristiques d'une transformation donnée par $z' = az + b$}

\begin{methode}[De $z' = az + b$ vers la nature de la transformation]
Soit $f$ la transformation d'écriture complexe $z' = az + b$ avec $a \in \mathbb{C}^*$, $b \in \mathbb{C}$.
\begin{enumerate}
\item \textbf{Comparer $a$ à $1$.}
\begin{itemize}
\item Si $a = 1$ : $f$ est la \cle{translation} de vecteur d'affixe $b$ (l'identité si $b = 0$).
\item Si $a \neq 1$ : $f$ est une \cle{similitude directe} de rapport $k = |a|$ et d'angle $\theta \equiv \arg(a) \ [2\pi]$.
\end{itemize}
\item \textbf{Préciser la nature :}
\begin{itemize}
\item si $a \in \mathbb{R}_+^* \setminus \{1\}$ : homothétie de rapport $a$ ;
\item si $a \in \mathbb{R}_-^*$ : homothétie de rapport $|a|$ (angle $\pi$) ;
\item si $|a| = 1$ et $a \neq 1$ : rotation d'angle $\arg(a)$ ;
\item sinon : similitude directe « générale » (composée d'une homothétie et d'une rotation de même centre).
\end{itemize}
\item \textbf{Déterminer le centre} (cas $a \neq 1$) : c'est l'unique point fixe $\Omega$, d'affixe solution de $\omega = a\omega + b$, soit
\[ \omega = \frac{b}{1 - a}. \]
\item Conclure en rédigeant une phrase complète : « $f$ est la similitude directe de centre $\Omega(\omega)$, de rapport $k = |a|$ et d'angle $\theta = \arg(a)$ ».
\end{enumerate}
\end{methode}

\begin{exoresolu}[Nature et éléments caractéristiques]
Le plan complexe est rapporté à un repère orthonormé direct $(O\,;\vec{u},\vec{v})$. On considère la transformation $f$ du plan qui à tout point $M$ d'affixe $z$ associe le point $M'$ d'affixe :
\[ z' = \left(\frac{1}{2} - i\frac{\sqrt{3}}{2}\right) z + 3 + i\sqrt{3}. \]
Déterminer la nature de $f$ et préciser ses éléments caractéristiques.
\tcblower
Posons $a = \dfrac{1}{2} - i\dfrac{\sqrt{3}}{2}$ et $b = 3 + i\sqrt{3}$.

\textbf{Étape 1 : comparaison de $a$ à $1$.} On a $a \neq 1$, donc $f$ est une similitude directe admettant un unique point fixe.

\textbf{Étape 2 : rapport.} 
\[ k = |a| = \sqrt{\left(\frac{1}{2}\right)^2 + \left(\frac{\sqrt{3}}{2}\right)^2} = \sqrt{\frac{1}{4} + \frac{3}{4}} = 1. \]
Comme $|a| = 1$ et $a \neq 1$, $f$ est une \textbf{rotation}.

\textbf{Étape 3 : angle.} On reconnaît $a = \cos\left(-\dfrac{\pi}{3}\right) + i\sin\left(-\dfrac{\pi}{3}\right) = e^{-i\pi/3}$. Donc
\[ \theta = \arg(a) \equiv -\frac{\pi}{3} \ [2\pi]. \]

\textbf{Étape 4 : centre.} Le centre $\Omega$ est l'unique point fixe, d'affixe $\omega = \dfrac{b}{1-a}$ :
\[ 1 - a = 1 - \frac{1}{2} + i\frac{\sqrt{3}}{2} = \frac{1}{2} + i\frac{\sqrt{3}}{2} = e^{i\pi/3}. \]
Donc
\[ \omega = \frac{3 + i\sqrt{3}}{\frac{1}{2} + i\frac{\sqrt{3}}{2}}. \]
Multiplions numérateur et dénominateur par le conjugué du dénominateur :
\begin{align*}
\omega &= \frac{(3 + i\sqrt{3})\left(\frac{1}{2} - i\frac{\sqrt{3}}{2}\right)}{\left(\frac{1}{2}\right)^2 + \left(\frac{\sqrt{3}}{2}\right)^2} = \frac{\frac{3}{2} - i\frac{3\sqrt{3}}{2} + i\frac{\sqrt{3}}{2} - i^2\frac{3}{2}}{1} \\
&= \frac{3}{2} + \frac{3}{2} + i\left(-\frac{3\sqrt{3}}{2} + \frac{\sqrt{3}}{2}\right) = 3 - i\sqrt{3}.
\end{align*}
\textbf{Vérification :} $a\omega + b = \left(\frac{1}{2} - i\frac{\sqrt{3}}{2}\right)(3 - i\sqrt{3}) + 3 + i\sqrt{3} = \frac{3}{2} - i\frac{3\sqrt{3}}{2} - i\frac{\sqrt{3}}{2} + i^2\frac{3}{2} + 3 + i\sqrt{3} = 3 - i\sqrt{3} = \omega$. Correct.

\textbf{Conclusion :} $f$ est la \textbf{rotation} de centre $\Omega$ d'affixe $\omega = 3 - i\sqrt{3}$ et d'angle $-\dfrac{\pi}{3}$.
\end{exoresolu}

\courssub{Méthode 3 : Déterminer l'image d'un point par une similitude directe}

\begin{methode}[Calcul de l'image d'un point]
Pour déterminer l'affixe $z'$ de l'image $M'$ d'un point $M$ d'affixe $z$ par la similitude $s : z' = az + b$ :
\begin{enumerate}
\item Substituer l'affixe de $M$ dans l'expression $z' = az + b$.
\item Effectuer le produit $az$ en développant : $(\alpha + i\beta)(x + iy) = (\alpha x - \beta y) + i(\alpha y + \beta x)$, en pensant à $i^2 = -1$.
\item Regrouper parties réelle et imaginaire, puis ajouter $b$.
\item Si l'on connaît plutôt les éléments caractéristiques $(\Omega, k, \theta)$, on peut aussi utiliser $z' - \omega = ke^{i\theta}(z - \omega)$, souvent plus rapide.
\end{enumerate}
\textbf{Réflexe de contrôle :} vérifier que $\Omega M' = k \cdot \Omega M$, c'est-à-dire $|z' - \omega| = k\,|z - \omega|$.
\end{methode}

\begin{exoresolu}[Image d'un point et contrôle géométrique]
Soit $s$ la similitude directe d'écriture complexe $z' = (1 + i)z + 2 - i$. On donne le point $B$ d'affixe $z_B = 1 - 2i$.
\begin{enumerate}
\item Déterminer l'affixe de $B'$, image de $B$ par $s$.
\item Déterminer le centre $\Omega$ de $s$, puis vérifier que $\Omega B' = \sqrt{2}\,\Omega B$.
\end{enumerate}
\tcblower
\textbf{1. Affixe de $B'$.}
\begin{align*}
z_{B'} &= (1+i)(1 - 2i) + 2 - i \\
&= 1 - 2i + i - 2i^2 + 2 - i \\
&= 1 + 2 + 2 + i(-2 + 1 - 1) \qquad (\text{car } -2i^2 = 2)\\
&= 5 - 2i.
\end{align*}
Donc $z_{B'} = 5 - 2i$.

\textbf{2. Centre et vérification.} Ici $a = 1+i \neq 1$, donc le centre est le point fixe :
\[ \omega = \frac{b}{1-a} = \frac{2 - i}{1 - (1+i)} = \frac{2 - i}{-i} = \frac{(2-i)\,i}{-i \cdot i} = \frac{2i - i^2}{1} = 1 + 2i. \]
Le rapport est $k = |a| = |1+i| = \sqrt{2}$. Vérifions :
\[ z_B - \omega = (1 - 2i) - (1 + 2i) = -4i, \qquad \Omega B = |{-4i}| = 4 ; \]
\[ z_{B'} - \omega = (5 - 2i) - (1 + 2i) = 4 - 4i, \qquad \Omega B' = |4 - 4i| = 4\sqrt{2}. \]
On a bien $\Omega B' = 4\sqrt{2} = \sqrt{2} \times 4 = k \cdot \Omega B$. Le contrôle est concluant.

\textbf{Conclusion :} $B'$ a pour affixe $5 - 2i$, et le centre de $s$ est $\Omega(1 + 2i)$, de rapport $\sqrt{2}$ (et d'angle $\arg(1+i) = \frac{\pi}{4}$).
\end{exoresolu}

\courssub{Méthode 4 : Déterminer l'image d'une droite ou d'un cercle par une similitude directe}

\begin{methode}[Image d'une droite, image d'un cercle]
\textbf{Image d'une droite $\mathcal{D}$ :}
\begin{enumerate}
\item Choisir deux points distincts $A$ et $B$ de $\mathcal{D}$.
\item Calculer leurs images $A' = s(A)$ et $B' = s(B)$.
\item Conclure : $s(\mathcal{D})$ est la droite $(A'B')$. De plus, si $\theta$ est l'angle de $s$, alors $\left(\overrightarrow{AB}\,;\overrightarrow{A'B'}\right) \equiv \theta \ [2\pi]$ : la droite image fait l'angle $\theta$ avec la droite initiale.
\end{enumerate}
\textbf{Image d'un cercle $\mathcal{C}$ de centre $I$ et de rayon $R$ :}
\begin{enumerate}
\item Calculer $I' = s(I)$.
\item L'image est le cercle de centre $I'$ et de rayon $R' = kR$, où $k$ est le rapport de $s$.
\end{enumerate}
\textbf{Alternative analytique :} exprimer $z$ en fonction de $z'$ via $z = \dfrac{z' - b}{a}$ et substituer dans l'équation de l'ensemble de départ.
\end{methode}

\begin{exoresolu}[Image d'une droite et d'un cercle]
Soit $s$ la similitude directe d'écriture complexe $z' = 2i\,z + 1 - i$.
\begin{enumerate}
\item Déterminer les éléments caractéristiques de $s$.
\item Soit $\mathcal{D}$ la droite passant par $A(0)$ et $B(1)$. Déterminer l'image de $\mathcal{D}$ par $s$.
\item Soit $\mathcal{C}$ le cercle de centre $I(1 + i)$ et de rayon $3$. Déterminer l'image de $\mathcal{C}$ par $s$.
\end{enumerate}
\tcblower
\textbf{1. Éléments caractéristiques.} On a $a = 2i \neq 1$ : $s$ est une similitude directe de rapport $k = |2i| = 2$ et d'angle $\theta = \arg(2i) \equiv \dfrac{\pi}{2}\ [2\pi]$. Son centre est le point fixe :
\[ \omega = \frac{b}{1-a} = \frac{1 - i}{1 - 2i} = \frac{(1-i)(1+2i)}{(1-2i)(1+2i)} = \frac{1 + 2i - i - 2i^2}{1 + 4} = \frac{3 + i}{5}. \]
Donc $\Omega\left(\dfrac{3}{5} + \dfrac{1}{5}i\right)$, rapport $2$, angle $\dfrac{\pi}{2}$.

\textbf{2. Image de la droite $\mathcal{D} = (AB)$.}
\[ z_{A'} = 2i \times 0 + 1 - i = 1 - i \ ; \qquad z_{B'} = 2i \times 1 + 1 - i = 1 + i. \]
L'image de $\mathcal{D}$ est la droite $(A'B')$ passant par $A'(1 - i)$ et $B'(1 + i)$ : c'est la droite d'équation $x = 1$ (droite verticale). Cohérent avec l'angle $\frac{\pi}{2}$ : l'horizontale $(AB)$ est envoyée sur une verticale.

\textbf{3. Image du cercle $\mathcal{C}$.} Le centre $I(1+i)$ a pour image :
\[ z_{I'} = 2i(1+i) + 1 - i = 2i + 2i^2 + 1 - i = -1 + i. \]
Le rayon image est $R' = kR = 2 \times 3 = 6$.

\textbf{Conclusion :} $s(\mathcal{D})$ est la droite d'équation $x = 1$ ; $s(\mathcal{C})$ est le cercle de centre $I'(-1 + i)$ et de rayon $6$.
\end{exoresolu}

\courssub{Méthode 5 : Déterminer une similitude directe connaissant deux points et leurs images}

\begin{methode}[Similitude définie par $s(A) = A'$ et $s(B) = B'$]
Une similitude directe est entièrement déterminée par la donnée de deux points distincts $A$, $B$ et de leurs images $A'$, $B'$.
\begin{enumerate}
\item Poser $z' = az + b$ et écrire le système :
\[ \begin{cases} z_{A'} = a\,z_A + b \\ z_{B'} = a\,z_B + b \end{cases} \]
\item Soustraire membre à membre pour éliminer $b$ :
\[ a = \frac{z_{A'} - z_{B'}}{z_A - z_B} \qquad (z_A \neq z_B). \]
\item En déduire $b = z_{A'} - a\,z_A$.
\item Conclure sur les éléments caractéristiques : rapport $k = |a| = \dfrac{A'B'}{AB}$, angle $\theta \equiv \arg(a) \equiv \left(\overrightarrow{AB}\,;\overrightarrow{A'B'}\right)\ [2\pi]$, et centre $\omega = \dfrac{b}{1-a}$ si $a \neq 1$.
\end{enumerate}
\end{methode}

\begin{exoresolu}[Similitude déterminée par deux couples de points]
Dans le plan complexe, on donne les points $A(1)$, $B(2 + i)$, $A'(2i)$ et $B'(-2 + 3i)$.
\begin{enumerate}
\item Montrer qu'il existe une unique similitude directe $s$ telle que $s(A) = A'$ et $s(B) = B'$, et déterminer son écriture complexe.
\item Donner le rapport, l'angle et le centre de $s$.
\end{enumerate}
\tcblower
\textbf{1. Écriture complexe de $s$.} Comme $A \neq B$, il existe une unique similitude directe envoyant $A$ sur $A'$ et $B$ sur $B'$ (résultat du cours). Cherchons $a$ et $b$ tels que $z' = az + b$ :
\[ \begin{cases} 2i = a \times 1 + b \\ -2 + 3i = a(2+i) + b \end{cases} \]
En soustrayant la première équation à la deuxième :
\[ -2 + 3i - 2i = a(2 + i) - a \iff -2 + i = a(1 + i) \iff a = \frac{-2 + i}{1 + i}. \]
Calculons :
\[ a = \frac{(-2+i)(1-i)}{(1+i)(1-i)} = \frac{-2 + 2i + i - i^2}{2} = \frac{-1 + 3i}{2} = -\frac{1}{2} + \frac{3}{2}i. \]
Puis $b = 2i - a = 2i + \dfrac{1}{2} - \dfrac{3}{2}i = \dfrac{1}{2} + \dfrac{1}{2}i$.
\[ \boxed{z' = \left(-\frac{1}{2} + \frac{3}{2}i\right) z + \frac{1}{2} + \frac{1}{2}i}. \]
\textbf{Vérification pour $B$ :} $a(2+i) + b = \left(-\frac{1}{2} + \frac{3}{2}i\right)(2+i) + \frac{1}{2} + \frac{1}{2}i = -1 - \frac{1}{2}i + 3i + \frac{3}{2}i^2 + \frac{1}{2} + \frac{1}{2}i = -1 - \frac{3}{2} + \frac{1}{2} + i\left(-\frac{1}{2} + 3 + \frac{1}{2}\right) = -2 + 3i = z_{B'}$. Correct.

\textbf{2. Éléments caractéristiques.}
\[ k = |a| = \sqrt{\frac{1}{4} + \frac{9}{4}} = \sqrt{\frac{10}{4}} = \frac{\sqrt{10}}{2}. \]
Pour l'angle : $a = \dfrac{-1 + 3i}{2}$. On a $\cos\theta = \dfrac{-1/2}{\sqrt{10}/2} = -\dfrac{1}{\sqrt{10}}$ et $\sin\theta = \dfrac{3/2}{\sqrt{10}/2} = \dfrac{3}{\sqrt{10}} > 0$, donc $\theta$ est l'angle du second quadrant défini par ces valeurs : $\theta = \pi - \arctan(3)$.
Le centre :
\[ \omega = \frac{b}{1-a} = \frac{\frac{1}{2} + \frac{1}{2}i}{\frac{3}{2} - \frac{3}{2}i} = \frac{1 + i}{3(1 - i)} = \frac{(1+i)(1+i)}{3 \times 2} = \frac{2i}{6} = \frac{i}{3}. \]
\textbf{Conclusion :} $s$ est la similitude directe de centre $\Omega\left(\dfrac{i}{3}\right)$, de rapport $\dfrac{\sqrt{10}}{2}$ et d'angle $\theta = \pi - \arctan 3$.
\end{exoresolu}

\courssub{Méthode 6 : Utiliser une similitude directe pour résoudre un problème de configuration}

\begin{methode}[Similitudes et configurations : le plan d'attaque]
Face à un problème de configuration (alignement, orthogonalité, triangle de forme imposée, cocyclicité) :
\begin{enumerate}
\item \textbf{Repérer la similitude cachée} : deux figures « de même forme » (triangles directement semblables) suggèrent une similitude directe envoyant l'une sur l'autre.
\item \textbf{Exploiter les propriétés de conservation} : une similitude directe conserve les angles orientés, le parallélisme, l'orthogonalité, le barycentre, et multiplie les distances par $k$.
\item \textbf{Traduire en complexe} : $s(M) = M' \iff z' = az + b$ ; l'angle de la similitude se lit dans $\arg(a)$, le rapport dans $|a|$.
\item \textbf{Conclure géométriquement} : revenir à la configuration (longueurs, angles, nature d'un triangle) en interprétant $|z' - z|$ comme une distance et $\arg\!\left(\dfrac{z_C - z_A}{z_B - z_A}\right)$ comme un angle.
\end{enumerate}
\textbf{Réflexe clé :} le triangle $ABC$ est rectangle isocèle direct en $A$ si et seulement si $z_C - z_A = i(z_B - z_A)$, c'est-à-dire si $C$ est l'image de $B$ par la rotation de centre $A$ et d'angle $\frac{\pi}{2}$.
\end{methode}

\begin{exoresolu}[Triangle isocèle et similitude]
Dans le plan complexe muni du repère orthonormé direct $(O\,;\vec{u},\vec{v})$, on considère les points $A(1)$, $B(3 + 2i)$ et $C(-1 + 4i)$.
\begin{enumerate}
\item Déterminer la nature du triangle $ABC$.
\item Soit $s$ la similitude directe de centre $A$ qui transforme $B$ en $C$. Déterminer son écriture complexe et ses éléments caractéristiques.
\item Soit $D$ le point tel que $ABDC$ soit un parallélogramme. Déterminer l'image de $D$ par $s$ et en déduire la nature du quadrilatère $AC'D'C$ où $C' = s(C)$ et $D' = s(D)$.
\end{enumerate}
\tcblower
\textbf{1. Nature du triangle $ABC$.} Calculons les affixes des vecteurs :
\[ z_B - z_A = 2 + 2i \ ; \qquad z_C - z_A = -2 + 4i \ ; \qquad z_C - z_B = -4 + 2i. \]
Les longueurs des côtés sont :
\[ AB = |2 + 2i| = 2\sqrt{2} \ ; \quad AC = |-2 + 4i| = \sqrt{4 + 16} = 2\sqrt{5} \ ; \quad BC = |-4 + 2i| = \sqrt{16 + 4} = 2\sqrt{5}. \]
On a $AC = BC = 2\sqrt{5}$ : le triangle est \textbf{isocèle en $C$}.
Vérifions l'orthogonalité en $A$ : $\dfrac{z_C - z_A}{z_B - z_A} = \dfrac{-2+4i}{2+2i} = \dfrac{(-2+4i)(2-2i)}{8} = \dfrac{-4 + 4i + 8i + 8}{8} = \dfrac{4 + 12i}{8} = \dfrac{1}{2} + \dfrac{3}{2}i$, qui n'est pas un imaginaire pur. Le triangle n'est donc pas rectangle en $A$.

\textbf{2. Similitude de centre $A$ envoyant $B$ sur $C$.} Une telle similitude vérifie $z_C - z_A = a(z_B - z_A)$, donc
\[ a = \frac{z_C - z_A}{z_B - z_A} = \frac{1}{2} + \frac{3}{2}i \quad (\text{calculé ci-dessus}). \]
L'écriture complexe est $z' - z_A = a(z - z_A)$, soit :
\[ z' = \left(\frac{1}{2} + \frac{3}{2}i\right) z + 1 - \left(\frac{1}{2} + \frac{3}{2}i\right) = \left(\frac{1}{2} + \frac{3}{2}i\right) z + \frac{1}{2} - \frac{3}{2}i. \]
Rapport : $k = |a| = \sqrt{\dfrac{1}{4} + \dfrac{9}{4}} = \dfrac{\sqrt{10}}{2}$ ; angle : $\theta = \arg(a) = \arctan 3$ (premier quadrant). Centre : $A(1)$.

\textbf{3. Image de $D$ et nature du quadrilatère.} $ABDC$ parallélogramme signifie $\overrightarrow{AB} = \overrightarrow{CD}$, d'où $z_D = z_C + z_B - z_A = -1 + 4i + 3 + 2i - 1 = 1 + 6i$.
\[ z_{D'} = \left(\frac{1}{2} + \frac{3}{2}i\right)(1 + 6i) + \frac{1}{2} - \frac{3}{2}i = \frac{1}{2} + 3i + \frac{3}{2}i + 9i^2 + \frac{1}{2} - \frac{3}{2}i = -8 + 3i. \]
Et $C' = s(C)$ : $z_{C'} = \left(\frac{1}{2} + \frac{3}{2}i\right)(-1+4i) + \frac{1}{2} - \frac{3}{2}i = -\frac{1}{2} + 2i - \frac{3}{2}i - 6 + \frac{1}{2} - \frac{3}{2}i = -6 - i$.
Comme une similitude directe conserve le parallélisme et transforme tout parallélogramme en parallélogramme, l'image du parallélogramme $ABDC$ est le parallélogramme $A C' D' C$ (car $s(A) = A$, $s(B) = C$, $s(D) = D'$, $s(C) = C'$). De plus, les rapports de longueurs étant multipliés par $k = \frac{\sqrt{10}}{2}$ et les angles conservés, ce parallélogramme a pour côtés $AC' = k \cdot AB = \frac{\sqrt{10}}{2} \times 2\sqrt{2} = 2\sqrt{5}$ et $CC' = k \cdot BC = \frac{\sqrt{10}}{2} \times 2\sqrt{5} = 5$.

\textbf{Conclusion :} le triangle $ABC$ est isocèle en $C$ ; $s$ a pour écriture $z' = \left(\frac{1}{2} + \frac{3}{2}i\right)z + \frac{1}{2} - \frac{3}{2}i$, de centre $A$, rapport $\frac{\sqrt{10}}{2}$, angle $\arctan 3$ ; l'image du parallélogramme $ABDC$ est le parallélogramme de sommets $A(1)$, $C'(-6-i)$, $D'(-8+3i)$, $C(-1+4i)$.
\end{exoresolu}

\courssub{Méthode 7 : Passer de l'écriture complexe à l'expression analytique (et réciproquement)}

\begin{methode}[Écriture complexe $\leftrightarrow$ expression analytique]
\textbf{Du complexe vers l'analytique :} soit $z' = az + b$ avec $a = \alpha + i\beta$, $b = \gamma + i\delta$, $z = x + iy$, $z' = x' + iy'$.
\begin{enumerate}
\item Développer : $x' + iy' = (\alpha + i\beta)(x + iy) + \gamma + i\delta$.
\item Identifier parties réelle et imaginaire :
\[ \begin{cases} x' = \alpha x - \beta y + \gamma \\ y' = \beta x + \alpha y + \delta \end{cases} \]
\end{enumerate}
\textbf{De l'analytique vers le complexe :} un système $\begin{cases} x' = \alpha x - \beta y + \gamma \\ y' = \beta x + \alpha y + \delta \end{cases}$ (même coefficient $\alpha$ sur la diagonale, coefficients opposés $\pm\beta$) correspond à $z' = az + b$ avec $a = \alpha + i\beta$ et $b = \gamma + i\delta$.
\begin{attention}[Forme reconnaissable]
Tout système de la forme $\begin{cases} x' = \alpha x - \beta y + \gamma \\ y' = \beta x + \alpha y + \delta \end{cases}$ définit une similitude directe de rapport $\sqrt{\alpha^2 + \beta^2}$. Si les coefficients n'ont pas cette structure (par exemple $x' = \alpha x + \beta y$, $y' = \beta x + \alpha y$), ce n'est \textbf{pas} une similitude directe.
\end{attention}
\end{methode}

\begin{exoresolu}[Passage complexe--analytique]
Soit $s$ la transformation du plan d'écriture complexe $z' = (1 - 2i)z + 3 + i$.
\begin{enumerate}
\item Déterminer l'expression analytique de $s$.
\item En déduire les coordonnées de l'image du point $M(2\,;-1)$.
\item Un point $N(x\,;y)$ a pour image $N'(0\,;5)$. Déterminer les coordonnées de $N$.
\end{enumerate}
\tcblower
\textbf{1. Expression analytique.} Posons $z = x + iy$ et $z' = x' + iy'$ :
\begin{align*}
x' + iy' &= (1 - 2i)(x + iy) + 3 + i \\
&= x + iy - 2ix - 2i^2 y + 3 + i \\
&= x + 2y + 3 + i(y - 2x + 1).
\end{align*}
Par identification des parties réelle et imaginaire :
\[ \boxed{\begin{cases} x' = x + 2y + 3 \\ y' = -2x + y + 1 \end{cases}} \]
On vérifie la structure attendue : $\alpha = 1$, $\beta = -2$ (car $a = 1 - 2i$), avec $x' = \alpha x - \beta y + \gamma = x + 2y + 3$. Correct.

\textbf{2. Image de $M(2\,;-1)$.}
\[ x' = 2 + 2(-1) + 3 = 3 \ ; \qquad y' = -2(2) + (-1) + 1 = -4. \]
Donc $M'(3\,;-4)$.

\textbf{3. Antécédent de $N'(0\,;5)$.} On résout le système :
\[ \begin{cases} x + 2y + 3 = 0 \\ -2x + y + 1 = 5 \end{cases} \iff \begin{cases} x + 2y = -3 \\ -2x + y = 4 \end{cases} \]
De la première équation : $x = -3 - 2y$. En substituant dans la deuxième : $-2(-3 - 2y) + y = 4 \iff 6 + 4y + y = 4 \iff 5y = -2 \iff y = -\dfrac{2}{5}$. Puis $x = -3 - 2\left(-\dfrac{2}{5}\right) = -3 + \dfrac{4}{5} = -\dfrac{11}{5}$.
\textbf{Vérification :} $x + 2y + 3 = -\frac{11}{5} - \frac{4}{5} + 3 = 0$ ; $-2x + y + 1 = \frac{22}{5} - \frac{2}{5} + 1 = 5$. Correct.

\textbf{Conclusion :} $s$ a pour expression analytique $\begin{cases} x' = x + 2y + 3 \\ y' = -2x + y + 1 \end{cases}$ ; $M(2;-1)$ a pour image $M'(3;-4)$, et $N'(0;5)$ a pour antécédent $N\left(-\dfrac{11}{5}\,;-\dfrac{2}{5}\right)$.
\end{exoresolu}

\begin{attention}[Les 5 erreurs qui coûtent des points au Bac]
\begin{enumerate}
\item \textbf{Oublier de discuter le cas $a = 1$.} Avant d'affirmer « le centre est $\omega = \dfrac{b}{1-a}$ », il faut avoir vérifié que $a \neq 1$. Si $a = 1$, la transformation est une translation (pas de centre !) et diviser par $1 - a$ est une division par zéro : faute grave.
\item \textbf{Confondre le rapport avec $a$ et l'angle avec $\arg(b)$.} Le rapport est $k = |a|$ (un réel strictement positif) et l'angle est $\theta \equiv \arg(a)\ [2\pi]$. Le nombre $b$ ne donne ni le rapport ni l'angle : il sert uniquement à déterminer le centre (ou le vecteur de la translation).
\item \textbf{Se tromper dans le calcul du point fixe.} L'équation est $\omega = a\omega + b$, d'où $\omega = \dfrac{b}{1-a}$ et non $\dfrac{b}{a-1}$ ni $\dfrac{-b}{1-a}$. Erreur de signe très fréquente ; vérifie toujours en recalculant $a\omega + b = \omega$.
\item \textbf{Annoncer un angle sans justifier le quadrant.} Écrire « $\theta = \arctan\left(\frac{\beta}{\alpha}\right)$ » est faux quand $\alpha < 0$. Il faut donner $\cos\theta = \dfrac{\alpha}{|a|}$ et $\sin\theta = \dfrac{\beta}{|a|}$ et conclure sur le bon quadrant, ou écrire $a$ sous forme exponentielle $ke^{i\theta}$.
\item \textbf{Rédaction incomplète de la conclusion.} « $s$ est une similitude » ne suffit pas : le Bac exige la \textbf{nature précise} et \textbf{tous} les éléments caractéristiques (centre d'affixe $\omega$, rapport $k$, angle $\theta$), plus une phrase de conclusion. De même, pour l'image d'un cercle, donner le centre image sans le rayon $kR$ fait perdre des points.
\end{enumerate}
\end{attention}

\newpage

\coursec{Exercices}

\coursec{Similitudes directes du plan}

\courssub{Je vérifie mes connaissances}

\begin{exercice}[QCM : nature d'une transformation]
Pour chaque écriture complexe $z' = az + b$, indiquer la nature de la transformation associée (translation, homothétie, rotation, similitude directe non réduite aux cas précédents). Une seule réponse par ligne.
\begin{enumerate}
\item $z' = 3z - 2 + i$ ;
\item $z' = z + 4 - 2i$ ;
\item $z' = -iz + 1$ ;
\item $z' = (1 + i)z$.
\end{enumerate}
\end{exercice}

\begin{exercice}[Vrai ou faux, en justifiant]
Dire si chacune des affirmations suivantes est vraie ou fausse, en justifiant la réponse.
\begin{enumerate}
\item Toute homothétie de rapport strictement positif est une similitude directe.
\item La transformation d'écriture complexe $z' = 2\bar{z} + 1$ est une similitude directe.
\item Une similitude directe de rapport $k$ transforme toute droite en une droite qui lui est parallèle.
\item La composée de deux rotations est toujours une rotation.
\item Une similitude directe qui fixe deux points distincts est l'identité.
\end{enumerate}
\end{exercice}

\begin{attention}[Pièges classiques sur la nature des transformations]
\begin{itemize}
\item La forme $z' = az + b$ avec $a \in \mathbb{C}^*$ caractérise les \textbf{similitudes directes}. Si $|a|=1$ et $a \neq 1$, c'est une rotation ; si $a \in \mathbb{R}^*_+$ et $a \neq 1$, c'est une homothétie ; si $a=1$, c'est une translation.
\item La présence de $\bar{z}$ (conjugaison) signale une similitude \textbf{indirecte} (symétrie axiale, inversion, etc.), jamais une similitude directe.
\item L'image d'une droite par une similitude directe est une droite qui forme avec l'originale un angle égal à l'angle de la similitude (modulo $\pi$). Elles sont parallèles \textbf{uniquement} si l'angle est $0$ ou $\pi$ (homothétie ou translation).
\item La composée de deux rotations d'angles $\alpha$ et $\beta$ est une rotation d'angle $\alpha+\beta$ \textbf{si} $\alpha+\beta \not\equiv 0 \pmod{2\pi}$ ; sinon, c'est une translation.
\item Une similitude directe non identité a au maximum un point fixe (son centre). Fixer deux points distincts force $a=1$ et $b=0$.
\end{itemize}
\end{attention}

\begin{exercice}[Lecture d'éléments caractéristiques]
Pour chacune des similitudes directes suivantes, donner le rapport $k$, l'angle $\theta$ (principal, dans $]-\pi,\pi]$) et l'affixe $\omega$ du centre $\Omega$.
\begin{enumerate}
\item $z' = 2i\,z + 3 - i$ ;
\item $z' = \left(\dfrac{1}{2} - \dfrac{\sqrt{3}}{2}i\right) z + 2$ ;
\item $z' = -3z + 8 + 4i$.
\end{enumerate}
\end{exercice}

\begin{methode}[Déterminer le centre d'une similitude directe]
Donnée $z' = az + b$ avec $a \neq 1$ :
\begin{enumerate}
\item Écrire $a$ sous forme exponentielle : $a = k e^{i\theta}$ ($k>0$). Le rapport est $k$, l'angle est $\theta$.
\item Le centre $\Omega$ est l'unique point fixe : $\omega = a\omega + b \iff \omega(1-a) = b \iff \omega = \dfrac{b}{1-a}$.
\item Calculer $\omega$ en rationalisant le dénominateur (multiplier par le conjugué de $1-a$).
\end{enumerate}
Si $a=1$, c'est une translation (pas de centre, ou centre à l'infini).
\end{methode}

\begin{exercice}[Calculs directs d'images]
Le plan complexe est muni d'un repère orthonormé direct $(O\,;\vec{u},\vec{v})$. On considère la similitude directe $s$ d'écriture complexe $z' = (1 + i)z - 2$.
\begin{enumerate}
\item Déterminer l'affixe de l'image par $s$ du point $A$ d'affixe $2 - i$.
\item Déterminer l'affixe de l'antécédent par $s$ du point $B'$ d'affixe $3 + i$.
\item Quelle est l'image du point $O$ ? Le point $O$ est-il un point fixe de $s$ ?
\end{enumerate}
\end{exercice}

\courssub{Je m'entraîne}

\begin{exercice}[Nature et éléments caractéristiques]
On considère la transformation $f$ du plan complexe d'écriture complexe :
\[ z' = (1 + i\sqrt{3})\,z + 2 - 2i\sqrt{3}. \]
\begin{enumerate}
\item Justifier que $f$ est une similitude directe et déterminer son rapport et son angle.
\item Déterminer l'affixe $\omega$ de son centre $\Omega$.
\item Montrer que $f$ est la composée d'une homothétie et d'une rotation de même centre que l'on précisera.
\end{enumerate}
\end{exercice}

\begin{exercice}[De l'écriture complexe à l'expression analytique]
Soit $s$ la similitude directe de centre $\Omega(1 - i)$, de rapport $2$ et d'angle $-\dfrac{\pi}{4}$.
\begin{enumerate}
\item Déterminer l'écriture complexe de $s$.
\item En déduire son expression analytique, c'est-à-dire les coordonnées $(x'\,; y')$ du point $M' = s(M)$ en fonction des coordonnées $(x\,; y)$ de $M$.
\item Calculer les coordonnées de l'image du point $A(2\,; 1)$ par $s$.
\end{enumerate}
\end{exercice}

\begin{exercice}[Image d'une droite par deux méthodes]
On considère la similitude directe $s$ d'écriture complexe $z' = 2i\,z - 1 + i$ et la droite $(d)$ d'équation $y = x + 1$.
\begin{enumerate}
\item \textbf{Méthode analytique.} Déterminer l'expression analytique de $s$, puis une équation cartésienne de la droite $(d')$, image de $(d)$ par $s$.
\item \textbf{Méthode géométrique.} Choisir deux points $A$ et $B$ de $(d)$, calculer les affixes de leurs images $A'$ et $B'$, puis en déduire une équation de $(d')$. Vérifier la cohérence avec la question 1.
\item Quel est l'angle entre les droites $(d)$ et $(d')$ ? Retrouver ce résultat à partir de l'angle de $s$.
\end{enumerate}
\end{exercice}

\begin{methode}[Image d'une droite par une similitude directe]
Deux méthodes équivalentes :
\begin{enumerate}
\item \textbf{Analytique :} Passer par l'expression analytique $(x',y')$ en fonction de $(x,y)$, substituer dans l'équation de la droite, éliminer $x,y$.
\item \textbf{Géométrique :} L'image d'une droite est une droite. Il suffit de calculer les images de \textbf{deux points distincts} de la droite (ou d'un point et du vecteur directeur). La droite image passe par ces deux images.
\end{enumerate}
\cle{Propriété} : Une similitude directe d'angle $\theta$ transforme toute droite en une droite qui forme un angle $\theta$ avec elle (modulo $\pi$).
\end{methode}

\begin{exercice}[Triangle rectangle isocèle et rotation]
Dans le plan complexe, on donne les points $A$, $B$ et $C$ d'affixes respectives :
\[ z_A = 1 + i, \qquad z_B = 4 + 2i, \qquad z_C = 4i. \]
\begin{enumerate}
\item Calculer $\dfrac{z_C - z_A}{z_B - z_A}$ et donner son module et son argument principal.
\item En déduire que le triangle $ABC$ est rectangle isocèle en $A$.
\item Déterminer l'écriture complexe de la rotation $r$ de centre $A$ qui transforme $B$ en $C$.
\item Le point $D$ est l'image de $C$ par $r$. Déterminer l'affixe de $D$ et la nature du quadrilatère $ABDC$.
\end{enumerate}
\end{exercice}

\begin{savaistu}[Le repérage complexe au Cameroun]
Les coordonnées GPS (longitude, latitude) utilisées par les applications de navigation (Google Maps, Waze) ou pour le suivi des flottes de transport (Camrail, bus urbains de Douala/Yaoundé) sont un repérage sphérique. En cartographie locale (cadastre, bornage, travaux publics), on utilise un repère plan projectif (Lambert, UTM) où chaque point a des coordonnées $(X,Y)$. L'identification $(X,Y) \leftrightarrow X+iY$ permet d'utiliser toute la puissance du calcul complexe pour les rotations (changement d'orientation d'un chantier), homothéties (changement d'échelle de plan) et similitudes (superposition de plans cadastraux).
\end{savaistu}

\courssub{Je me prépare au Bac}

\begin{exercice}[Type Bac : étude complète d'une similitude et image d'un cercle]
Le plan complexe est rapporté à un repère orthonormé direct $(O\,;\vec{u},\vec{v})$, unité graphique $\SI{2}{cm}$. On considère la transformation $s$ du plan qui, à tout point $M$ d'affixe $z$, associe le point $M'$ d'affixe :
\[ z' = (1 + i\sqrt{3})\,z + 2 - 2i\sqrt{3}. \]
\begin{enumerate}
\item \begin{enumerate}
\item Écrire $1 + i\sqrt{3}$ sous forme exponentielle.
\item En déduire la nature exacte de $s$ ainsi que son rapport et son angle.
\item Déterminer l'affixe $\omega$ du centre $\Omega$ de $s$.
\end{enumerate}
\item On pose $A$ le point d'affixe $z_A = 1$ et $B$ le point d'affixe $z_B = 2 + i\sqrt{3}$.
\begin{enumerate}
\item Calculer les affixes $z_{A'}$ et $z_{B'}$ des images $A'$ et $B'$ de $A$ et $B$ par $s$.
\item Le point $B$ appartient-il au cercle $(\mathcal{C})$ de centre $O$ et de rayon $1$ ? Sinon, placer $B$ par rapport à ce cercle en calculant $OB$.
\item Déterminer l'image $(\mathcal{C}')$ du cercle $(\mathcal{C})$ de centre $O$ et de rayon $1$ par $s$ : on précisera son centre et son rayon.
\end{enumerate}
\item Soit $M$ un point d'affixe $z$ et $M'$ son image d'affixe $z'$.
\begin{enumerate}
\item Démontrer que $|z' - \omega| = 2\,|z - \omega|$.
\item En déduire que si $M$ décrit le cercle de centre $\Omega$ et de rayon $1$, alors $M'$ décrit un cercle que l'on précisera.
\end{enumerate}
\item Un technicien de CAMTEL modélise l'agrandissement d'une zone de couverture réseau autour d'un relais placé en $\Omega$ par la transformation $s$. La zone initiale est le disque de centre $O$ et de rayon $1$ (unité : $\SI{1}{km}$). Quelle est l'aire, en $\mathrm{km}^2$, de la zone après transformation ?
\end{enumerate}
\end{exercice}

\begin{exercice}[Type Bac : détermination d'une similitude définie par deux couples de points]
Le plan complexe est muni d'un repère orthonormé direct $(O\,;\vec{u},\vec{v})$. Une entreprise de Bafoussam matérialise quatre points du terrain par leurs affixes :
\[ z_A = -1 + 2i, \qquad z_B = 1 + 4i, \qquad z_C = 3 + i, \qquad z_D = 5 - i. \]
On cherche la similitude directe $s$ telle que $s(A) = C$ et $s(B) = D$.
\begin{enumerate}
\item On note $z' = az + b$ l'écriture complexe de $s$, où $a \in \mathbb{C}^*$ et $b \in \mathbb{C}$.
\begin{enumerate}
\item Montrer que $a$ et $b$ vérifient le système :
$\begin{cases} a(-1 + 2i) + b = 3 + i \\ a(1 + 4i) + b = 5 - i \end{cases}$
\item Résoudre ce système et en déduire l'écriture complexe de $s$.
\end{enumerate}
\item \begin{enumerate}
\item Déterminer le rapport et l'angle de $s$.
\item Calculer l'affixe $\omega$ du centre $\Omega$ de $s$.
\item Donner l'expression analytique de $s$.
\end{enumerate}
\item \begin{enumerate}
\item Calculer les longueurs $AB$ et $CD$, et vérifier que $CD = k \cdot AB$ où $k$ est le rapport de $s$.
\item Déterminer une mesure de l'angle $\left(\overrightarrow{AB}\,;\overrightarrow{CD}\right)$ et la comparer à l'angle de $s$. Conclure.
\end{enumerate}
\item Le point $E$ est le milieu du segment $[AB]$. Déterminer l'affixe de $E' = s(E)$ et vérifier que $E'$ est le milieu de $[CD]$. Quelle propriété des similitudes cela illustre-t-il ?
\end{enumerate}
\end{exercice}

\begin{exercice}[Type Bac : suite de points et configuration géométrique]
Le plan complexe est rapporté à un repère orthonormé direct $(O\,;\vec{u},\vec{v})$. On considère la similitude directe $s$ d'écriture complexe :
\[ z' = \dfrac{1}{2}\,e^{i\frac{\pi}{3}}\,z + 3. \]
On définit la suite de points $(M_n)$ par $M_0 = O$ et, pour tout entier naturel $n$, $M_{n+1} = s(M_n)$. On note $z_n$ l'affixe de $M_n$.
\begin{enumerate}
\item \begin{enumerate}
\item Déterminer le rapport, l'angle et l'affixe $\omega$ du centre $\Omega$ de $s$.
\item Calculer $z_1$ et $z_2$.
\end{enumerate}
\item \begin{enumerate}
\item Démontrer que pour tout entier naturel $n$ : $z_{n+1} - \omega = \dfrac{1}{2}e^{i\frac{\pi}{3}}\,(z_n - \omega)$.
\item On pose $u_n = z_n - \omega$. Montrer que la suite $(u_n)$ est géométrique ; préciser sa raison et son premier terme.
\item En déduire l'expression de $z_n$ en fonction de $n$, puis la limite de la suite $(z_n)$. Interpréter géométriquement ce résultat.
\end{enumerate}
\item \begin{enumerate}
\item Étudier la nature du triangle $\Omega M_n M_{n+1}$ pour tout entier naturel $n$.
\item Calculer la longueur $\Omega M_n$ en fonction de $n$ et en déduire à partir de quel rang $n$ on a $\Omega M_n < 10^{-3}$.
\end{enumerate}
\item \textbf{Application.} Un opérateur télécoms ajuste itérativement la position d'une antenne relais à Yaoundé : chaque nouvelle position est l'image de la précédente par $s$ (unité : $\SI{100}{m}$). Vers quel point l'antenne se stabilise-t-elle ? Justifier.
\end{enumerate}
\end{exercice}

\begin{popfigure}
\begin{tikzpicture}[scale=0.8, >=Stealth]
% Axes
\draw[->, popmuted] (-1,0) -- (5,0) node[right] {$x$};
\draw[->, popmuted] (0,-1) -- (0,4) node[above] {$y$};
% Center Omega
\coordinate (O) at (3,1.732); % 3 + i*sqrt(3)
\fill[popblue] (O) circle (2pt) node[above right] {$\Omega$};
% Spiral points
\foreach \n/\r/\ang in {0/3.464/0, 1/1.732/60, 2/0.866/120, 3/0.433/180, 4/0.216/240} {
    \pgfmathsetmacro{\x}{3 + \r*cos(\ang)}
    \pgfmathsetmacro{\y}{1.732 + \r*sin(\ang)}
    \coordinate (M\n) at (\x,\y);
    \fill[poporange] (M\n) circle (1.5pt);
    \node[poporange, font=\tiny] at (\x,\y) {$M_{\n}$};
}
% Segments
\draw[poporange, thick] (M0) -- (M1) -- (M2) -- (M3) -- (M4);
% Circle initial radius
\draw[popblue, dashed] (O) circle (3.464);
\draw[popblue, dashed] (O) circle (1.732);
\draw[popblue, dashed] (O) circle (0.866);
% Vectors for triangle n=0
\draw[popgreen, thick, ->] (O) -- (M0) node[midway, left] {$\Omega M_0$};
\draw[popgreen, thick, ->] (O) -- (M1) node[midway, right] {$\Omega M_1$};
\draw[popgreen, thick, ->] (M1) -- (M0) node[midway, below] {$M_1 M_0$};
% Right angle mark
\draw[popgreen] (M1) ++(-0.2,0) -- ++(0,-0.2) -- ++(0.2,0);
\end{tikzpicture}
\legende{Convergence spiralaire de la suite $(M_n)$ vers le centre $\Omega$ (Exercice 3, Bac).}
\end{popfigure}

\begin{aretenir}[Bilan : Similitude directe $z' = az + b$ ($a \neq 0$)]
\begin{itemize}
\item \textbf{Nature} : Similitude directe. Cas particuliers : translation ($a=1$), rotation ($|a|=1, a\neq 1$), homothétie ($a \in \mathbb{R}^*_+, a\neq 1$).
\item \textbf{Rapport} : $k = |a|$. \textbf{Angle} : $\theta = \arg(a)$ (principal dans $]-\pi,\pi]$).
\item \textbf{Centre} $\Omega(\omega)$ : Point fixe unique si $a \neq 1$. $\omega = \dfrac{b}{1-a}$.
\item \textbf{Expression analytique} : Si $a = k(\cos\theta + i\sin\theta)$, $b = b_x + i b_y$, $\omega = \omega_x + i\omega_y$ :
\[ \begin{cases} x' = \omega_x + k(x\cos\theta - y\sin\theta) - k(\omega_x\cos\theta - \omega_y\sin\theta) \\ y' = \omega_y + k(x\sin\theta + y\cos\theta) - k(\omega_x\sin\theta + \omega_y\cos\theta) \end{cases} \]
Ou plus simplement : $\begin{pmatrix} x'-\omega_x \\ y'-\omega_y \end{pmatrix} = k \begin{pmatrix} \cos\theta & -\sin\theta \\ \sin\theta & \cos\theta \end{pmatrix} \begin{pmatrix} x-\omega_x \\ y-\omega_y \end{pmatrix}$.
\item \textbf{Conservation} : Angles (orientés), parallélisme, alignement, rapports de longueurs ($k$), milieux, barycentres.
\item \textbf{Images} : Droite $\to$ Droite (angle $\theta$). Cercle $\to$ Cercle (rayon multiplié par $k$, centre image du centre).
\end{itemize}
\end{aretenir}

\newpage

\coursec{Corrigés des exercices}

\begin{corrige}[Exercice 1 — QCM : nature d'une transformation]
On utilise le critère fondamental : pour $z' = az + b$ avec $a \in \mathbb{C}^*$, on examine $a$.
\begin{enumerate}
\item $a = 3 \in \mathbb{R}^*_+$, $a \neq 1$ : c'est une \cle{homothétie} de rapport $3$, de centre d'affixe $\omega = \dfrac{b}{1-a} = \dfrac{-2+i}{1-3} = \dfrac{-2+i}{-2} = 1 - \dfrac{i}{2}$.
\item $a = 1$ : c'est une \cle{translation} de vecteur d'affixe $b = 4 - 2i$.
\item $a = -i$, $|a| = 1$ et $a \neq 1$ : c'est une \cle{rotation} d'angle $\arg(-i) = -\dfrac{\pi}{2}$.
\item $a = 1+i$, $|a| = \sqrt{2} \neq 1$ et $a \notin \mathbb{R}$ : c'est une \cle{similitude directe} non réduite aux cas précédents (rapport $\sqrt{2}$, angle $\dfrac{\pi}{4}$).
\end{enumerate}
\end{corrige}

\begin{corrige}[Exercice 2 — Vrai ou faux, en justifiant]
\begin{enumerate}
\item \textbf{Vrai.} Une homothétie de rapport $k > 0$ s'écrit $z' - \omega = k(z - \omega)$, soit $z' = az + b$ avec $a = k \in \mathbb{R}^*_+$. C'est bien une similitude directe (de rapport $k$ et d'angle $0$).
\item \textbf{Faux.} L'écriture fait intervenir $\bar{z}$ : c'est une similitude \emph{indirecte} (elle renverse les angles). Une similitude directe a une écriture de la forme $z' = az + b$, sans conjugaison.
\item \textbf{Faux.} L'image d'une droite $(d)$ par une similitude directe d'angle $\theta$ est une droite $(d')$ faisant un angle $\theta$ avec $(d)$. Contre-exemple : la rotation d'angle $\dfrac{\pi}{2}$ envoie l'axe des abscisses sur l'axe des ordonnées, qui lui est perpendiculaire, non parallèle. Il n'y a parallélisme que si $\theta \equiv 0 \pmod{\pi}$.
\item \textbf{Faux.} Si $r_1$ et $r_2$ sont de centres $\Omega_1, \Omega_2$ et d'angles $\alpha, \beta$ avec $\alpha + \beta \equiv 0 \pmod{2\pi}$ (et $\Omega_1 \neq \Omega_2$), alors $r_2 \circ r_1$ est une \emph{translation}. Exemple : la composée de deux symétries centrales de centres distincts est une translation.
\item \textbf{Vrai.} Si $M_1(z_1)$ et $M_2(z_2)$, $z_1 \neq z_2$, sont fixes, alors $z_1 = az_1 + b$ et $z_2 = az_2 + b$. Par soustraction : $z_1 - z_2 = a(z_1 - z_2)$, d'où $a = 1$ puisque $z_1 \neq z_2$. Alors $z_1 = z_1 + b$ donne $b = 0$ : c'est l'identité.
\end{enumerate}
\end{corrige}

\begin{corrige}[Exercice 3 — Lecture d'éléments caractéristiques]
On applique : $k = |a|$, $\theta = \arg(a)$, $\omega = \dfrac{b}{1-a}$.
\begin{enumerate}
\item $z' = 2i\,z + 3 - i$ : $a = 2i = 2e^{i\frac{\pi}{2}}$, donc $k = 2$ et $\theta = \dfrac{\pi}{2}$.
\[ \omega = \frac{3 - i}{1 - 2i} = \frac{(3-i)(1+2i)}{(1-2i)(1+2i)} = \frac{3 + 6i - i - 2i^2}{1+4} = \frac{5 + 5i}{5} = 1 + i. \]
\item $a = \dfrac{1}{2} - \dfrac{\sqrt{3}}{2}i = \cos\left(-\dfrac{\pi}{3}\right) + i\sin\left(-\dfrac{\pi}{3}\right) = e^{-i\frac{\pi}{3}}$ : $k = 1$, $\theta = -\dfrac{\pi}{3}$ (c'est une rotation).
\[ \omega = \frac{2}{1 - a} = \frac{2}{\dfrac{1}{2} + \dfrac{\sqrt{3}}{2}i} = \frac{2\left(\dfrac{1}{2} - \dfrac{\sqrt{3}}{2}i\right)}{\dfrac{1}{4} + \dfrac{3}{4}} = 1 - \sqrt{3}\,i. \]
\item $a = -3$ : $k = |-3| = 3$, $\theta = \arg(-3) = \pi$ (homothétie de rapport $-3$).
\[ \omega = \frac{8 + 4i}{1 - (-3)} = \frac{8 + 4i}{4} = 2 + i. \]
\end{enumerate}
\end{corrige}

\begin{corrige}[Exercice 4 — Calculs directs d'images]
\begin{enumerate}
\item $z_{A'} = (1+i)(2-i) - 2 = 2 - i + 2i - i^2 - 2 = 2 + i + 1 - 2 = 1 + i$. Donc $A'$ a pour affixe $1 + i$.
\item On résout $3 + i = (1+i)z - 2$, soit $z = \dfrac{5 + i}{1+i} = \dfrac{(5+i)(1-i)}{(1+i)(1-i)} = \dfrac{5 - 5i + i - i^2}{2} = \dfrac{6 - 4i}{2} = 3 - 2i$. L'antécédent de $B'$ est le point d'affixe $3 - 2i$.
\item $z_{O'} = (1+i)\times 0 - 2 = -2$. L'image de $O$ est le point d'affixe $-2$. Comme $s(O) \neq O$, le point $O$ n'est \textbf{pas} un point fixe de $s$. (Le seul point fixe est le centre $\Omega$ d'affixe $\omega = \dfrac{-2}{1-(1+i)} = \dfrac{-2}{-i} = \dfrac{2}{i} = -2i$.)
\end{enumerate}
\end{corrige}

\begin{corrige}[Exercice 5 — Nature et éléments caractéristiques]
\begin{enumerate}
\item L'écriture $z' = (1+i\sqrt{3})z + 2 - 2i\sqrt{3}$ est de la forme $z' = az + b$ avec $a = 1 + i\sqrt{3} \neq 0$ : $f$ est une \cle{similitude directe}. Or $a = 2\left(\dfrac{1}{2} + i\dfrac{\sqrt{3}}{2}\right) = 2e^{i\frac{\pi}{3}}$. Donc le rapport est $k = |a| = 2$ et l'angle est $\theta = \arg(a) = \dfrac{\pi}{3}$.
\item Le centre est l'unique point fixe : $\omega = \dfrac{b}{1-a} = \dfrac{2 - 2i\sqrt{3}}{-i\sqrt{3}}$. On multiplie numérateur et dénominateur par $i\sqrt{3}$ :
\[ \omega = \frac{(2 - 2i\sqrt{3})(i\sqrt{3})}{(-i\sqrt{3})(i\sqrt{3})} = \frac{2i\sqrt{3} - 2i^2 \times 3}{3} = \frac{6 + 2i\sqrt{3}}{3} = 2 + \frac{2\sqrt{3}}{3}i. \]
\item Toute similitude directe de centre $\Omega$, rapport $k$ et angle $\theta$ est la composée $f = h \circ r = r \circ h$ de l'homothétie $h$ de centre $\Omega$ et de rapport $2$ et de la rotation $r$ de centre $\Omega$ et d'angle $\dfrac{\pi}{3}$ (les deux commutent car elles ont même centre). En effet, pour tout $M$ : $z' - \omega = 2e^{i\frac{\pi}{3}}(z - \omega)$, ce qui se lit « on tourne de $\dfrac{\pi}{3}$ autour de $\Omega$ puis on dilate d'un facteur $2$ » (ou l'inverse).
\end{enumerate}
\end{corrige}

\begin{corrige}[Exercice 6 — De l'écriture complexe à l'expression analytique]
\begin{enumerate}
\item On a $a = k e^{i\theta} = 2e^{-i\frac{\pi}{4}} = 2\left(\dfrac{\sqrt{2}}{2} - i\dfrac{\sqrt{2}}{2}\right) = \sqrt{2} - i\sqrt{2}$. L'écriture réduite $z' - \omega = a(z - \omega)$ avec $\omega = 1 - i$ donne $b = \omega(1 - a)$ :
\begin{align*}
b &= (1 - i)\left(1 - \sqrt{2} + i\sqrt{2}\right) = \left(1 - \sqrt{2} + i\sqrt{2}\right) - i\left(1 - \sqrt{2} + i\sqrt{2}\right) \\
&= 1 - \sqrt{2} + i\sqrt{2} - i + i\sqrt{2} + \sqrt{2} = 1 + i\left(2\sqrt{2} - 1\right).
\end{align*}
D'où l'écriture complexe : $z' = \left(\sqrt{2} - i\sqrt{2}\right)z + 1 + i\left(2\sqrt{2} - 1\right)$.
\item En posant $z = x + iy$ : $(\sqrt{2} - i\sqrt{2})(x + iy) = \sqrt{2}\,x + \sqrt{2}\,y + i\left(-\sqrt{2}\,x + \sqrt{2}\,y\right)$. Donc :
\[ \begin{cases} x' = \sqrt{2}\,x + \sqrt{2}\,y + 1 \\ y' = -\sqrt{2}\,x + \sqrt{2}\,y + 2\sqrt{2} - 1 \end{cases} \]
\item Pour $A(2\,;1)$ : $x' = 2\sqrt{2} + \sqrt{2} + 1 = 3\sqrt{2} + 1$ et $y' = -2\sqrt{2} + \sqrt{2} + 2\sqrt{2} - 1 = \sqrt{2} - 1$. Donc $A'\left(3\sqrt{2} + 1\,;\ \sqrt{2} - 1\right)$.
\end{enumerate}
\textbf{Vérification} : $\Omega$ est fixe : $x' = \sqrt{2} - \sqrt{2} + 1 = 1$ et $y' = -\sqrt{2} - \sqrt{2} + 2\sqrt{2} - 1 = -1$ : on retrouve bien $\Omega(1\,;-1)$.
\end{corrige}

\begin{corrige}[Exercice 7 — Image d'une droite par deux méthodes]
\begin{enumerate}
\item \textbf{Méthode analytique.} Avec $z = x + iy$ : $2i(x+iy) - 1 + i = -2y - 1 + i(2x + 1)$, donc $\begin{cases} x' = -2y - 1 \\ y' = 2x + 1 \end{cases}$. On inverse : $y = \dfrac{-x' - 1}{2}$ et $x = \dfrac{y' - 1}{2}$. En substituant dans $y = x + 1$ :
\[ \frac{-x'-1}{2} = \frac{y'-1}{2} + 1 \iff -x' - 1 = y' + 1 \iff y' = -x' - 2. \]
La droite image $(d')$ a pour équation $y = -x - 2$.
\item \textbf{Méthode géométrique.} Choisissons $A(0\,;1)$ et $B(1\,;2)$ sur $(d)$.
$A$ d'affixe $i$ : $z_{A'} = 2i \times i - 1 + i = -3 + i$, soit $A'(-3\,;1)$.
$B$ d'affixe $1 + 2i$ : $z_{B'} = 2i(1+2i) - 1 + i = 2i - 4 - 1 + i = -5 + 3i$, soit $B'(-5\,;3)$.
La droite $(A'B')$ a pour coefficient directeur $\dfrac{3-1}{-5-(-3)} = \dfrac{2}{-2} = -1$ et passe par $A'$ : $y - 1 = -(x + 3)$, soit $y = -x - 2$. \textbf{Cohérent} avec la question 1. (Vérification : $A'$ et $B'$ satisfont bien $y = -x-2$ : $3 - 2 = 1$ et $5 - 2 = 3$.)
\item $(d)$ a pour vecteur directeur $\vec{u}(1\,;1)$ (angle $\dfrac{\pi}{4}$ avec l'axe des abscisses) et $(d')$ a pour vecteur directeur $\vec{u}'(1\,;-1)$ (angle $-\dfrac{\pi}{4}$). L'angle entre $(d)$ et $(d')$ est donc $\dfrac{\pi}{4} - \left(-\dfrac{\pi}{4}\right) = \dfrac{\pi}{2}$ : les droites sont \textbf{perpendiculaires}. On retrouve ce résultat car l'angle de $s$ est $\theta = \arg(2i) = \dfrac{\pi}{2}$ : une similitude directe fait tourner toute droite de son angle.
\end{enumerate}
\end{corrige}

\begin{corrige}[Exercice 8 — Triangle rectangle isocèle et rotation]
\begin{enumerate}
\item $z_C - z_A = 4i - 1 - i = -1 + 3i$ et $z_B - z_A = 4 + 2i - 1 - i = 3 + i$. Donc :
\[ \frac{z_C - z_A}{z_B - z_A} = \frac{-1 + 3i}{3 + i} = \frac{(-1+3i)(3-i)}{(3+i)(3-i)} = \frac{-3 + i + 9i - 3i^2}{9 + 1} = \frac{10i}{10} = i. \]
Module : $|i| = 1$ ; argument principal : $\arg(i) = \dfrac{\pi}{2}$.
\item De $\dfrac{z_C - z_A}{z_B - z_A} = i$, on tire d'une part $\dfrac{AC}{AB} = \left|\dfrac{z_C - z_A}{z_B - z_A}\right| = 1$, donc $AC = AB$ (triangle isocèle en $A$), et d'autre part $\left(\overrightarrow{AB}\,;\overrightarrow{AC}\right) = \arg(i) = \dfrac{\pi}{2}$ (angle droit en $A$). Le triangle $ABC$ est donc \cle{rectangle isocèle en $A$}.
\item La rotation $r$ de centre $A$ qui envoie $B$ sur $C$ a pour angle $\dfrac{\pi}{2}$ (d'après la question 1). Son écriture complexe est $z' - z_A = e^{i\frac{\pi}{2}}(z - z_A)$, soit :
\[ z' = iz + (1+i)(1 - i) = iz + 2. \]
\item $z_D = i \times z_C + 2 = i \times 4i + 2 = -4 + 2 = -2$. Donc $D$ a pour affixe $-2$.
Remarquons que $z_D - z_A = -2 - 1 - i = -3 - i = -(z_B - z_A)$ : c'est normal, car $D = r(r(B))$ et $r \circ r$ est le demi-tour de centre $A$. Ainsi $A$ est le \textbf{milieu de} $[BD]$. De plus $AC = AB = AD = \sqrt{10}$ et $(AC) \perp (BD)$. Les points $B$, $A$, $D$ étant alignés, considérons le triangle $BCD$ : on a $BC = \left|z_C - z_B\right| = |-4 + 2i| = \sqrt{20}$, $CD = \left|z_D - z_C\right| = |-2 - 4i| = \sqrt{20}$ et $BD = \left|z_D - z_B\right| = |-6 - 2i| = \sqrt{40} = 2\sqrt{10}$, d'où $BC^2 + CD^2 = 20 + 20 = 40 = BD^2$ : le triangle $BCD$ est \cle{rectangle isocèle en $C$}, et $A$, milieu de son hypoténuse $[BD]$, est le centre du cercle circonscrit à ce triangle.
\end{enumerate}
\end{corrige}

\begin{corrige}[Exercice 9 — Type Bac : étude complète d'une similitude et image d'un cercle]
\textbf{1.a)} $1 + i\sqrt{3}$ : module $\sqrt{1+3} = 2$ ; $\cos\theta = \dfrac{1}{2}$, $\sin\theta = \dfrac{\sqrt{3}}{2}$ donc $\theta = \dfrac{\pi}{3}$. Ainsi $1 + i\sqrt{3} = 2e^{i\frac{\pi}{3}}$.

\textbf{1.b)} $s$ est de la forme $z' = az + b$ avec $a = 2e^{i\frac{\pi}{3}} \neq 0$ : c'est une \cle{similitude directe} de rapport $k = 2$ et d'angle $\theta = \dfrac{\pi}{3}$.

\textbf{1.c)} $\omega = \dfrac{b}{1-a} = \dfrac{2 - 2i\sqrt{3}}{-i\sqrt{3}} = \dfrac{(2-2i\sqrt{3})(i\sqrt{3})}{3} = \dfrac{6 + 2i\sqrt{3}}{3} = 2 + \dfrac{2\sqrt{3}}{3}i$.

\textbf{2.a)} $z_{A'} = (1+i\sqrt{3})(1) + 2 - 2i\sqrt{3} = 3 - i\sqrt{3}$.
$z_{B'} = (1+i\sqrt{3})(2+i\sqrt{3}) + 2 - 2i\sqrt{3} = 2 + 3i\sqrt{3} - 3 + 2 - 2i\sqrt{3} = 1 + i\sqrt{3}$.

\textbf{2.b)} $OB = |z_B| = |2 + i\sqrt{3}| = \sqrt{4+3} = \sqrt{7}$. Comme $\sqrt{7} > 1$, $B$ n'appartient pas à $(\mathcal{C})$ : il est \textbf{à l'extérieur} du cercle.

\textbf{2.c)} L'image d'un cercle par une similitude directe de rapport $k$ est un cercle de rayon multiplié par $k$, dont le centre est l'image du centre. Ici $s(O)$ a pour affixe $z' = 2 - 2i\sqrt{3}$. Donc $(\mathcal{C}')$ est le cercle de centre $O'\left(2 - 2i\sqrt{3}\right)$ et de rayon $k \times 1 = 2$.

\textbf{3.a)} Comme $\Omega$ est fixe, $\omega = a\omega + b$, donc $b = \omega - a\omega$. Alors :
\[ z' - \omega = az + b - \omega = az - a\omega = a(z - \omega). \]
En passant aux modules : $|z' - \omega| = |a|\,|z - \omega| = 2\,|z - \omega|$.

\textbf{3.b)} Si $M$ décrit le cercle de centre $\Omega$ et de rayon $1$, alors $|z - \omega| = 1$, donc $|z' - \omega| = 2$ : $M'$ décrit le \textbf{cercle de centre $\Omega$ et de rayon $2$}.

\textbf{4)} Une similitude de rapport $k$ multiplie les aires par $k^2$. L'aire initiale est $\pi \times 1^2 = \pi\ \mathrm{km}^2$ ; l'aire après transformation est $4\pi\ \mathrm{km}^2 \approx \num{12,57}\ \mathrm{km}^2$.

\textbf{Barème indicatif (sur 10)} : 1.a) 1 pt ; 1.b) 1 pt ; 1.c) 1,5 pt ; 2.a) 1 pt ; 2.b) 1 pt ; 2.c) 1,5 pt ; 3.a) 1 pt ; 3.b) 0,5 pt ; 4) 1,5 pt.

\textbf{Points de vigilance} : justifier la nature en citant la forme $z' = az + b$, $a \neq 0$ ; pour l'image du cercle, préciser explicitement que le centre de l'image est \emph{l'image du centre} (et non $\Omega$, sauf si $O = \Omega$) ; en 3.a), la relation $z' - \omega = a(z-\omega)$ doit être démontrée à partir du point fixe, pas seulement citée ; en 4), citer la propriété « les aires sont multipliées par $k^2$ ».
\end{corrige}

\begin{corrige}[Exercice 10 — Type Bac : similitude définie par deux couples de points]
\textbf{1.a)} $s(A) = C$ s'écrit $a\,z_A + b = z_C$, c'est-à-dire $a(-1+2i) + b = 3 + i$ ; de même $s(B) = D$ donne $a(1+4i) + b = 5 - i$. D'où le système annoncé.

\textbf{1.b)} Par soustraction membre à membre : $a\left[(1+4i) - (-1+2i)\right] = (5-i) - (3+i)$, soit $a(2+2i) = 2 - 2i$, d'où :
\[ a = \frac{2-2i}{2+2i} = \frac{1-i}{1+i} = \frac{(1-i)^2}{(1+i)(1-i)} = \frac{1 - 2i + i^2}{2} = \frac{-2i}{2} = -i. \]
Puis $b = 3 + i - a(-1+2i) = 3 + i - (-i)(-1+2i) = 3 + i - (i + 2) = 1$.
Vérification avec la deuxième équation : $(-i)(1+4i) + 1 = -i - 4i^2 + 1 = 5 - i$ : c'est bien l'affixe de $D$. Donc $s$ a pour écriture complexe $z' = -iz + 1$.

\textbf{2.a)} $a = -i = e^{-i\frac{\pi}{2}}$ : $|a| = 1$ et $a \neq 1$, donc $s$ est une \cle{rotation} : rapport $k = 1$, angle $\theta = -\dfrac{\pi}{2}$.

\textbf{2.b)} $\omega = \dfrac{b}{1-a} = \dfrac{1}{1+i} = \dfrac{1-i}{(1+i)(1-i)} = \dfrac{1-i}{2}$.

\textbf{2.c)} Avec $z = x + iy$ : $-i(x+iy) + 1 = y + 1 - ix$. Donc $\begin{cases} x' = y + 1 \\ y' = -x \end{cases}$.

\textbf{3.a)} $AB = |z_B - z_A| = |2 + 2i| = \sqrt{4+4} = 2\sqrt{2}$ ; $CD = |z_D - z_C| = |2 - 2i| = 2\sqrt{2}$. On a bien $CD = k \cdot AB$ avec $k = 1$ (la rotation conserve les distances).

\textbf{3.b)} $\left(\overrightarrow{AB}\,;\overrightarrow{CD}\right) = \arg\left(\dfrac{z_D - z_C}{z_B - z_A}\right) = \arg\left(\dfrac{2-2i}{2+2i}\right) = \arg(-i) = -\dfrac{\pi}{2}$. Cet angle est égal à l'angle de $s$ : c'est conforme, car une similitude directe d'angle $\theta$ envoie tout vecteur $\overrightarrow{AB}$ sur un vecteur $\overrightarrow{A'B'}$ tel que $\left(\overrightarrow{AB}\,;\overrightarrow{A'B'}\right) = \theta$.

\textbf{4)} $z_E = \dfrac{z_A + z_B}{2} = \dfrac{(-1+2i) + (1+4i)}{2} = 3i$. Alors $z_{E'} = -i(3i) + 1 = 3 + 1 = 4$. Or le milieu de $[CD]$ a pour affixe $\dfrac{z_C + z_D}{2} = \dfrac{(3+i)+(5-i)}{2} = 4 = z_{E'}$. Donc $E'$ est bien le milieu de $[CD]$ : cela illustre la \cle{conservation des milieux} par les similitudes (plus généralement, conservation des barycentres).

\textbf{Barème indicatif (sur 10)} : 1.a) 1 pt ; 1.b) 2 pts ; 2.a) 1 pt ; 2.b) 1 pt ; 2.c) 1 pt ; 3.a) 1 pt ; 3.b) 1,5 pt ; 4) 1,5 pt.

\textbf{Points de vigilance} : ne pas oublier de vérifier la solution $(a,b)$ dans les deux équations ; reconnaître qu'ici la similitude est une rotation ($|a|=1$) et le dire explicitement ; pour l'angle entre vecteurs, utiliser l'argument d'un quotient d'affixes et non un calcul « au feeling » ; conclure la question 4 en nommant la propriété de conservation utilisée.
\end{corrige}

\begin{corrige}[Exercice 11 — Type Bac : suite de points et configuration géométrique]
\textbf{1.a)} $a = \dfrac{1}{2}e^{i\frac{\pi}{3}}$ : rapport $k = \dfrac{1}{2}$, angle $\theta = \dfrac{\pi}{3}$. Le centre :
\[ \omega = \frac{3}{1 - \frac{1}{2}e^{i\frac{\pi}{3}}} = \frac{3}{1 - \frac{1}{4} - i\frac{\sqrt{3}}{4}} = \frac{3}{\frac{3}{4} - i\frac{\sqrt{3}}{4}} = \frac{3\left(\frac{3}{4} + i\frac{\sqrt{3}}{4}\right)}{\frac{9}{16} + \frac{3}{16}} = \frac{3\left(\frac{3}{4} + i\frac{\sqrt{3}}{4}\right)}{\frac{3}{4}} = 3 + i\sqrt{3}. \]

\textbf{1.b)} $z_1 = \dfrac{1}{2}e^{i\frac{\pi}{3}} \times 0 + 3 = 3$. Puis :
\[ z_2 = \frac{1}{2}e^{i\frac{\pi}{3}} \times 3 + 3 = 3\left(\frac{1}{4} + i\frac{\sqrt{3}}{4}\right) + 3 = \frac{15}{4} + \frac{3\sqrt{3}}{4}i. \]

\textbf{2.a)} $\Omega$ étant fixe, $\omega = a\omega + 3$ avec $a = \dfrac{1}{2}e^{i\frac{\pi}{3}}$. Alors pour tout $n \in \mathbb{N}$ :
\[ z_{n+1} - \omega = a z_n + 3 - \omega = a z_n - a\omega = a\left(z_n - \omega\right) = \frac{1}{2}e^{i\frac{\pi}{3}}\left(z_n - \omega\right). \]

\textbf{2.b)} Posant $u_n = z_n - \omega$, la relation précédente s'écrit $u_{n+1} = \dfrac{1}{2}e^{i\frac{\pi}{3}}\,u_n$ : $(u_n)$ est \cle{géométrique} de raison $q = \dfrac{1}{2}e^{i\frac{\pi}{3}}$ et de premier terme $u_0 = z_0 - \omega = -3 - i\sqrt{3}$.

\textbf{2.c)} Pour tout $n$ : $u_n = q^n u_0$, donc $z_n = \omega + \left(\dfrac{1}{2}\right)^n e^{i\frac{n\pi}{3}}\left(-3 - i\sqrt{3}\right)$, c'est-à-dire :
\[ z_n = 3 + i\sqrt{3} - \left(\frac{1}{2}\right)^n e^{i\frac{n\pi}{3}}\left(3 + i\sqrt{3}\right). \]
Comme $|q| = \dfrac{1}{2} < 1$, $q^n \to 0$, donc $\lim\limits_{n\to+\infty} z_n = \omega = 3 + i\sqrt{3}$. \textbf{Interprétation} : les points $M_n$ se rapprochent indéfiniment du centre $\Omega$ en spiralant autour de lui (à chaque étape, la distance à $\Omega$ est divisée par $2$ et le point tourne de $\dfrac{\pi}{3}$) : c'est la \emph{convergence spiralaire} illustrée par la figure.

\textbf{3.a)} Comme $M_{n+1} = s(M_n)$ et $\Omega$ fixe : $\Omega M_{n+1} = \dfrac{1}{2}\Omega M_n$ et $\left(\overrightarrow{\Omega M_n}\,;\overrightarrow{\Omega M_{n+1}}\right) = \dfrac{\pi}{3}$. Posons $L = \Omega M_n$. Par la formule d'Al-Kashi dans le triangle $\Omega M_n M_{n+1}$ :
\[ M_nM_{n+1}^2 = L^2 + \frac{L^2}{4} - 2 \times L \times \frac{L}{2}\cos\frac{\pi}{3} = \frac{5L^2}{4} - \frac{L^2}{2} = \frac{3L^2}{4}. \]
Alors $\Omega M_{n+1}^2 + M_nM_{n+1}^2 = \dfrac{L^2}{4} + \dfrac{3L^2}{4} = L^2 = \Omega M_n^2$ : d'après la réciproque du théorème de Pythagore, le triangle $\Omega M_n M_{n+1}$ est \cle{rectangle en $M_{n+1}$}, avec un angle de $\dfrac{\pi}{3}$ en $\Omega$ (c'est un « demi-triangle équilatéral »).

\textbf{3.b)} $\Omega M_n = |u_n| = |q|^n\,|u_0| = \left(\dfrac{1}{2}\right)^n \times \sqrt{9+3} = \dfrac{2\sqrt{3}}{2^n}$. On cherche $n$ tel que $\dfrac{2\sqrt{3}}{2^n} < 10^{-3}$, c'est-à-dire $2^n > 2000\sqrt{3} \approx \num{3464,1}$. Or $2^{11} = 2048 < \num{3464,1}$ et $2^{12} = 4096 > \num{3464,1}$. Donc $\Omega M_n < 10^{-3}$ à partir du rang $n = 12$.

\textbf{4)} D'après 2.c), la suite des positions converge vers le point fixe $\Omega$ d'affixe $3 + i\sqrt{3}$, c'est-à-dire le point de coordonnées $\left(3\,;\sqrt{3}\right)$ dans le repère (soit environ $(\SI{300}{m}\,;\SI{173}{m})$ en grandeur réelle). L'antenne se \textbf{stabilise en $\Omega$}, le centre de la similitude : c'est l'unique point fixe, attracteur car le rapport $\dfrac{1}{2}$ est strictement inférieur à $1$.

\textbf{Barème indicatif (sur 10)} : 1.a) 2 pts ; 1.b) 1 pt ; 2.a) 1 pt ; 2.b) 1 pt ; 2.c) 2 pts ; 3.a) 1,5 pt ; 3.b) 1 pt ; 4) 0,5 pt.

\textbf{Points de vigilance} : la relation de récurrence $z_{n+1} - \omega = a(z_n - \omega)$ doit être \emph{démontrée} en utilisant le point fixe ; pour la limite, justifier par $|q| < 1$ ; pour la nature du triangle, une démonstration par le calcul (Pythagore ou produit scalaire) est exigée, pas un constat sur figure ; pour le rang, résoudre proprement l'inéquation (les logarithmes ne sont pas indispensables : un encadrement par puissances de $2$ suffit).
\end{corrige}

\newpage

\coursec{Fiche bilan}

\coursec{Similitudes directes du plan}

\begin{popfigure}
\begin{tikzpicture}[mindmap, grow cyclic,
  every node/.style={concept, text=white, font=\small\bfseries},
  root concept/.append style={concept color=popdark, font=\large\bfseries},
  level 1/.append style={level distance=3.8cm, sibling angle=60},
  level 2/.append style={level distance=2.4cm, sibling angle=45, font=\scriptsize\bfseries}]
\node[root concept]{Similitudes directes}
  child[concept color=popblue]{ node {Rappels transformations}
    child { node {Translation : $z'=z+b$} }
    child { node {Homoth\'etie : $z'-\omega=k(z-\omega)$} }
    child { node {Rotation : $z'-\omega=e^{i\theta}(z-\omega)$} }
  }
  child[concept color=poporange]{ node {D\'efinition}
    child { node {$z'=az+b$, $a\in\mathbb{C}^*$} }
    child { node {Compos\'ee $h_{\Omega,k}\circ r_{\Omega,\theta}$} }
  }
  child[concept color=popgreen]{ node {\'El\'ements caract\'eristiques}
    child { node {Centre $\Omega$ : $\omega=\dfrac{b}{1-a}\ (a\neq 1)$} }
    child { node {Rapport $k=|a|>0$} }
    child { node {Angle $\theta=\arg(a)\ [2\pi]$} }
  }
  child[concept color=poppurple]{ node {Forme r\'eduite}
    child { node {$z'-\omega=a(z-\omega)$} }
    child { node {$a=ke^{i\theta}$} }
  }
  child[concept color=poppink]{ node {Propri\'et\'es}
    child { node {Conserve angles orient\'es} }
    child { node {Conserve alignement, parall\'elisme} }
    child { node {Longueurs $\times k$, Aires $\times k^2$} }
  }
  child[concept color=popgold]{ node {Images \& Applications}
    child { node {Droite $\to$ droite} }
    child { node {Cercle $\to$ cercle (centre image, rayon $\times k$)} }
    child { node {Triangle $\to$ triangle semblable direct} }
    child { node {Probl\`emes configuration} }
  };
\end{tikzpicture}
\legende{Carte mentale de la le\c{c}on : les six id\'ees forces des similitudes directes.}
\end{popfigure}

\courssub{Rappels : \'ecritures complexes des transformations usuelles}

\begin{definition}[Translation]
Soit $\vec{v}$ un vecteur du plan d'affixe $b \in \mathbb{C}^*$. La translation de vecteur $\vec{v}$ est l'application $T_{\vec{v}}$ qui \`a tout point $M(z)$ associe $M'(z')$ tel que :
\[
z' = z + b.
\]
En coordonn\'ees : si $b = u+iv$, alors $\begin{cases} x' = x + u \\ y' = y + v \end{cases}$.
\end{definition}

\begin{definition}[Homoth\'etie]
Soit $\Omega(\omega)$ un point du plan et $k \in \mathbb{R}^*$ un nombre r\'eel non nul. L'homoth\'etie de centre $\Omega$ et de rapport $k$ est l'application $H_{\Omega,k}$ qui \`a tout point $M(z)$ associe $M'(z')$ tel que :
\[
z' - \omega = k(z - \omega) \quad \iff \quad z' = kz + (1-k)\omega.
\]
En coordonn\'ees : $\begin{cases} x' = kx + (1-k)x_\Omega \\ y' = ky + (1-k)y_\Omega \end{cases}$.
\end{definition}

\begin{definition}[Rotation]
Soit $\Omega(\omega)$ un point du plan et $\theta \in \mathbb{R}$ un angle (orient\'e modulo $2\pi$). La rotation de centre $\Omega$ et d'angle $\theta$ est l'application $R_{\Omega,\theta}$ qui \`a tout point $M(z)$ associe $M'(z')$ tel que :
\[
z' - \omega = e^{i\theta}(z - \omega) \quad \iff \quad z' = e^{i\theta}z + \omega(1-e^{i\theta}).
\]
En coordonn\'ees : si $e^{i\theta} = \cos\theta + i\sin\theta$, alors
\[
\begin{cases}
x' = x_\Omega + (x-x_\Omega)\cos\theta - (y-y_\Omega)\sin\theta \\
y' = y_\Omega + (x-x_\Omega)\sin\theta + (y-y_\Omega)\cos\theta
\end{cases}
\]
\end{definition}

\begin{propriete}[Composition homoth\'etie - rotation]
Soient $H_{\Omega,k}$ une homoth\'etie ($k>0$) et $R_{\Omega,\theta}$ une rotation de m\^eme centre $\Omega$. Leur compos\'ee (dans n'importe quel ordre car elles commutent) est une transformation d'\'ecriture complexe :
\[
z' - \omega = k e^{i\theta}(z - \omega).
\]
C'est la forme \emph{r\'eduite} d'une similitude directe de centre $\Omega$, rapport $k$ et angle $\theta$.
\end{propriete}

\begin{exemplebox}[Exemple num\'erique : transformations \'el\'ementaires]
Soit le point $A(1+i)$.
\begin{enumerate}
\item Translation de vecteur $\vec{v}(2-i)$ : $z_A' = (1+i)+(2-i) = 3 \Rightarrow A'(3,0)$.
\item Homoth\'etie de centre $\Omega(0)$, rapport $k=2$ : $z_A' = 2(1+i) = 2+2i \Rightarrow A'(2,2)$.
\item Rotation de centre $\Omega(0)$, angle $\theta = \frac{\pi}{2}$ : $z_A' = e^{i\pi/2}(1+i) = i(1+i) = -1+i \Rightarrow A'(-1,1)$.
\end{enumerate}
\end{exemplebox}

\courssub{D\'efinition et \'ecriture complexe d'une similitude directe}

\begin{definition}[Similitude directe]
On appelle \cle{similitude directe} du plan toute application $S$ qui \`a tout point $M(z)$ associe un point $M'(z')$ d\'efini par une relation de la forme :
\[
\boxed{z' = a z + b}
\]
o\`u $a \in \mathbb{C}^*$ (c'est-\`a-dire $a \neq 0$) et $b \in \mathbb{C}$.
\end{definition}

\begin{propriete}[D\'ecomposition g\'eom\'etrique]
Toute similitude directe $S : z' = az+b$ avec $a \neq 1$ est la compos\'ee, dans n'importe quel ordre, d'une homoth\'etie $H_{\Omega,k}$ et d'une rotation $R_{\Omega,\theta}$ de m\^eme centre $\Omega$, o\`u :
\[
k = |a| > 0, \qquad \theta = \arg(a)\ [2\pi], \qquad \Omega(\omega) \text{ tel que } \omega = \frac{b}{1-a}.
\]
On \'ecrit $S = H_{\Omega,k} \circ R_{\Omega,\theta} = R_{\Omega,\theta} \circ H_{\Omega,k}$.
\end{propriete}

\begin{propriete}[Expression analytique]
Soit $S : z' = az+b$ avec $a = \alpha + i\beta$ et $b = u + iv$ ($\alpha,\beta,u,v \in \mathbb{R}$). L'expression analytique (cart\'esienne) de $S$ est :
\[
\boxed{\begin{cases}
x' = \alpha x - \beta y + u \\
y' = \beta x + \alpha y + v
\end{cases}}
\]
La matrice lin\'eaire associ\'ee est $\begin{pmatrix} \alpha & -\beta \\ \beta & \alpha \end{pmatrix} = |a| \begin{pmatrix} \cos\theta & -\sin\theta \\ \sin\theta & \cos\theta \end{pmatrix}$. C'est une matrice de similitude (produit d'une homoth\'etie et d'une rotation).
\end{propriete}

\begin{aretenir}[Cas particuliers (identification imm\'ediate)]
\begin{center}
\begin{tabularx}{\linewidth}{|l|X|X|}
\hline
\rowcolor{popblueL} \textbf{Condition sur $a$} & \textbf{Nature de $S$} & \textbf{Caract\'eristiques} \\
\hline
$a = 1$ & \cle{Translation} & Vecteur d'affixe $b$. Pas de point fixe. \\
\hline
$a \in \mathbb{R}^*$, $a \neq 1$ & \cle{Homoth\'etie} & Centre $\omega = \frac{b}{1-a}$, rapport $k=a$. Angle $0\ [2\pi]$. \\
\hline
$|a| = 1$, $a \neq 1$ & \cle{Rotation} & Centre $\omega = \frac{b}{1-a}$, angle $\theta = \arg(a)$. Rapport $k=1$. \\
\hline
$a \notin \mathbb{R}$, $|a| \neq 1$ & \cle{Similitude directe ``vraie''} & Centre $\omega = \frac{b}{1-a}$, rapport $k=|a|$, angle $\theta=\arg(a)$. \\
\hline
\end{tabularx}
\end{center}
\end{aretenir}

\begin{attention}[Pi\`ege classique : $a=1$]
La formule du centre $\omega = \frac{b}{1-a}$ n'est valide que si $a \neq 1$. Si $a=1$, c'est une translation : il n'y a \textbf{pas de centre} (aucun point fixe). Ne jamais \'ecrire $\omega = \frac{b}{0}$.
\end{attention}

\courssub{\'El\'ements caract\'eristiques : centre, rapport, angle}

\begin{propriete}[Centre d'une similitude directe ($a \neq 1$)]
Le centre $\Omega$ est l'unique point fixe de la transformation : $z = az+b \iff (1-a)z = b$. Son affixe est :
\[
\boxed{\omega = \frac{b}{1-a}}
\]
La forme r\'eduite (centr\'ee en $\Omega$) s'\'ecrit alors :
\[
\boxed{z' - \omega = a(z - \omega)}
\]
\end{propriete}

\begin{propriete}[Rapport et angle]
Pour une similitude directe $z' = az+b$ ($a \neq 0$), on d\'efinit :
\begin{itemize}
\item Le \cle{rapport} $k = |a| \in \mathbb{R}_+^*$. C'est le rapport des longueurs : $\forall M,N,\ M'N' = k \cdot MN$.
\item L'\cle{angle} $\theta = \arg(a)\ [2\pi]$. C'est l'angle orient\'e $(\overrightarrow{MN}, \overrightarrow{M'N'})$ pour tout couple de points distincts $M,N$.
\end{itemize}
On a donc la d\'ecomposition polaire fondamentale : $\boxed{a = k e^{i\theta}}$.
\end{propriete}

\begin{methode}[D\'eterminer la nature et les \'el\'ements de $z'=az+b$]
\begin{enumerate}
\item \textbf{Regarder $a$ :}
      \begin{itemize}
      \item Si $a=1$ : Translation de vecteur d'affixe $b$. \textbf{Arr\^eter l\`a.}
      \item Sinon ($a \neq 1$) : C'est une similitude directe de centre $\omega = \frac{b}{1-a}$.
      \end{itemize}
\item \textbf{Calculer le rapport :} $k = |a| = \sqrt{\alpha^2+\beta^2}$.
\item \textbf{Calculer l'angle :} $\theta = \arg(a)$ (rep\'er\'e dans $[0, 2\pi[$ ou $]-\pi,\pi]$). Utiliser $\cos\theta = \frac{\alpha}{k}$, $\sin\theta = \frac{\beta}{k}$.
\item \textbf{Conclure pr\'ecis\'ement :} ``$S$ est la similitude directe de centre $\Omega(\omega)$, rapport $k$, angle $\theta$.''
\end{enumerate}
\end{methode}

\begin{methode}[\'Ecrire l'\'equation complexe connaissant $\Omega$, $k$, $\theta$]
\begin{enumerate}
\item Calculer $a = k e^{i\theta} = k(\cos\theta + i\sin\theta)$.
\item Calculer $b = \omega(1-a)$.
\item \'Ecrire $z' = az + b$.
\end{enumerate}
\end{methode}

\begin{exemplebox}[Application : d\'etermination des \'el\'ements]
Soit $S$ d'\'equation $z' = (1+i\sqrt{3})z + 2-2i\sqrt{3}$.
\begin{enumerate}
\item $a = 1+i\sqrt{3} \neq 1$ : similitude directe (pas translation).
\item Centre : $\omega = \frac{2-2i\sqrt{3}}{1-(1+i\sqrt{3})} = \frac{2(1-i\sqrt{3})}{-i\sqrt{3}} = \frac{2(1-i\sqrt{3})i\sqrt{3}}{3} = \frac{2(i\sqrt{3}+3)}{3} = 2 + i\frac{2\sqrt{3}}{3}$.
\item Rapport : $k = |a| = \sqrt{1^2+(\sqrt{3})^2} = 2$.
\item Angle : $\cos\theta = \frac{1}{2}, \sin\theta = \frac{\sqrt{3}}{2} \Rightarrow \theta = \frac{\pi}{3}\ [2\pi]$.
\item Conclusion : $S$ est la similitude directe de centre $\Omega\left(2;\frac{2\sqrt{3}}{3}\right)$, rapport $2$, angle $\frac{\pi}{3}$.
\end{enumerate}
\end{exemplebox}

\courssub{Propri\'et\'es de conservation et images des configurations}

\begin{propriete}[Propri\'et\'es m\'etriques et graphiques]
Soit $S$ une similitude directe de rapport $k=|a|$ et d'angle $\theta=\arg(a)$.
\begin{enumerate}
\item \textbf{Conservation des angles orient\'es :} Pour tous points $A,B,C$ distincts,
      \[
      (\overrightarrow{AB}, \overrightarrow{AC}) = (\overrightarrow{A'B'}, \overrightarrow{A'C'})\ [2\pi].
      \]
\item \textbf{Conservation de l'alignement et du parall\'elisme :} L'image d'une droite est une droite ; des droites parall\'el\`es ont des images parall\'el\`es.
\item \textbf{Rapport des longueurs :} $\forall M,N,\quad M'N' = k \cdot MN$.
\item \textbf{Rapport des aires :} Pour toute figure plane d'aire $\mathcal{A}$, l'image a pour aire $\mathcal{A}' = k^2 \mathcal{A}$.
\end{enumerate}
\end{propriete}

\begin{propriete}[Image d'une droite]
L'image d'une droite $D$ par une similitude directe est une droite $D'$.
\textbf{M\'ethode :} D\'eterminer les images $A', B'$ de deux points distincts $A,B \in D$. Alors $D' = (A'B')$.
\end{propriete}

\begin{propriete}[Image d'un cercle]
L'image d'un cercle $\mathcal{C}(\Omega_C, R)$ par une similitude directe de rapport $k$ est le cercle $\mathcal{C}'(\Omega_{C'}, kR)$ o\`u $\Omega_{C'}$ est l'image du centre $\Omega_C$.
\emph{Preuve :} $M \in \mathcal{C} \iff \Omega_C M = R \iff k\cdot \Omega_C M = kR \iff \Omega_{C'} M' = kR \iff M' \in \mathcal{C}'$.
\end{propriete}

\begin{propriete}[Image d'un triangle et d'un parall\'elogramme]
\begin{itemize}
\item L'image d'un triangle est un triangle \textbf{semblable} et \textbf{directement orient\'e} (m\^eme sens de parcours des sommets).
\item L'image d'un parall\'elogramme est un parall\'elogramme (conservation du parall\'elisme et des milieux).
\end{itemize}
\end{propriete}

\begin{exemplebox}[Image d'une droite et d'un cercle]
Soit $S : z' = (1+i)z + 1$ (rapport $k=\sqrt{2}$, angle $\pi/4$, centre $\omega = \frac{1}{1-(1+i)} = \frac{1}{-i} = i$).
\begin{enumerate}
\item Image de la droite $D : \text{Re}(z)=0$ (axe des ordonn\'ees).
      Points $A(0), B(i)$. $A'(1), B' = (1+i)i+1 = i-1+1 = i$.
      $D'$ est la droite passant par $A'(1,0)$ et $B'(0,1)$ : \'equation $x+y=1$.
\item Image du cercle $\mathcal{C} : |z-2| = 3$ (centre $C(2)$, rayon $3$).
      Centre image : $C' = S(2) = (1+i)2+1 = 3+2i$.
      Rayon image : $R' = k \times 3 = 3\sqrt{2}$.
      Cercle image : $\mathcal{C}' : |z - (3+2i)| = 3\sqrt{2}$.
\end{enumerate}
\end{exemplebox}

\begin{savaistu}[Application concr\`ete : Cartographie et GPS]
Les similitudes directes mod\'elisent le passage d'une carte \`a une autre \`a \'echelle diff\'erente (homoth\'etie) combin\'ee \`a une rotation (orientation Nord diff\'erente). Si une carte de Yaound\'e (échelle 1/10 000) est num\'eris\'ee et qu'on veut l'afficher sur un \'ecran de t\'el\'ephone (échelle 1/50 000) avec une rotation de $90^\circ$ pour le mode paysage, la transformation des coordonn\'ees GPS (affixes complexes) est une similitude directe $z' = az+b$. Le rapport $k=1/5$ r\'edit les distances, l'angle $\theta=\pi/2$ tourne l'affichage. Le centre $\Omega$ correspond au point qui reste au centre de l'\'ecran lors du zoom/rotation.
\end{savaistu}

\courssub{D\'etermination d'une similitude et applications aux configurations}

\begin{methode}[Cas 1 : On conna\^it le centre $\Omega$, le rapport $k$, l'angle $\theta$]
\begin{enumerate}
\item $a = k e^{i\theta} = k(\cos\theta + i\sin\theta)$.
\item $b = \omega(1-a)$.
\item \'Equation : $z' = az + b$.
\end{enumerate}
\end{methode}

\begin{methode}[Cas 2 : On conna\^it les images de deux points distincts $A(z_A) \to A'(z_{A'})$, $B(z_B) \to B'(z_{B'})$]
\begin{enumerate}
\item V\'erifier $z_A \neq z_B$ (sinon infinit\'e de solutions).
\item Calculer $a = \dfrac{z_{B'} - z_{A'}}{z_B - z_A}$.
      \emph{Justification :} $z_{B'}-z_{A'} = a(z_B-z_A) \Rightarrow a = \frac{\overrightarrow{A'B'}}{\overrightarrow{AB}}$. Le module donne $k$, l'argument donne $\theta$.
\item Calculer $b = z_{A'} - a z_A$ (ou $b = z_{B'} - a z_B$).
\item \'Equation : $z' = az + b$.
\item \textbf{V\'erifier :} Tester sur $A$ et $B$.
\end{enumerate}
\end{methode}

\begin{attention}[Condition n\'ecessaire pour Cas 2]
Si $z_A = z_B$ mais $z_{A'} \neq z_{B'}$, \textbf{aucune similitude directe} n'envoie $A$ sur $A'$ et $B$ sur $B'$ (un point ne peut pas avoir deux images distinctes). Si $z_A=z_B$ et $z_{A'}=z_{B'}$, il y a une infinit\'e de similitudes (toutes celles qui envoient $A$ sur $A'$).
\end{attention}

\begin{exoresolu}[D\'etermination par deux couples de points]
D\'eterminer la similitude directe $S$ qui envoie $A(1+i)$ sur $A'(3)$ et $B(2)$ sur $B'(1+2i)$.
\tcblower
\textbf{Solution.}
\begin{enumerate}
\item $z_A = 1+i,\ z_{A'}=3$ ; $z_B=2,\ z_{B'}=1+2i$.
\item $a = \dfrac{z_{B'}-z_{A'}}{z_B-z_A} = \dfrac{(1+2i)-3}{2-(1+i)} = \dfrac{-2+2i}{1-i} = \dfrac{2(-1+i)}{1-i}$.
      Multiplier num\'erateur et d\'enominateur par le conjugu\'e $1+i$ :
      $a = \dfrac{2(-1+i)(1+i)}{(1-i)(1+i)} = \dfrac{2(-1 -i +i +i^2)}{2} = \dfrac{2(-2)}{2} = -2$.
\item $b = z_{A'} - a z_A = 3 - (-2)(1+i) = 3 + 2 + 2i = 5+2i$.
\item \'Equation : $\boxed{z' = -2z + 5+2i}$.
\item \textbf{V\'erification :} $S(B) = -2(2) + 5+2i = 1+2i = z_{B'}$. \checkmark
\item \textbf{Nature :} $a=-2 \in \mathbb{R}^*, a\neq 1 \Rightarrow$ \textbf{Homoth\'etie} de rapport $k=|a|=2$, angle $\theta=\pi$.
      Centre : $\omega = \frac{5+2i}{1-(-2)} = \frac{5+2i}{3} = \frac{5}{3} + i\frac{2}{3}$.
\end{enumerate}

\end{exoresolu}

\begin{exoresolu}[Probl\`eme de configuration : Triangles semblables]
Soient $A(0), B(4), C(1+2i)$ et $A'(1+i), B'(3+3i), C'(2+4i)$.
Montrer qu'il existe une similitude directe envoyant $A,B,C$ sur $A',B',C'$ et la d\'eterminer.
\tcblower
\textbf{Solution.}
\begin{enumerate}
\item On cherche $S$ telle que $S(A)=A', S(B)=B'$. (Deux points suffisent pour d\'efinir $S$, le troisi\`eme servira de v\'erification).
\item $a = \frac{z_{B'}-z_{A'}}{z_B-z_A} = \frac{(3+3i)-(1+i)}{4-0} = \frac{2+2i}{4} = \frac{1}{2}+\frac{1}{2}i = \frac{\sqrt{2}}{2}e^{i\pi/4}$.
      Donc $k = \frac{\sqrt{2}}{2}$, $\theta = \frac{\pi}{4}$.
\item $b = z_{A'} - a z_A = 1+i - 0 = 1+i$.
\item \'Equation candidate : $z' = \left(\frac{1}{2}+\frac{1}{2}i\right)z + 1+i$.
\item \textbf{V\'erification sur $C$ :}
      $z_C' = \left(\frac{1}{2}+\frac{1}{2}i\right)(1+2i) + 1+i = \frac{1}{2}(1+i)(1+2i) + 1+i$
      $= \frac{1}{2}(1+2i+i+2i^2) + 1+i = \frac{1}{2}(-1+3i) + 1+i = -\frac{1}{2}+\frac{3}{2}i + 1+i = \frac{1}{2} + \frac{5}{2}i$.
      Mais $z_{C'} = 2+4i = \frac{4}{2} + \frac{8}{2}i \neq \frac{1}{2} + \frac{5}{2}i$.
\item \textbf{Conclusion :} \textbf{Aucune similitude directe} n'envoie $ABC$ sur $A'B'C'$. Les triangles ne sont pas directement semblables (ou les correspondances de sommets ne sont pas les bonnes).
      \emph{Note :} Si on avait $C'(0.5+2.5i)$, cela aurait fonctionn\'e.
\end{enumerate}

\end{exoresolu}

\begin{methode}[R\'esolution d'un probl\`eme de configuration par similitude]
\begin{enumerate}
\item \textbf{Identifier} les deux figures suppos\'ees semblables (ex: deux triangles, un triangle et son image par une rotation/homoth\'etie).
\item \textbf{Choisir} une correspondance de sommets (ordre direct !).
\item \textbf{D\'eterminer} la similitude $S$ envoyant deux sommets de la premi\`ere figure sur les correspondants de la seconde (M\'ethode Cas 2).
\item \textbf{V\'erifier} que $S$ envoie bien le troisi\`eme sommet (et les autres \'el\'ements) au bon endroit.
\item \textbf{Exploiter} les propri\'et\'es : conservation des angles, rapport des longueurs $k$, position du centre $\Omega$ (souvent intersection de cercles circonscrits ou de droites particuli\`eres).
\end{enumerate}
\end{methode}

\begin{exemplebox}[Construction du centre d'une similitude envoyant $AB$ sur $A'B'$]
Soient $A,B,A',B'$ quatre points tels que $AB$ et $A'B'$ ne soient pas parall\`eles. Le centre $\Omega$ de la similitude directe envoyant $A$ sur $A'$ et $B$ sur $B'$ est l'intersection des droites $(AB)$ et $(A'B')$ \textbf{si elle existe}, sinon c'est l'intersection des cercles circonscrits aux triangles $AA'X$ et $BB'X$... \emph{(Rappel : $\Omega$ est le centre de la similitude qui transforme le segment $AB$ en $A'B'$. Il est l'intersection des droites $(AA')$ et $(BB')$ si ces droites ne sont pas parall\`eles, sinon c'est le point \`a l'infini de la direction commune $\to$ translation).}
\textbf{Propri\'et\'e utile :} $\Omega$ est le point d'intersection des droites $(AB)$ et $(A'B')$ \textbf{si et seulement si} $A,B,A',B'$ sont cocycliques (cas de la similitude ``inverse'' ou homoth\'etie). Pour une similitude directe quelconque, $\Omega$ est l'intersection des droites $(AA')$ et $(BB')$ (si non parall\`eles).
\end{exemplebox}

\begin{exercice}[Entra\^inement Bac - Niveau 1 : Identification]
Soient les transformations suivantes :
\begin{enumerate}
\item $S_1 : z' = 2z + 1 - i$
\item $S_2 : z' = -z + 3$
\item $S_3 : z' = \frac{\sqrt{2}}{2}(1+i)z + 2$
\item $S_4 : z' = z + 4i$
\item $S_5 : z' = (3-4i)z$
\end{enumerate}
Pour chacune, pr\'eciser la nature, le centre (si existant), le rapport et l'angle principal.
\end{exercice}

\begin{exercice}[Entra\^inement Bac - Niveau 2 : \'Equation et image]
Soit $S$ la similitude directe de centre $\Omega(1+i)$, de rapport $k=2$ et d'angle $\theta = -\frac{\pi}{2}$.
\begin{enumerate}
\item D\'eterminer son \'equation complexe $z'=az+b$ et son expression analytique.
\item D\'eterminer l'image du cercle $\mathcal{C} : |z-2| = 1$.
\item D\'eterminer l'image de la droite $D : y = x$.
\end{enumerate}
\end{exercice}

\begin{exercice}[Entra\^inement Bac - Niveau 3 : Probl\`eme de configuration]
Dans le plan muni d'un rep\`ere orthonorm\'e, on consid\`ere les points $A(0), B(4), C(1+3i)$.
On cherche une similitude directe $S$ qui transforme le triangle $ABC$ en un triangle $A'B'C'$ tel que $A'(2+2i)$ et $B'(4)$.
\begin{enumerate}
\item D\'eterminer l'\'equation complexe de $S$.
\item D\'eterminer les coordonn\'ees de $C'$.
\item Calculer l'aire du triangle $A'B'C'$ sachant que l'aire de $ABC$ est $6$ unit\'es d'aire.
\item Donner la nature de $S$, son centre, son rapport et son angle.
\end{enumerate}
\end{exercice}

\begin{corrige}[Exercice Niveau 1]
\begin{enumerate}
\item $S_1: a=2 \in \mathbb{R}^*, a\neq 1 \Rightarrow$ \textbf{Homoth\'etie} centre $\omega=\frac{1-i}{1-2}=-1+i$, rapport $k=2$, angle $0$.
\item $S_2: a=-1 \in \mathbb{R}^*, a\neq 1 \Rightarrow$ \textbf{Homoth\'etie} centre $\omega=\frac{3}{1-(-1)}=1.5$, rapport $k=1$, angle $\pi$ (sym\'etrie centrale).
\item $S_3: a=\frac{\sqrt{2}}{2}(1+i)=e^{i\pi/4}, |a|=1, a\neq 1 \Rightarrow$ \textbf{Rotation} centre $\omega=\frac{2}{1-e^{i\pi/4}}$, rapport $k=1$, angle $\pi/4$.
\item $S_4: a=1 \Rightarrow$ \textbf{Translation} vecteur $\vec{v}(0,4)$.
\item $S_5: a=3-4i, |a|=5, \arg(a)=-\arctan(4/3) \Rightarrow$ \textbf{Similitude directe} centre $\omega=0$, rapport $5$, angle $-\arctan(4/3)$.
\end{enumerate}
\end{corrige}

\begin{corrige}[Exercice Niveau 2]
\begin{enumerate}
\item $a = 2 e^{-i\pi/2} = -2i$. $b = \omega(1-a) = (1+i)(1+2i) = 1+2i+i+2i^2 = -1+3i$.
      $S: z' = -2i z -1+3i$.
      Analytique : $x' = 2y -1$, $y' = -2x + 3$.
\item Cercle $\mathcal{C}(2, 1)$. Centre image $C' = S(2) = -2i(2)-1+3i = -4i-1+3i = -1-i$. Rayon $R' = 2\times 1 = 2$.
      $\mathcal{C}' : |z+1+i| = 2$.
\item Droite $D: y=x \iff z = t(1+i), t\in\mathbb{R}$.
      $z' = -2i(t(1+i)) -1+3i = -2i(t+ti) -1+3i = -2it+2t -1+3i = (2t-1) + i(3-2t)$.
      $x' = 2t-1, y' = 3-2t \Rightarrow x'+y' = 2$. Image : droite $x+y=2$.
\end{enumerate}
\end{corrige}

\begin{corrige}[Exercice Niveau 3]
\begin{enumerate}
\item $a = \frac{z_{B'}-z_{A'}}{z_B-z_A} = \frac{4-(2+2i)}{4-0} = \frac{2-2i}{4} = \frac{1}{2} - \frac{1}{2}i = \frac{\sqrt{2}}{2}e^{-i\pi/4}$.
      $b = z_{A'} - a z_A = 2+2i$.
      $S: z' = \left(\frac{1}{2}-\frac{1}{2}i\right)z + 2+2i$.
\item $z_{C'} = \left(\frac{1}{2}-\frac{1}{2}i\right)(1+3i) + 2+2i = \frac{1}{2}(1-i)(1+3i) + 2+2i$
      $= \frac{1}{2}(1+3i-i-3i^2) + 2+2i = \frac{1}{2}(4+2i) + 2+2i = 2+i + 2+2i = 4+3i$.
      $C'(4,3)$.
\item $k = |a| = \frac{\sqrt{2}}{2}$. Aire $A'B'C' = k^2 \times 6 = \frac{1}{2} \times 6 = 3$.
\item Nature : Similitude directe vraie ($a\notin\mathbb{R}, |a|\neq 1$).
      Centre $\omega = \frac{b}{1-a} = \frac{2+2i}{1-(\frac{1}{2}-\frac{1}{2}i)} = \frac{2+2i}{\frac{1}{2}+\frac{1}{2}i} = \frac{2(1+i)}{\frac{1}{2}(1+i)} = 4$.
      Rapport $k = \frac{\sqrt{2}}{2}$. Angle $\theta = -\frac{\pi}{4}\ [2\pi]$ (ou $\frac{7\pi}{4}$).
\end{enumerate}
\end{corrige}

\begin{propriete}[Formules essentielles \`a conna\^itre par c\oe ur]
\begin{center}
\begin{tabularx}{\linewidth}{|l|X|}
\hline
\rowcolor{popblueL} \textbf{Objet} & \textbf{Formule} \\
\hline
\'Ecriture complexe g\'en\'erale & $z' = a z + b$, \quad $a \in \mathbb{C}^*,\ b \in \mathbb{C}$ \\
\hline
Forme r\'eduite (si $a \neq 1$) & $z' - \omega = a(z - \omega)$ \quad avec \quad $\omega = \dfrac{b}{1-a}$ \\
\hline
Rapport & $k = |a|$ \\
\hline
Angle (principal) & $\theta = \arg(a)\ [2\pi]$ \\
\hline
D\'ecomposition polaire & $a = k e^{i\theta} = k(\cos\theta + i\sin\theta)$ \\
\hline
Expression analytique & $\begin{cases} x' = \alpha x - \beta y + u \\ y' = \beta x + \alpha y + v \end{cases}$ \quad ($a=\alpha+i\beta,\ b=u+iv$) \\
\hline
Image d'un segment & $M'N' = k \cdot MN$ \quad ; \quad $(\overrightarrow{MN}, \overrightarrow{M'N'}) = \theta$ \\
\hline
Aire image & $\mathcal{A}' = k^2 \mathcal{A}$ \\
\hline
$S$ d\'efinie par $A\to A',\ B\to B'$ & $a = \dfrac{z_{B'}-z_{A'}}{z_B-z_A}$, \quad $b = z_{A'} - a z_A$ \\
\hline
\end{tabularx}
\end{center}
\end{propriete}

\begin{aretenir}[Checklist finale avant le Bac]
\begin{itemize}
  \item[$\square$] Je sais qu'une similitude directe s'\'ecrit $z'=az+b$ avec $a\neq 0$.
  \item[$\square$] J'identifie imm\'ediatement les cas particuliers : translation ($a=1$), homoth\'etie ($a\in\mathbb{R}^*\setminus\{1\}$), rotation ($|a|=1, a\neq 1$).
  \item[$\square$] Je calcule le centre par $\omega=\dfrac{b}{1-a}$ \textbf{uniquement si $a\neq 1$}.
  \item[$\square$] Je sais que le rapport est $k=|a|$ et l'angle $\theta=\arg(a)$.
  \item[$\square$] Je passe de l'\'ecriture complexe \`a l'analytique via $a=\alpha+i\beta, b=u+iv$.
  \item[$\square$] Je d\'etermine $a$ et $b$ \`a partir des images de deux points distincts : $a=\frac{z_{B'}-z_{A'}}{z_B-z_A}$.
  \item[$\square$] Je sais que les longueurs sont multipli\'ees par $k$ et les aires par $k^2$.
  \item[$\square$] Je sais que l'image d'une droite est une droite (par 2 points) et celle d'un cercle un cercle (centre image, rayon $\times k$).
  \item[$\square$] Je sais qu'une similitude directe conserve les angles orient\'es, l'alignement, le parall\'elisme et le rapport des longueurs.
  \item[$\square$] Je v\'erifie syst\'ematiquement mes calculs en testant l'image d'un point connu.
  \item[$\square$] Je r\'edige proprement ma conclusion : \emph{Nature, Centre, Rapport, Angle}.
  \item[$\square$] Je n'oublie pas qu'une similitude directe transforme un triangle en un triangle \emph{semblable et directement orient\'e}.
\end{itemize}
\end{aretenir}

\findecours

\end{document}
